% Options for packages loaded elsewhere
\PassOptionsToPackage{unicode}{hyperref}
\PassOptionsToPackage{hyphens}{url}
\PassOptionsToPackage{dvipsnames,svgnames,x11names}{xcolor}
\PassOptionsToPackage{space}{xeCJK}
%
\documentclass[
  9pt,
  a4paper,
]{article}
\usepackage{amsmath,amssymb}
\usepackage{iftex}
\ifPDFTeX
  \usepackage[T1]{fontenc}
  \usepackage[utf8]{inputenc}
  \usepackage{textcomp} % provide euro and other symbols
\else % if luatex or xetex
  \usepackage{fontspec}
  \defaultfontfeatures{Scale=MatchLowercase}
  \defaultfontfeatures[\rmfamily]{Ligatures=TeX,Scale=1}
\fi
\usepackage{lmodern}
\ifPDFTeX\else
  % xetex/luatex font selection
  \setmainfont[]{Noto Sans TC}
  \setmonofont[]{Consolas}
  \ifXeTeX
    \usepackage{xeCJK}
    \setCJKmainfont[]{Noto Sans TC}
    \setCJKmonofont[]{Noto Sans TC}
          \fi
  \ifLuaTeX
    \usepackage[]{luatexja-fontspec}
    \setmainjfont[]{Noto Sans TC}
  \fi
\fi
% Use upquote if available, for straight quotes in verbatim environments
\IfFileExists{upquote.sty}{\usepackage{upquote}}{}
\IfFileExists{microtype.sty}{% use microtype if available
  \usepackage[]{microtype}
  \UseMicrotypeSet[protrusion]{basicmath} % disable protrusion for tt fonts
}{}
\makeatletter
\@ifundefined{KOMAClassName}{% if non-KOMA class
  \IfFileExists{parskip.sty}{%
    \usepackage{parskip}
  }{% else
    \setlength{\parindent}{0pt}
    \setlength{\parskip}{6pt plus 2pt minus 1pt}}
}{% if KOMA class
  \KOMAoptions{parskip=half}}
\makeatother
\usepackage{xcolor}
\usepackage[margin=15mm]{geometry}
\usepackage{color}
\usepackage{fancyvrb}
\newcommand{\VerbBar}{|}
\newcommand{\VERB}{\Verb[commandchars=\\\{\}]}
\DefineVerbatimEnvironment{Highlighting}{Verbatim}{commandchars=\\\{\}}
% Add ',fontsize=\small' for more characters per line
\newenvironment{Shaded}{}{}
\newcommand{\AlertTok}[1]{\textcolor[rgb]{1.00,0.00,0.00}{\textbf{#1}}}
\newcommand{\AnnotationTok}[1]{\textcolor[rgb]{0.38,0.63,0.69}{\textbf{\textit{#1}}}}
\newcommand{\AttributeTok}[1]{\textcolor[rgb]{0.49,0.56,0.16}{#1}}
\newcommand{\BaseNTok}[1]{\textcolor[rgb]{0.25,0.63,0.44}{#1}}
\newcommand{\BuiltInTok}[1]{\textcolor[rgb]{0.00,0.50,0.00}{#1}}
\newcommand{\CharTok}[1]{\textcolor[rgb]{0.25,0.44,0.63}{#1}}
\newcommand{\CommentTok}[1]{\textcolor[rgb]{0.38,0.63,0.69}{\textit{#1}}}
\newcommand{\CommentVarTok}[1]{\textcolor[rgb]{0.38,0.63,0.69}{\textbf{\textit{#1}}}}
\newcommand{\ConstantTok}[1]{\textcolor[rgb]{0.53,0.00,0.00}{#1}}
\newcommand{\ControlFlowTok}[1]{\textcolor[rgb]{0.00,0.44,0.13}{\textbf{#1}}}
\newcommand{\DataTypeTok}[1]{\textcolor[rgb]{0.56,0.13,0.00}{#1}}
\newcommand{\DecValTok}[1]{\textcolor[rgb]{0.25,0.63,0.44}{#1}}
\newcommand{\DocumentationTok}[1]{\textcolor[rgb]{0.73,0.13,0.13}{\textit{#1}}}
\newcommand{\ErrorTok}[1]{\textcolor[rgb]{1.00,0.00,0.00}{\textbf{#1}}}
\newcommand{\ExtensionTok}[1]{#1}
\newcommand{\FloatTok}[1]{\textcolor[rgb]{0.25,0.63,0.44}{#1}}
\newcommand{\FunctionTok}[1]{\textcolor[rgb]{0.02,0.16,0.49}{#1}}
\newcommand{\ImportTok}[1]{\textcolor[rgb]{0.00,0.50,0.00}{\textbf{#1}}}
\newcommand{\InformationTok}[1]{\textcolor[rgb]{0.38,0.63,0.69}{\textbf{\textit{#1}}}}
\newcommand{\KeywordTok}[1]{\textcolor[rgb]{0.00,0.44,0.13}{\textbf{#1}}}
\newcommand{\NormalTok}[1]{#1}
\newcommand{\OperatorTok}[1]{\textcolor[rgb]{0.40,0.40,0.40}{#1}}
\newcommand{\OtherTok}[1]{\textcolor[rgb]{0.00,0.44,0.13}{#1}}
\newcommand{\PreprocessorTok}[1]{\textcolor[rgb]{0.74,0.48,0.00}{#1}}
\newcommand{\RegionMarkerTok}[1]{#1}
\newcommand{\SpecialCharTok}[1]{\textcolor[rgb]{0.25,0.44,0.63}{#1}}
\newcommand{\SpecialStringTok}[1]{\textcolor[rgb]{0.73,0.40,0.53}{#1}}
\newcommand{\StringTok}[1]{\textcolor[rgb]{0.25,0.44,0.63}{#1}}
\newcommand{\VariableTok}[1]{\textcolor[rgb]{0.10,0.09,0.49}{#1}}
\newcommand{\VerbatimStringTok}[1]{\textcolor[rgb]{0.25,0.44,0.63}{#1}}
\newcommand{\WarningTok}[1]{\textcolor[rgb]{0.38,0.63,0.69}{\textbf{\textit{#1}}}}
\usepackage{longtable,booktabs,array}
\usepackage{calc} % for calculating minipage widths
% Correct order of tables after \paragraph or \subparagraph
\usepackage{etoolbox}
\makeatletter
\patchcmd\longtable{\par}{\if@noskipsec\mbox{}\fi\par}{}{}
\makeatother
% Allow footnotes in longtable head/foot
\IfFileExists{footnotehyper.sty}{\usepackage{footnotehyper}}{\usepackage{footnote}}
\makesavenoteenv{longtable}
\setlength{\emergencystretch}{3em} % prevent overfull lines
\providecommand{\tightlist}{%
  \setlength{\itemsep}{0pt}\setlength{\parskip}{0pt}}
\setcounter{secnumdepth}{-\maxdimen} % remove section numbering
\ifLuaTeX
\usepackage[bidi=basic]{babel}
\else
\usepackage[bidi=default]{babel}
\fi
\babelprovide[main,import]{chinese-hant}
\ifPDFTeX
\else
\babelfont{rm}[]{Noto Sans TC}
\fi
% get rid of language-specific shorthands (see #6817):
\let\LanguageShortHands\languageshorthands
\def\languageshorthands#1{}
\usepackage{fvextra}
\fvset{breaklines=true,breakanywhere=true,fontsize=\small}
\usepackage{microtype}
\usepackage{booktabs,longtable,array}
\usepackage{ragged2e}
\usepackage{etoolbox}
\usepackage{setspace}
\setstretch{1.04}
\setlength{\emergencystretch}{5em}
\setlength{\tabcolsep}{3pt}
\renewcommand{\arraystretch}{1.15}
\AtBeginEnvironment{longtable}{\footnotesize}
\sloppy
\usepackage{enumitem}
\setlist{nosep,leftmargin=*}
\usepackage{titlesec}
\titlespacing*{\section}{0pt}{1.2em}{0.5em}
\titlespacing*{\subsection}{0pt}{1em}{0.35em}
\usepackage{caption}
\captionsetup{font=small}
\ifLuaTeX
  \usepackage{selnolig}  % disable illegal ligatures
\fi
\usepackage{bookmark}
\IfFileExists{xurl.sty}{\usepackage{xurl}}{} % add URL line breaks if available
\urlstyle{same}
\hypersetup{
  pdftitle={排球賽事標註、AI 串接與教練監看系統},
  pdflang={zh-Hant},
  colorlinks=true,
  linkcolor={teal},
  filecolor={Maroon},
  citecolor={Blue},
  urlcolor={teal},
  pdfcreator={LaTeX via pandoc}}

\title{排球賽事標註、AI 串接與教練監看系統}
\usepackage{etoolbox}
\makeatletter
\providecommand{\subtitle}[1]{% add subtitle to \maketitle
  \apptocmd{\@title}{\par {\large #1 \par}}{}{}
}
\makeatother
\subtitle{實作總規格 v3.2}
\author{}
\date{2026-08-07}

\begin{document}
\maketitle

\renewcommand*\contentsname{目錄}
{
\hypersetup{linkcolor=}
\setcounter{tocdepth}{3}
\tableofcontents
}
\begin{quote}
Repository：\texttt{volleyball-monitoring-ai}\\
主要讀者：負責統籌實作的主 PM／架構 Agent，以及最多三個受委派
subagent。\\
文件目的：把產品行為、跨團隊 Schema、媒體時間軸、外部 AI 串接、Nuxt
UI、後端、容器與開發順序定成可實作的共同基線。\\
核心邊界：\textbf{本 repository 不實作任何 AI
模型；只實作中央系統、影音流程、前端、外部 AI 介面、Fake Provider 與可由
GitHub 安裝的 Python SDK。}
\end{quote}

\begin{center}\rule{0.5\linewidth}{0.5pt}\end{center}

\subsection{0. 主 Agent
必須先理解的事情}\label{0-ux4e3b-agent-ux5fc5ux9808ux5148ux7406ux89e3ux7684ux4e8bux60c5}

這不是單純的影片播放器、標註工具或 AI
Dashboard，而是一條完整且必須可追溯的資料鏈：

\begin{Shaded}
\begin{Highlighting}[]
\NormalTok{直播訊號}
\NormalTok{→ 伺服器完整錄製與建立 canonical timeline}
\NormalTok{→ 多人標註端在 live／歷史回放畫面上建立人工 key point}
\NormalTok{→ 人工確認並提交 immutable rally submission}
\NormalTok{→ 後端依 authoritative source time 裁切影片並換算 clip{-}local time}
\NormalTok{→ 外部 AI 透過固定 Job Schema 收件}
\NormalTok{→ 外部 AI 透過 callback 回傳 Analysis JSON + FlatBuffers overlay}
\NormalTok{→ 中央系統驗證、保存、正規化與聚合}
\NormalTok{→ 教練／裁判端查看歷史、影片 overlay、2D 球路、熱點與統計}
\end{Highlighting}
\end{Shaded}

主 Agent 的第一責任是確保這條鏈中的
ID、時間、版本與資料語意不會斷裂。任何只完成 UI、只完成 API、只完成 DB
table 或只完成 demo 的工作，都不能宣稱該流程已完成。

主 Agent最多同時委派三個 subagent：

\begin{enumerate}
\def\labelenumi{\arabic{enumi}.}
\tightlist
\item
  \textbf{Contracts／Python SDK Agent}
\item
  \textbf{Backend／Media／Infra Agent}
\item
  \textbf{Nuxt Luna Frontend Agent}
\end{enumerate}

主 Agent保留以下不可委派的最終責任：

\begin{itemize}
\tightlist
\item
  產品語意與 Schema 變更核准。
\item
  跨目錄修改與 merge 衝突。
\item
  時間軸與 passthrough invariant 審查。
\item
  端到端驗收。
\item
  是否接受 subagent 提出的架構調整。
\end{itemize}

\begin{center}\rule{0.5\linewidth}{0.5pt}\end{center}

\section{1. 固定產品行為}\label{1-ux56faux5b9aux7522ux54c1ux884cux70ba}

\subsection{1.1
標註快捷鍵與介面控制}\label{11-ux6a19ux8a3bux5febux6377ux9375ux8207ux4ecbux9762ux63a7ux5236}

下列 command 語意與畫面按鈕是產品合約，不得由實作者自行改成其他語意。表中鍵位是預設值，使用者可以在設定選單修改實體快捷鍵；觸控按鈕與目前鍵盤綁定必須呼叫同一個 command path。

\begin{longtable}[]{@{}>{\RaggedRight\arraybackslash}p{0.11\linewidth}>{\RaggedRight\arraybackslash}p{0.19\linewidth}>{\RaggedRight\arraybackslash}p{0.31\linewidth}>{\RaggedRight\arraybackslash}p{0.33\linewidth}@{}}
\toprule\noalign{}
預設介面顯示 & 預設快捷鍵 & 固定語意 & 重要限制 \\
\midrule\noalign{}
\endhead
\bottomrule\noalign{}
\endlastfoot
\texttt{Z\ 發球} & \texttt{Z} & 建立新 rally，並在目前已呈現影片 frame
建立第一個人工 \texttt{service} key point & 已有未提交 rally
時不可另開新 rally \\
\texttt{Space\ 擊球} & \texttt{Space} & 在目前已呈現影片 frame 建立一般
\texttt{contact} key point & 播放器正在 seek、cursor stale 或位於 gap
時不可建立 \\
\texttt{\textless{}\ 左側得分} & \texttt{\textless{}} &
以單一 \texttt{CLOSE\_\allowbreak{}RALLY} atomically把目前server-confirmed最後key
point標為terminal，並將rally-level outcome設為\texttt{resolved/left} &
Command必須帶\texttt{target\_\allowbreak{}key\_\allowbreak{}point\_\allowbreak{}id}；若它已不是最後一點，回傳revision conflict；不建立新時間或得分事件 \\
\texttt{\textgreater{}\ 右側得分} & \texttt{\textgreater{}} &
同樣atomically關閉rally並將rally-level outcome設為\texttt{resolved/right} &
同上；方向鍵仍只供逐幀／播放器控制 \\
\texttt{?\ 未知} & \texttt{?} &
同樣atomically關閉rally並將rally-level outcome設為\texttt{unknown/null} &
不是\texttt{pending}；可以提交AI但不納入勝負統計；不建立新時間或得分事件 \\
\texttt{Enter\ 提交} & \texttt{Enter} & 將目前 draft 建立成 immutable
\texttt{RallySubmission}，啟動裁切與 AI 串接 & \texttt{pending}
不可提交；\texttt{resolved} 或 \texttt{unknown} 可提交 \\
\end{longtable}

\subsubsection{快捷鍵自訂邊界}

\begin{itemize}
\tightlist
\item
  使用者可以修改 annotation 與播放器 command 的實體鍵位；修改鍵位不得改變 command 語意、WebSocket command kind 或 server-side validation。
\item
  每個可用 command 必須保留一個有效鍵位。設定選單使用集中式 command registry、快捷鍵錄製器與衝突檢查，不得由各 component 各自監聽或保存。
\item
  設定選單必須提供「還原所有預設快捷鍵」。還原後回到本節表格與方向鍵的預設值。
\item
  相同 scope 不得有重複鍵位；瀏覽器保留鍵、文字輸入焦點、Dialog/Select scope 與平台鍵盤差異必須被辨識並清楚提示。
\item
  Control deck、設定選單與快捷鍵說明必須以 TanStack Hotkeys \texttt{formatForDisplay} 顯示目前綁定，讓 macOS 使用符號、Windows/Linux 使用平台慣用標籤；不得自行拼接鍵帽文字。六個觸控動作不可被移除，且不受鍵盤自訂影響。
\end{itemize}

\texttt{service} 與 \texttt{contact} 只是人工 marker kind，不是 AI
action label，也沒有人工 confidence。第一個 key point 被標成
\texttt{service} 只代表標註流程由 Z 開始；AI 是否輸出任何 action、action
label 長什麼樣、是否有 confidence，均屬待 AI 團隊提供真實輸出後確認的
optional extension。

\subsubsection{按鍵模式優先順序}\label{ux6309ux9375ux6a21ux5f0fux512aux5148ux9806ux5e8f}

\begin{enumerate}
\def\labelenumi{\arabic{enumi}.}
\tightlist
\item
  文字輸入、Dialog 或 Select 開啟時，禁止攔截全域快捷鍵。
\item
  \texttt{OPEN}且存在server-confirmed最後key point時，左側得分、右側得分、未知等
  \texttt{CLOSE\_\allowbreak{}RALLY} command依目前綁定觸發；預設為
  \texttt{\textless{}}、\texttt{\textgreater{}}、\texttt{?}。
\item
  前後 canonical frame／播放器移動 command 預設綁定
  \texttt{←}／\texttt{→}；使用者可改鍵，但該 command 不得變成得分或 annotation data mutation。
\item
  其他狀態下播放器 command 只控制播放器，不可暗中改資料。
\item
  每次 destructive action 都必須等待 server ack；optimistic UI 可顯示
  pending，但不得把暫態當 canonical revision。
\end{enumerate}

\subsection{1.2 Rally Mask
與狀態呈現}\label{12-rally-mask-ux8207ux72c0ux614bux5448ux73fe}

\begin{itemize}
\tightlist
\item
  \texttt{OPEN}、\texttt{READY}：灰色半透明
  mask，表示尚未提交且仍可修正。
\item
  \texttt{SUBMITTED}：綠色半透明 mask，表示某個 immutable submission
  已建立。
\item
  綠色不代表 AI 完成；裁切、排隊、處理中、完成、失敗以 mask 上的 status
  badge 顯示。
\item
  Mask 上緣必須有文字狀態，不可只靠顏色：

  \begin{itemize}
  \tightlist
  \item
    \texttt{左側得分}
  \item
    \texttt{右側得分}
  \item
    \texttt{?\ 得分未知}
  \item
    \texttt{回合進行中}
  \end{itemize}
\item
  Mask 範圍永遠是 service key point 到 terminal key
  point；pre-roll/post-roll 只屬裁切範圍，不改變 mask。
\item
  Key point marker 至少可辨識 \texttt{Z/service}、一般
  \texttt{contact}、terminal、selected、possible duplicate、server
  pending。
\item
  顏色、線寬與動畫屬前端衍生視覺，不保存為 canonical DB 欄位。
\end{itemize}

\subsection{1.3
教練端不是單頁}\label{13-ux6559ux7df4ux7aefux4e0dux662fux55aeux9801}

教練／裁判端採多頁資訊架構，iPad
使用底部導覽列；頂部狀態列可返回主選單並開啟過去保存的紀錄。

建議底部導覽：

\begin{enumerate}
\def\labelenumi{\arabic{enumi}.}
\tightlist
\item
  \texttt{現場}
\item
  \texttt{球路}
\item
  \texttt{球員}
\item
  \texttt{統計}
\item
  \texttt{紀錄}
\end{enumerate}

頂部列至少包含：

\begin{itemize}
\tightlist
\item
  返回賽事／場次選單。
\item
  場次、局數、比分與 live 狀態。
\item
  當前是否位於 live edge。
\item
  同步／離線狀態。
\item
  Overlay 快捷設定。
\item
  進入歷史 rally、保存分析視圖與設定。
\end{itemize}

\subsection{1.4 AI 邊界}\label{14-ai-ux908aux754c}

外部 AI
團隊已有可利用的球場、球員、球、個體動作，以及可能存在的群體動作／比賽
phase work。本 repository：

\begin{itemize}
\tightlist
\item
  \textbf{不實作、不重訓、不替換這些模型。}
\item
  不規定 detector、tracker、action model 的架構。
\item
  不規定時間窗、IOU threshold、融合權重、補點參數或模型內部順序。
\item
  只規定目標導向的方法骨架、輸入、輸出、錯誤、版本與串接方式。
\item
  action taxonomy、action confidence、群體 phase 的 label 與實際用途，需
  AI 團隊提供真實輸出樣本後另行確認；baseline 不能硬編碼。
\end{itemize}

\begin{center}\rule{0.5\linewidth}{0.5pt}\end{center}

\section{2.
實作與使用角度稽核：已修正問題}\label{2-ux5be6ux4f5cux8207ux4f7fux7528ux89d2ux5ea6ux7a3dux6838ux5df2ux4feeux6b63ux554fux984c}

以下為原規格若直接實作會產生的問題，以及本文件採用的修正版。

\begin{longtable}[]{@{}>{\RaggedRight\arraybackslash}p{0.11\linewidth}>{\RaggedRight\arraybackslash}p{0.19\linewidth}>{\RaggedRight\arraybackslash}p{0.31\linewidth}>{\RaggedRight\arraybackslash}p{0.33\linewidth}@{}}
\toprule\noalign{}
嚴重度 & 問題 & 使用／實作風險 & 修正後基線 \\
\midrule\noalign{}
\endhead
\bottomrule\noalign{}
\endlastfoot
Critical & \texttt{left/right} 被當成永久隊伍 ID &
換邊後歷史得分、熱點與球員統計錯隊 & \texttt{court\_\allowbreak{}side} 與
\texttt{team\_\allowbreak{}id} 分離；每 rally 綁定當下的 side assignment \\
Critical & key point 直接信任 browser \texttt{currentTime} 或
\texttt{currentTime\ *\ fps} & HLS、VFR、seek、延遲與轉碼後會指向錯幀 &
Client 送 \texttt{PlaybackCursor}；後端依 playback mapping 解出
authoritative PTS/time/frame \\
Critical & 只提供幾分鐘回放 buffer &
比賽後段無法返回開場；重新整理後資料消失 & 伺服器完整錄整場；前端只
windowed lazy-load，目前 buffer 幾分鐘但完整 timeline 可 seek \\
Critical & \texttt{court\_\allowbreak{}pos} 被限制 \texttt{{[}0,1{]}} &
發球員、救球、出界點會被 clamp 到場內 & 球場內以 0--1
為基準，但允許負值與大於 1；前端可裁切顯示，資料不可 clamp \\
Critical & \texttt{multiple} 同時表示共同參與與無法分辨 &
兩名重疊球員會被錯算成共同觸球 &
\texttt{resolved\_\allowbreak{}single}、\texttt{resolved\_\allowbreak{}multiple}、\texttt{ambiguous}、\texttt{unresolved}
分開；actors/candidates 分開 \\
Critical & tracker ID 被當成跨 rally 球員 ID & 同一整場球員統計會錯綁 &
\texttt{track\_\allowbreak{}id} 只在單 rally／analysis run 內有效；中央另建 identity
assignment \\
Critical & 每次編輯都複製完整 revision snapshot & 大量 DB
寫入與多人同步成本不必要 & Draft 使用 mutable rows + operation log +
revision；只有 Enter 建 immutable submission snapshot \\
High & \texttt{?} 與「尚未選得分」都用 null & UI
無法區分使用者已明確標未知 & 新增 \texttt{score\_\allowbreak{}resolution\_\allowbreak{}state}:
\texttt{pending/resolved/unknown}；\texttt{?} 後可提交，pending
不可提交 \\
High & 舊流程以獨立end-rally動作建立另一個timestamp &
與使用者定義衝突，也會多一個虛假事件 &
\texttt{CLOSE\_\allowbreak{}RALLY}只標記指定的既有最後key point為terminal並保存rally-level
outcome；command不含新時間 \\
High & 沒有 reopen／改 terminal 的流程 & 關閉rally誤按後只能刪整個rally &
Draft提供\texttt{REOPEN\_\allowbreak{}RALLY}；重開後以新的
\texttt{CLOSE\_\allowbreak{}RALLY}重新指定目前最後key point與outcome；submitted必須建立correction draft \\
High & Rally 只有單一 status & annotation、clip、AI 狀態互相覆寫 &
\texttt{annotation\_\allowbreak{}status} 與 \texttt{processing\_\allowbreak{}status} 分離 \\
High & submitted draft 可被原地修改 & AI 結果無法知道對應哪個版本 &
submission 永久 immutable；修正建立新 submission，舊 job/result 標
superseded \\
High & Score 修正必定重跑 AI & 只改勝方卻浪費推論 & 比對 submission
content hash；只改 outcome 可重用 clip/AI geometry，重算分析聚合 \\
High & AI Job 混入 \texttt{court\_\allowbreak{}definition}、模型 requirements、上傳
URI & AI request 像是在遠端配置 pipeline，耦合中央儲存 &
固定座標規格只寫一次；Job 僅含 clip、key points、outcome、callback 與
passthrough IDs \\
High & AI Result 的 \texttt{x\_\allowbreak{}m/y\_\allowbreak{}m/x\_\allowbreak{}norm/y\_\allowbreak{}norm} 多套命名 &
前後端容易混用或方向錯誤 & 對外只用
\texttt{frame\_\allowbreak{}pos}、\texttt{frame\_\allowbreak{}bbox}、\texttt{court\_\allowbreak{}pos}；公尺值由固定公式需要時衍生 \\
High & AI actor 只有單一 \texttt{actor\_\allowbreak{}track\_\allowbreak{}id} &
不能表達雙人攔網、共同事件或不確定候選 & \texttt{actors{[}{]}} 與
\texttt{candidates{[}{]}}；每個 contact event 具有 association state \\
High & Action label 被寫死 & AI 團隊尚未確認 taxonomy，會被假規格綁住 &
action 是 optional extension，label 是 provider 提供字串，帶 taxonomy
ID/version \\
High & Confidence 被視為必填 & 實際模型未必輸出可校準 confidence & 所有
confidence optional；缺少時不得前端自行補 0 或 1 \\
High & AI input/output IDs 沒有 passthrough 規則 & callback 對不到
submission/key point &
\texttt{ai\_\allowbreak{}job\_\allowbreak{}id/rally\_\allowbreak{}id/submission\_\allowbreak{}id/annotation\_\allowbreak{}revision/key\_\allowbreak{}point\_\allowbreak{}id}
必須原樣回傳 \\
High & 一次 WebSocket 傳整段逐幀 JSON & iPad 記憶體、解析與 GC 壓力過大
& JSON 管 metadata；FlatBuffers 管 per-frame overlay；HTTP lazy
chunk，WS 只通知 ready/version \\
High & Clip URL 與 callback token 共用同一 bearer & SDK 若把 callback
token帶到物件儲存，會擴大secret暴露面 & \texttt{clip.download\_\allowbreak{}url}
是獨立、短期簽名URL；\texttt{callback.token}只可送到中央callback，兩者不可混用且都不進browser \\
High & callback token 真正一次性 & 網路重試會失敗 & job-scoped、TTL
內可重試；每次 callback 有唯一 \texttt{callback\_\allowbreak{}id} 做 idempotency \\
High & 沒有 media discontinuity/gap & 來源重連後時間軸假裝連續，key
point 會錯段 & capture timeline 保存 discontinuity；UI 顯示 gap，不允許
seek/mark 到缺檔區 \\
High & GraphQL \texttt{Int} 保存 PTS／microseconds／size & GraphQL Int
僅 32-bit；JS Number 也有精度風險 & 使用 \texttt{BigInt} scalar，wire
一律 decimal string；DB 使用 PostgreSQL BIGINT \\
High & MinIO credential 交給 browser & 物件儲存暴露 & 所有 S3 操作
server-side；前端只取得受權限控制的短期 URL或同源 proxy route \\
Medium & Client 自行選 pre/post-roll & 每個人裁切結果不一致 &
match/system clip profile 決定；submit 只提交 revision \\
Medium & 只保存色彩，不保存狀態 & 改版後資料無法解讀 & 保存狀態
enum，顏色純 UI \\
Medium & Action 不存在時仍顯示 attack 指標 & 顯示空值或誤導數字 & API
回傳 \texttt{feature\_\allowbreak{}availability} 與
\texttt{available\_\allowbreak{}dimensions/metrics} \\
Medium & 統計把 rally win 當成某擊球直接得分 & 中間 attack 會被誤算 kill
& 分離 \texttt{rally\_\allowbreak{}outcome} 與 optional
\texttt{direct\_\allowbreak{}outcome}；baseline 只保證前者 \\
Medium & WebSocket 斷線後只接續增量 & 錯過 revision 後狀態永久錯誤 &
reconnect 先比較 revision；有 gap 立即 GraphQL refetch
snapshot，再訂閱增量 \\
Medium & 高頻操作透過 GraphQL subscription 與大型 snapshot &
resolver/loading/serialization 不必要 & GraphQL 管結構化查詢；自訂 WS 管
annotation command/revision/presence \\
\end{longtable}

\begin{center}\rule{0.5\linewidth}{0.5pt}\end{center}

\subsection{2.1 歷史稽核：v3.0
首次補齊的缺口}\label{21-ux6b77ux53f2ux7a3dux6838v30-ux9996ux6b21ux88dcux9f4aux7684ux7f3aux53e3}

\begin{longtable}[]{@{}>{\RaggedRight\arraybackslash}p{0.18\linewidth}>{\RaggedRight\arraybackslash}p{0.34\linewidth}>{\RaggedRight\arraybackslash}p{0.42\linewidth}@{}}
\toprule\noalign{}
嚴重度 & 缺口 & 修正 \\
\midrule\noalign{}
\endhead
\bottomrule\noalign{}
\endlastfoot
Critical & 原始 PTS 在來源重連後可能從零或跳變，只有 discontinuity
number 不足以關聯 key point & 新增 \texttt{CaptureEpoch}；所有 resolved
marker 保存
\texttt{capture\_\allowbreak{}epoch\_\allowbreak{}id\ +\ source\_\allowbreak{}pts\ +\ capture\_\allowbreak{}time\_\allowbreak{}us}。\texttt{capture\_\allowbreak{}time\_\allowbreak{}us}
是整個 capture session 單調座標，raw PTS 只在 epoch 內有效。 \\
Critical & terminal key point 可能是球落地、出界或無球員歸屬，卻被
schema 強迫一定有 actor & AI event 新增
\texttt{no\_\allowbreak{}player}/\texttt{unresolved} 狀態，\texttt{actors{[}{]}}
可以為空；\texttt{court\_\allowbreak{}pos} 也可為 null並帶 quality flag。 \\
Critical & AI result 的 \texttt{court\_\allowbreak{}side} 容易被當成永久 team
identity & AI 只輸出 \texttt{court\_\allowbreak{}side:\ left/right/unknown}；中央依
submission 的 side-assignment snapshot解析 \texttt{team\_\allowbreak{}id}。 \\
Critical & clip key point若引用 mutable draft ID，correction後
callback無法追溯 & AI Job只引用 immutable
\texttt{RallySubmissionKeyPoint.id}；另保留
\texttt{source\_\allowbreak{}draft\_\allowbreak{}key\_\allowbreak{}point\_\allowbreak{}id} 供 audit。 \\
High & 得分欄位只有「誰得分」，沒有可重放的比分 ledger & 新增 immutable
\texttt{PointAward}，保存 submission、scoring team、set score
before/after與\texttt{score\_\allowbreak{}revision\_\allowbreak{}before/after}；每set的after
revision唯一，比分由ledger衍生。 \\
High & 同一場次可能停止再啟動錄影；單一 capture session假設會使
timeline斷裂 & Match可有多個 \texttt{CaptureSession}，但 baseline標註
room指定一個 active capture；跨 capture replay由 match timeline聚合。 \\
High & 只保存 callback token，未定義 retry/idempotency & token是
job-scoped且有TTL，可重試；每次 callback傳
\texttt{callback\_\allowbreak{}id}，中央以唯一索引與payload checksum去重。 \\
High & 只定義 AI complete callback，無明確 accepted/failed行為 &
SDK與REST合約定義 \texttt{accepted}、optional
\texttt{progress}、\texttt{failed}、\texttt{completed}；progress不影響
correctness，failed保存可重試分類。 \\
High & AI event count沒有明確與 key point 對應規則 & Required
invariant：每個 submission key point恰好一個分析 event；segment數通常為
\texttt{max(event\_\allowbreak{}count-1,\ 0)}，只有service且terminal的service
error合法為0段。 \\
High & \texttt{frame\_\allowbreak{}pos}/\texttt{court\_\allowbreak{}pos} 缺失時沒有語意 & 兩者為
nullable；\texttt{quality\_\allowbreak{}flags}與\texttt{association\_\allowbreak{}state}表達缺失，禁止以
\texttt{(0,0)} 代表 unknown。 \\
High & Overlay binary
schema若把court座標量化到uint16，無法表示場外負值/大於1 &
\texttt{frame\_\allowbreak{}pos/frame\_\allowbreak{}bbox}可uint16量化；\texttt{court\_\allowbreak{}pos}使用float32，允許場外值。 \\
High & 沒有 API authorization boundary &
新增角色/permission基線：administrator、operator、annotator、coach/viewer、integration
service；MinIO與AI token永不下發瀏覽器。 \\
High & GraphQL source與SDL可能雙向手改而漂移 & Pothos
source唯一；SDL為產物，CI重生後必須clean，並執行breaking-change檢查。 \\
High & GraphQL scalar BigInt若被 codegen成number仍會精度損失 & Web
codegen將 \texttt{BigInt} 映射為 \texttt{string}；DB為BIGINT；JSON
Schema也用decimal string。 \\
Medium & 使用者按鍵可能在播放器seek尚未完成時建立marker & 播放器 state
machine新增 \texttt{cursor\_\allowbreak{}status=ready/seeking/stale}；只有ready
cursor可標記，其他狀態按鍵顯示阻擋原因。 \\
Medium &
多位標註者快速按同一碰撞會產生重複點，若自動去重可能誤刪真實短間隔觸球 &
只標示 \texttt{possible\_\allowbreak{}duplicate}，不靜默合併；由人工刪除/合併。 \\
Medium & 保存 overlay preset而沒有schema/version & Saved view保存
\texttt{filter\_\allowbreak{}schema\_\allowbreak{}version}與\texttt{overlay\_\allowbreak{}preset\_\allowbreak{}version}；未知欄位忽略但保留原始JSON。 \\
Medium & 歷史回放只有 media window、沒有完整 timeline cursor &
URL/route保存
\texttt{capture\_\allowbreak{}session\_\allowbreak{}id\ +\ capture\_\allowbreak{}time\_\allowbreak{}us}，切換window後仍回到相同canonical時間。 \\
\end{longtable}

\subsection{2.1.1 歷史稽核：v3.0
的實作修正}\label{211-ux6b77ux53f2ux7a3dux6838v30-ux7684ux5be6ux4f5cux4feeux6b63}

\begin{longtable}[]{@{}>{\RaggedRight\arraybackslash}p{0.18\linewidth}>{\RaggedRight\arraybackslash}p{0.34\linewidth}>{\RaggedRight\arraybackslash}p{0.42\linewidth}@{}}
\toprule\noalign{}
嚴重度 & 缺口 & 修正後基線 \\
\midrule\noalign{}
\endhead
\bottomrule\noalign{}
\endlastfoot
Critical & GraphQL 文件使用 \texttt{BigInt}，code-first builder卻註冊
\texttt{Int64} & 全部統一為 \texttt{BigInt}；Web
codegen映射為\texttt{string}。 \\
Critical & Annotation WebSocket只定義自由格式 \texttt{payload} &
每一個Z／Space／\texttt{CLOSE\_\allowbreak{}RALLY}／move／delete／reopen／Enter
command都具有strict payload schema；close必須帶target ID與strict outcome。 \\
Critical & \texttt{DvrSegment.presentationStartUs}容易被誤認為player/HLS
presentation time &
DB改名\texttt{captureStartUs/captureEndUs}，只表示整場canonical
timeline。 \\
High & SDK內附的FlatBuffers schema與canonical schema欄位／版本不一致 &
SDK
wheel強制內含與\texttt{packages/contracts/flatbuffers/overlay.fbs}完全相同的檔案，CI
byte-compare。 \\
High & callback文件使用\texttt{kind}，schema/SDK使用\texttt{status} &
Callback metadata統一使用\texttt{kind=processing/failed/completed}。 \\
High & Score ledger沒有revision序列 &
\texttt{MatchSet.scoreRevision}與\texttt{PointAward.scoreRevisionBefore/After}形成CAS
ledger；after revision在set內唯一。 \\
High & AI ContactEvent只保存key point UUID字串，沒有DB relation &
ContactEvent外鍵指向immutable
\texttt{RallySubmissionKeyPoint}；仍由domain
validator確認同一submission。 \\
High & PlaybackWindow與frame-step只存在文字範例，沒有versioned schema &
新增request/response分離的media
schemas，避免同一\texttt{oneOf}讓錯誤方向的payload也通過驗證。 \\
High & Safari原生HLS未必提供segment ID & Browser cursor不要求segment
ID；後端依window mapping與sample index解析。 \\
Medium & Coach底部導覽漏掉統計頁 &
固定五項：現場、球路、球員、統計、紀錄；頂部可返回場次選單與歷史。 \\
\end{longtable}

\subsection{2.1.2 v3.2
實作／使用稽核新增修正}\label{212-v32-ux5be6ux4f5cux4f7fux7528ux7a3dux6838ux65b0ux589eux4feeux6b63}

\begin{longtable}[]{@{}>{\RaggedRight\arraybackslash}p{0.18\linewidth}>{\RaggedRight\arraybackslash}p{0.34\linewidth}>{\RaggedRight\arraybackslash}p{0.42\linewidth}@{}}
\toprule\noalign{}
嚴重度 & 缺口 & 修正後基線 \\
\midrule\noalign{}
\endhead
\bottomrule\noalign{}
\endlastfoot
Critical & Playback request與response放在同一個\texttt{oneOf} schema &
拆成\texttt{playback-window-request}／\texttt{descriptor}、\texttt{playback-cursor}／\texttt{resolved-media-anchor}、\texttt{frame-step-request}／\texttt{canonical-frame-anchor}；OpenAPI每個方向只引用正確shape。 \\
Critical & Annotation command由client傳\texttt{user\_\allowbreak{}id}並被當成可信身份
& 使用HTTP/WebSocket upgrade完成身份驗證並綁定device
session；command不再攜帶可偽造的\texttt{user\_\allowbreak{}id}。 \\
Critical & CREATE\_\allowbreak{}SERVICE要求server既有rally ID，但Z本身又負責建立rally
& 前端在Z按下時產生client UUID作\texttt{rally\_\allowbreak{}id}與idempotency
key；server在同一transaction接受／建立draft，不另做會改變時間點的前置API。 \\
High & ACK只有自由格式\texttt{snapshot\_\allowbreak{}patch}，前端拿不到created key
point／submission ID &
\texttt{command\_\allowbreak{}ack.effects}明確回傳created／terminal／deleted key
point、submission、annotation status與score state；另廣播typed rally
snapshot。 \\
High & PlaybackWindow只回目前window範圍 &
Descriptor新增整場\texttt{timeline\_\allowbreak{}capture\_\allowbreak{}start\_\allowbreak{}us/end\_\allowbreak{}us}及\texttt{has\_\allowbreak{}more\_\allowbreak{}before/after}，全場進度條不需把整場media塞進client。 \\
High & AI Provider capabilities與202 Accepted response沒有正式schema &
新增\texttt{capabilities.schema.json}與\texttt{job-accepted.schema.json}，Python
SDK也提供typed models。 \\
High & AI result只檢查數量，未保證segment連接相鄰event或A/B一致 & SDK
validator要求event/segment index連續、segment引用相鄰key
point、A/B等於對應contact event representative positions、terminal
flag一致。 \\
High & \texttt{observed}
ball可以沒有座標，而\texttt{missing}卻可殘留座標 &
observed/interpolated必須帶sample
frame與\texttt{frame\_\allowbreak{}pos}；missing兩者皆不得帶。 \\
High & Submission可保存\texttt{PENDING}，且clip
policy只有版本沒有實際roll值 &
Submission使用只含RESOLVED／UNKNOWN的enum，並snapshot\texttt{clip\_\allowbreak{}pre\_\allowbreak{}roll\_\allowbreak{}us}／\texttt{clip\_\allowbreak{}post\_\allowbreak{}roll\_\allowbreak{}us}及service/terminal
IDs。 \\
High & MediaAsset在UPLOADING狀態就強制要求checksum/size &
\texttt{byte\_\allowbreak{}length}與\texttt{sha256}在upload完成前可null；轉READY時domain
transaction必須補齊。 \\
Medium & Browser OverlayChunk內嵌完整OverlaySequence & Chunk
schema改成獨立SoA
payload，只保留chunk範圍與陣列，不重複整段metadata。 \\
\end{longtable}

\subsection{2.2 Schema
欄位來源標記}\label{22-schema-ux6b04ux4f4dux4f86ux6e90ux6a19ux8a18}

所有跨系統欄位表必須標示以下 ownership，而不是只列型別：

\begin{itemize}
\tightlist
\item
  \texttt{USER\_\allowbreak{}INPUT}：由標註者或管理者明確輸入。
\item
  \texttt{CLIENT\_\allowbreak{}OBSERVED}：瀏覽器觀測值，例如 PlaybackCursor；不可當
  authoritative。
\item
  \texttt{BACKEND\_\allowbreak{}DERIVED}：後端從mapping/assignment/ledger推導。
\item
  \texttt{PREPROCESS\_\allowbreak{}DERIVED}：裁切與sample mapping產生。
\item
  \texttt{AI\_\allowbreak{}GENERATED}：外部 AI 分析輸出。
\item
  \texttt{PASSTHROUGH}：上游建立，接收端必須原樣回傳/關聯。
\item
  \texttt{CONSTANT}：固定於版本化規格，不在每個 request重複傳送。
\item
  \texttt{OPTIONAL\_\allowbreak{}EXTENSION}：只有 capability聲明支援時才可用。
\end{itemize}

主 Agent在新增任何 public field前，必須先寫出 producer、consumer、source
marker、nullability、versioning與失敗語意。

\begin{center}\rule{0.5\linewidth}{0.5pt}\end{center}

\section{3. 三個團隊與 Schema
Ownership}\label{3-ux4e09ux500bux5718ux968aux8207-schema-ownership}

\subsection{3.1 AI 團隊}\label{31-ai-ux5718ux968a}

\subsubsection{負責}\label{ux8ca0ux8cac}

\begin{itemize}
\tightlist
\item
  以 SDK／HTTP 接收 AI Job。
\item
  使用既有 AI work 產生 tracks、contact events、participants、A/B
  segments 與 overlay。
\item
  保持所有 PASSTHROUGH 欄位不變。
\item
  呼叫中央 callback。
\end{itemize}

\subsubsection{不負責}\label{ux4e0dux8ca0ux8cac}

\begin{itemize}
\tightlist
\item
  直播採集、整場錄影、時間軸或影片裁切。
\item
  Nuxt UI、GraphQL、PostgreSQL 或 MinIO 的中央資料模型。
\item
  跨 rally 球員身份。
\item
  中央統計公式與教練頁面。
\end{itemize}

\subsubsection{擁有 Schema}\label{ux64c1ux6709-schema}

AI 團隊不單方面擁有任何 public schema。它與後端共同消費：

\begin{itemize}
\tightlist
\item
  \texttt{packages/contracts/ai/job.schema.json}
\item
  \texttt{packages/contracts/ai/result.schema.json}
\item
  \texttt{packages/contracts/ai/callback.schema.json}
\item
  \texttt{packages/contracts/flatbuffers/overlay.fbs}
\end{itemize}

Schema 由主 Agent核准；AI 團隊透過 Python SDK 使用。

\subsection{3.2
後端系統團隊}\label{32-ux5f8cux7aefux7cfbux7d71ux5718ux968a}

包含中央 API、影音採集、完整 DVR、裁切、AI orchestration、callback
ingest、MinIO、PostgreSQL、Redis 與 worker。

\subsubsection{擁有}\label{ux64c1ux6709}

\begin{itemize}
\tightlist
\item
  Prisma schema/migrations。
\item
  Pothos code-first schema modules、resolver 與產生後的 SDL。
\item
  REST/OpenAPI。
\item
  WebSocket command/event protocol server 實作。
\item
  媒體 timeline 與 PlaybackCursor resolution。
\item
  Clip preprocessing 與 AI Job 建立。
\item
  Result 驗證、正規化、分析聚合與 artifact 管理。
\item
  容器與部署。
\end{itemize}

\subsubsection{必須視為 Passthrough
的欄位}\label{ux5fc5ux9808ux8996ux70ba-passthrough-ux7684ux6b04ux4f4d}

AI request 建立後，下列值在 AI result/callback 不得改寫：

\begin{itemize}
\tightlist
\item
  \texttt{ai\_\allowbreak{}job\_\allowbreak{}id}
\item
  \texttt{rally\_\allowbreak{}id}
\item
  \texttt{rally\_\allowbreak{}submission\_\allowbreak{}id}
\item
  \texttt{annotation\_\allowbreak{}revision}
\item
  \texttt{clip\_\allowbreak{}id}
\item
  每個 \texttt{key\_\allowbreak{}point\_\allowbreak{}id}
\item
  每個 key point 的 \texttt{sequence\_\allowbreak{}index}
\item
  \texttt{marker\_\allowbreak{}kind}
\item
  \texttt{is\_\allowbreak{}terminal}
\end{itemize}

\subsection{3.3
前端設計團隊}\label{33-ux524dux7aefux8a2dux8a08ux5718ux968a}

\subsubsection{負責}\label{ux8ca0ux8cac-1}

\begin{itemize}
\tightlist
\item
  PC-first Nuxt 多頁式標註編輯器、教練端與歷史紀錄；標註端以鍵盤、滑鼠與高資訊密度操作為主要驗收。
\item
  只有教練／裁判顯示面板要求 installable iPad landscape PWA；標註端保留六個 touch action 的同 command-path parity，但不以 iPad PWA 作主要版面或驗收裝置。
\item
  Full-session timeline + lazy playback windows。
\item
  快捷鍵 state machine 與灰／綠 mask。
\item
  GraphQL query/mutation client。
\item
  Annotation WebSocket command/reconciliation。
\item
  Video + Canvas overlay 與 2D court。
\item
  Feature availability／資料品質／樣本數的正確呈現。
\end{itemize}

\subsubsection{不負責}\label{ux4e0dux8ca0ux8cac-1}

\begin{itemize}
\tightlist
\item
  自行換算 authoritative source PTS/frame。
\item
  自行從球影像座標推導 court A/B。
\item
  自行認定 action taxonomy。
\item
  直接存取 MinIO credential。
\item
  把 track ID 當 player ID。
\end{itemize}

\subsection{3.4 Producer／Consumer
Matrix}\label{34-producerconsumer-matrix}

\begin{longtable}[]{@{}>{\RaggedRight\arraybackslash}p{0.11\linewidth}>{\RaggedRight\arraybackslash}p{0.19\linewidth}>{\RaggedRight\arraybackslash}p{0.31\linewidth}>{\RaggedRight\arraybackslash}p{0.33\linewidth}@{}}
\toprule\noalign{}
Contract & Producer & Consumer & 傳輸 \\
\midrule\noalign{}
\endhead
\bottomrule\noalign{}
\endlastfoot
Match／Set／Rally structured data & Backend & Web & GraphQL \\
Annotation command & Web & Backend & WebSocket；GraphQL fallback \\
Annotation ack/snapshot/conflict & Backend & Web & WebSocket \\
Timeline availability/live edge & Backend & Web & GraphQL + WebSocket
delta \\
Playback window & Backend media & Web player & REST + HLS/fMP4/Range \\
Clip job & Backend & Worker & PostgreSQL durable job \\
AI Job & Backend AI gateway & External AI & REST JSON \\
AI callback metadata & External AI SDK & Backend & REST
JSON/multipart \\
Full per-frame overlay & External AI & Backend & FlatBuffers
multipart \\
Overlay manifest/chunks & Backend & Web & REST JSON + FlatBuffers \\
Analysis query & Backend & Web & GraphQL \\
Presence/soft lock & Backend & Web & WebSocket \\
\end{longtable}

\begin{center}\rule{0.5\linewidth}{0.5pt}\end{center}

\section{4. 最終 Repository
與技術架構}\label{4-ux6700ux7d42-repository-ux8207ux6280ux8853ux67b6ux69cb}

\subsection{4.1 Monorepo}\label{41-monorepo}

\begin{Shaded}
\begin{Highlighting}[]
\NormalTok{volleyball{-}monitoring{-}ai/}
\NormalTok{├── AGENTS.md}
\NormalTok{├── README.md}
\NormalTok{├── package.json}
\NormalTok{├── bunfig.toml}
\NormalTok{├── .env.example}
\NormalTok{├── .codex/}
\NormalTok{│   ├── config.toml}
\NormalTok{│   └── agents/}
\NormalTok{│       ├── contracts{-}sdk{-}worker.toml}
\NormalTok{│       ├── backend{-}worker.toml}
\NormalTok{│       └── luna{-}worker.toml}
\NormalTok{├── docs/}
\NormalTok{│   ├── MASTER\_\allowbreak{}IMPLEMENTATION\_\allowbreak{}SPEC.md}
\NormalTok{│   ├── SYSTEM\_\allowbreak{}SPEC\_\allowbreak{}V3\_\allowbreak{}2.pdf}
\NormalTok{│   ├── MAIN\_\allowbreak{}AGENT\_\allowbreak{}PROMPT.md}
\NormalTok{│   ├── LUNA\_\allowbreak{}WORKER\_\allowbreak{}SETUP\_\allowbreak{}PROMPT.md}
\NormalTok{│   ├── requirements{-}matrix.md}
\NormalTok{│   ├── progress.md}
\NormalTok{│   ├── open{-}decisions.md}
\NormalTok{│   └── adr/}
\NormalTok{├── packages/}
\NormalTok{│   ├── contracts/}
\NormalTok{│   │   ├── annotation/}
\NormalTok{│   │   ├── media/}
\NormalTok{│   │   ├── ai/}
\NormalTok{│   │   ├── openapi/}
\NormalTok{│   │   ├── flatbuffers/}
\NormalTok{│   │   ├── graphql/}
\NormalTok{│   │   └── fixtures/}
\NormalTok{│   └── db/}
\NormalTok{│       ├── prisma/schema.prisma}
\NormalTok{│       ├── prisma/migrations/}
\NormalTok{│       └── src/}
\NormalTok{├── sdk/}
\NormalTok{│   ├── pyproject.toml}
\NormalTok{│   ├── src/volleyball\_\allowbreak{}monitoring\_\allowbreak{}ai/}
\NormalTok{│   ├── examples/}
\NormalTok{│   └── tests/}
\NormalTok{├── web/}
\NormalTok{│   ├── app/}
\NormalTok{│   ├── graphql/}
\NormalTok{│   ├── tests/}
\NormalTok{│   └── nuxt.config.ts}
\NormalTok{├── server/}
\NormalTok{│   ├── src/graphql/}
\NormalTok{│   ├── src/domain/}
\NormalTok{│   ├── src/rest/}
\NormalTok{│   ├── src/realtime/}
\NormalTok{│   └── tests/}
\NormalTok{├── worker/}
\NormalTok{│   ├── src/roles/}
\NormalTok{│   └── tests/}
\NormalTok{├── examples/}
\NormalTok{│   └── fake\_\allowbreak{}ai\_\allowbreak{}provider/}
\NormalTok{├── infra/}
\NormalTok{│   ├── compose.yaml}
\NormalTok{│   ├── docker/}
\NormalTok{│   ├── mediamtx/}
\NormalTok{│   └── k8s/}
\NormalTok{├── scripts/}
\NormalTok{└── .github/workflows/ci.yml}
\end{Highlighting}
\end{Shaded}

採 top-level \texttt{server/} 與 \texttt{worker/}，避免同一文件同時出現
\texttt{backend/api}、\texttt{backend/worker} 與另一套 root 路徑。

\subsection{4.2 Web}\label{42-web}

\begin{itemize}
\tightlist
\item
  Nuxt 4、Vue 3、TypeScript strict。
\item
  UI 可以參考舊 MVP 使用的 Nuxt／Vant／Tailwind／Motion／Lucide
  組合，但不得沿用其 domain flow、memory store 或手動事件模型。
\item
  Vant 4 用於教練端 iPad/mobile interaction 與必要的 touch parity；PC-first 標註器使用高資訊密度 desktop layout；Tailwind 4 建立 design
  tokens；Motion for Vue 只用於必要狀態與球路動畫。
\item
  GraphQL Code Generator \texttt{client-preset} 由產生後 SDL 與前端
  operations 建立 \texttt{TypedDocumentNode}；禁止手寫另一套 domain
  interface。
\item
  Video element 上疊透明 Canvas；影片本身不逐 frame 畫入 Canvas。
\end{itemize}

\subsection{4.3 API：Fastify + Yoga + Pothos +
Prisma}\label{43-apifastify--yoga--pothos--prisma}

中央 API 採獨立 TypeScript service，不把長生命週期 WebSocket、callback
大檔與 FFmpeg 工作塞入 Nuxt/Nitro。

\begin{itemize}
\tightlist
\item
  \textbf{Fastify}：HTTP server、REST routes、WebSocket
  upgrade、lifecycle 與 health endpoints。
\item
  \textbf{GraphQL Yoga}：GraphQL over HTTP 執行層。
\item
  \textbf{Pothos}：唯一 GraphQL code-first schema
  source；\texttt{builder.toSchema()} 產生標準 \texttt{GraphQLSchema}
  交給 Yoga。
\item
  \textbf{Pothos Prisma plugin}：建立 Prisma-backed
  object/relations，並利用 selection information 避免常見 N+1。
\item
  \textbf{Prisma 7 + PostgreSQL}：canonical persistence；使用
  \texttt{prisma.config.ts}、\texttt{prisma-client} generator 與
  PostgreSQL driver adapter。
\item
  \textbf{GraphQL Code Generator}：從 Pothos schema 輸出 checked-in
  SDL，再由 client preset 產生 Nuxt typed operations。
\item
  \textbf{JSON Schema/Zod}：REST、WebSocket 與 callback validation。
\item
  \textbf{Redis}：presence、soft lock、pub/sub fan-out；不得放 durable
  canonical state。
\item
  \textbf{pg-boss 或等價 PostgreSQL-backed queue}：durable media/AI
  jobs。
\end{itemize}

\subsubsection{Code-first
單一真實來源}\label{code-first-ux55aeux4e00ux771fux5be6ux4f86ux6e90}

\begin{Shaded}
\begin{Highlighting}[]
\NormalTok{server/src/graphql/**/*.ts (Pothos source)}
\NormalTok{    → builder.toSchema()}
\NormalTok{    → packages/contracts/graphql/schema.graphql}
\NormalTok{    → GraphQL Inspector / CI contract diff}
\NormalTok{    → web/codegen.ts}
\NormalTok{    → web/app/gql/** typed documents}
\end{Highlighting}
\end{Shaded}

規則：

\begin{enumerate}
\def\labelenumi{\arabic{enumi}.}
\tightlist
\item
  不可手改 \texttt{graphql/schema.graphql}。
\item
  Prisma model 不等於 GraphQL public contract；只有 Pothos 明確 expose
  的欄位會公開。
\item
  GraphQL resolver 必須呼叫 domain/service layer，不能繞過
  annotation、time mapping 或 submission invariant。
\item
  高頻 annotation collaboration 使用專用 WebSocket command
  protocol；GraphQL mutation只作 fallback，不強迫使用 GraphQL
  subscription。
\item
  Binary、影片、FlatBuffers、AI multipart callback 不進 GraphQL。
\end{enumerate}

\subsection{4.4 Worker}\label{44-worker}

同一個 worker image 可以依 command/environment 執行不同角色：

\begin{itemize}
\tightlist
\item
  \texttt{media-indexer}
\item
  \texttt{playback-packager}
\item
  \texttt{clip-worker}
\item
  \texttt{ai-dispatcher}
\item
  \texttt{analysis-ingest}
\end{itemize}

它們共享 \texttt{packages/db} 與 \texttt{packages/contracts}，但每個
worker 只能 claim 自己的 job type。FFmpeg subprocess 必須有
timeout、cancel、stderr capture、temporary directory quota 與 idempotent
output key。

\subsection{4.5 Media}\label{45-media}

\begin{itemize}
\tightlist
\item
  MediaMTX：接收來源、live LL-HLS/WebRTC、fMP4 recording 與基礎
  playback。
\item
  FFmpeg/ffprobe：索引、archive playback window、canonical clip、preview
  與 sample mapping。
\item
  MinIO：完整錄影片段、playback windows、clips、AI artifacts。
\item
  PostgreSQL：capture epochs、segment index、available ranges、mapping
  version、job state。
\item
  標註用播放 profile 必須是可建立 deterministic media-time →
  capture-time mapping 的單一 rendition；若未來開放 ABR，每個 rendition
  必須各自有 mapping，不能共用 \texttt{currentTime\ ×\ FPS}。
\end{itemize}

\subsection{4.6 固定 Framework／Library
選型}\label{46-ux56faux5b9a-frameworklibrary-ux9078ux578b}

第一版 starter 固定以下基線，避免主 Agent與三個 subagent各自選一套：

\begin{longtable}[]{@{}>{\RaggedRight\arraybackslash}p{0.18\linewidth}>{\RaggedRight\arraybackslash}p{0.34\linewidth}>{\RaggedRight\arraybackslash}p{0.42\linewidth}@{}}
\toprule\noalign{}
層 & 選型 & 用途 \\
\midrule\noalign{}
\endhead
\bottomrule\noalign{}
\endlastfoot
Runtime／Package & Bun \texttt{1.3.14} & workspace
install、script、server/worker runtime、Docker image \\
Web & Nuxt 4、Vue 3、TypeScript & PC-first Annotation workstation與Coach多頁PWA \\
UI & Vant 4、Tailwind CSS、Motion-V、Lucide、VueUse &
觸控元件、版面、物理動畫、icon與composable \\
PWA & \texttt{@vite-pwa/nuxt} & Coach/viewer standalone manifest、service
worker、更新提示、iPad加入主畫面 \\
GraphQL Client & URQL + GraphQL Code Generator &
讀取結構化domain與generated type \\
Media Client & native Safari HLS + HLS.js fallback & live/archive
HLS、bounded client buffer \\
API Server & Fastify 5 + GraphQL Yoga 5 & 單一HTTP
server內同時掛GraphQL、REST與WebSocket upgrade \\
Code-first GraphQL & Pothos + Prisma plugin & TypeScript code-first
schema；SDL只作CI產物 \\
Data & Prisma 7 + PostgreSQL 17 & domain、revision、media
index、submission、job與analysis metadata \\
Job／Realtime & pg-boss + Redis & durable worker job與WS
fan-out／presence \\
Object Storage & MinIO S3 & raw recording、DVR、canonical
clip、analysis與overlay artifacts \\
Media & MediaMTX + FFmpeg/ffprobe &
ingest、LL-HLS、完整錄製、index、playback window與裁切 \\
Edge & Traefik 3.7.9 & 本地Docker同源routing、WebSocket與HLS反向代理 \\
AI Integration & Python 3.11+ SDK & 外部AI
team驗證Job／Result、下載clip、callback與FlatBuffers \\
\end{longtable}

不採Caddy。瀏覽器使用Traefik同源路徑
\texttt{/graphql}、\texttt{/api/v1}、\texttt{/ws}、\texttt{/hls}，避免把\texttt{localhost}編入任何 browser
bundle。MinIO credentials、AI provider bearer與callback
token不可進前端。

\subsection{4.7 Coach／Viewer iPad PWA
固定行為}\label{47-ipad-pwa-ux56faux5b9aux884cux70ba}

本節只約束教練／裁判顯示面板。PC-first annotation editor 不得為了 iPad 單欄版面犧牲 timeline、selected-point inspector、快捷鍵與多區工作流；共用 media／command 元件仍須維持 touch parity 與可及性。

\begin{itemize}
\tightlist
\item
  PWA \texttt{display=standalone}、landscape-first、支援safe-area與Apple
  touch icon。
\item
  iPad安裝與Service Worker必須使用可信任的HTTPS origin；LAN上的純HTTP
  IP只能作一般瀏覽器預覽，不得當作PWA驗收環境。
\item
  \texttt{orientation=landscape}只是一個manifest偏好，UI仍要在直向時顯示「請旋轉iPad」提示，不得依賴瀏覽器一定鎖定方向。
\item
  頂部狀態列與底部五分頁導覽在standalone模式仍完整可用。
\item
  Service worker可cache app shell與已讀分析資料；GraphQL
  mutation、annotation WebSocket、live/DVR media、callback與signed
  artifact一律NetworkOnly。
\item
  Offline時可以查看已cache的歷史資料，但未獲server
  ACK的annotation只能顯示local pending，不得冒充canonical revision。
\item
  PWA更新使用auto-update提示；執行中的annotation不可無提示強制reload。
\end{itemize}

\subsection{4.7.1
本地Docker與iPad可信任HTTPS}\label{471-ux672cux5730dockerux8207ipadux53efux4fe1ux4efbhttps}

本地Docker仍由Traefik統一入口。Desktop可先用HTTP檢查服務，但iPad「加入主畫面」、Service
Worker與正式PWA驗收必須使用可信任HTTPS：

\begin{enumerate}
\def\labelenumi{\arabic{enumi}.}
\tightlist
\item
  為LAN hostname或IP建立本地開發憑證；starter提供
  \texttt{scripts/generate-local-tls.sh}，預設使用 \texttt{mkcert}。
\item
  將本地CA安裝到iPad，並在iOS的憑證信任設定中啟用完整信任。
\item
  Traefik的\texttt{web} entrypoint只做HTTP到HTTPS
  redirect；PWA、GraphQL、REST、WebSocket與HLS全部走\texttt{websecure}同源入口。
\item
  憑證私鑰與本地CA不得提交Git；repository只保存產生腳本與範例設定。
\item
  若部署到正式DNS，改用組織既有PKI或ACME，不沿用本地CA。
\end{enumerate}

\subsection{4.8 Serialization 分工}\label{48-serialization-ux5206ux5de5}

\begin{longtable}[]{@{}>{\RaggedRight\arraybackslash}p{0.18\linewidth}>{\RaggedRight\arraybackslash}p{0.34\linewidth}>{\RaggedRight\arraybackslash}p{0.42\linewidth}@{}}
\toprule\noalign{}
類型 & 格式 & 原因 \\
\midrule\noalign{}
\endhead
\bottomrule\noalign{}
\endlastfoot
一般資料查詢／mutation & GraphQL JSON & 關聯查詢、code-first
contract、Nuxt type safety \\
Media／binary／callback & REST & range、stream、multipart 與大型 payload
不適合 GraphQL \\
即時 annotation command／status & WebSocket JSON & command
ID、revision、低延遲、重連 \\
AI event summary & JSON Schema & 可讀、易驗證、資料量有限 \\
逐幀 tracking/ball/action & FlatBuffers & 避免巨大 JSON object、降低
iPad parse/GC 成本 \\
\end{longtable}

\begin{center}\rule{0.5\linewidth}{0.5pt}\end{center}

\section{5. 容器與
Infrastructure}\label{5-ux5bb9ux5668ux8207-infrastructure}

\subsection{5.1 Baseline Containers}\label{51-baseline-containers}

\begin{longtable}[]{@{}>{\RaggedRight\arraybackslash}p{0.18\linewidth}>{\RaggedRight\arraybackslash}p{0.34\linewidth}>{\RaggedRight\arraybackslash}p{0.42\linewidth}@{}}
\toprule\noalign{}
Container & 責任 & 是否有 durable state \\
\midrule\noalign{}
\endhead
\bottomrule\noalign{}
\endlastfoot
\texttt{web} & Nuxt UI/SSR/SPA & 否 \\
\texttt{server} & GraphQL、REST、WS、auth、domain command & 否 \\
\texttt{worker-media-indexer} / \texttt{worker-playback} /
\texttt{worker-clip} & index、playback window、clip、FFmpeg & local
spool 只是暫存 \\
\texttt{worker-ai-dispatcher} / \texttt{worker-analysis-ingest} &
submit、retry、callback ingest／normalize & 否 \\
\texttt{mediamtx} & ingest/live/record & recording spool 暫存 \\
\texttt{postgres} & canonical metadata、revision、job、analysis & 是 \\
\texttt{minio} & media、clip、analysis artifact & 是 \\
\texttt{redis} & presence、soft lock、fan-out & 否，可重建 \\
\texttt{minio-init} & 建 buckets/lifecycle & 否 \\
\end{longtable}

Baseline 不建立 AI 推論容器。開發環境可另啟動
\texttt{fake-ai-provider}，它只驗證 contract 與產生 fixtures。

\subsection{5.2 MinIO Buckets}\label{52-minio-buckets}

\begin{Shaded}
\begin{Highlighting}[]
\NormalTok{raw{-}media/}
\NormalTok{  matches/\{match\_\allowbreak{}id\}/captures/\{capture\_\allowbreak{}id\}/segments/\{discontinuity\}/\{segment\_\allowbreak{}id\}.mp4}

\NormalTok{playback/}
\NormalTok{  matches/\{match\_\allowbreak{}id\}/captures/\{capture\_\allowbreak{}id\}/windows/\{window\_\allowbreak{}id\}/...}

\NormalTok{rally{-}media/}
\NormalTok{  matches/\{match\_\allowbreak{}id\}/rallies/\{rally\_\allowbreak{}id\}/submissions/\{submission\_\allowbreak{}id\}/}
\NormalTok{    canonical.mp4}
\NormalTok{    preview.mp4}
\NormalTok{    timing{-}manifest.json}

\NormalTok{analysis{-}artifacts/}
\NormalTok{  matches/\{match\_\allowbreak{}id\}/rallies/\{rally\_\allowbreak{}id\}/runs/\{analysis\_\allowbreak{}run\_\allowbreak{}id\}/}
\NormalTok{    result.json}
\NormalTok{    overlay{-}sequence.fb}
\NormalTok{    overlay{-}manifest.json}
\NormalTok{    chunks/\{index\}.fb}
\end{Highlighting}
\end{Shaded}

物件 key 不作資料庫主鍵。DB 永遠保存 bucket、object key、checksum、byte
size、MIME、schema version 與 lifecycle state。

\subsection{5.3 Durable Job}\label{53-durable-job}

\begin{itemize}
\tightlist
\item
  Clip、playback package、AI submit、artifact ingest 都必須是 DB-backed
  durable job。
\item
  Worker claim 使用 queue library 的 PostgreSQL
  transaction/lease；不得只用 Redis list。
\item
  每個 job 必須有：

  \begin{itemize}
  \tightlist
  \item
    \texttt{id}
  \item
    \texttt{job\_\allowbreak{}type}
  \item
    \texttt{status}
  \item
    \texttt{deduplication\_\allowbreak{}key}
  \item
    \texttt{attempt\_\allowbreak{}count}
  \item
    \texttt{max\_\allowbreak{}attempts}
  \item
    \texttt{available\_\allowbreak{}at}
  \item
    \texttt{leased\_\allowbreak{}until}
  \item
    \texttt{last\_\allowbreak{}error\_\allowbreak{}code/message}
  \item
    \texttt{created\_\allowbreak{}at/updated\_\allowbreak{}at/completed\_\allowbreak{}at}
  \end{itemize}
\item
  Worker crash 後 lease 到期可由另一 worker 接手。
\end{itemize}

\begin{center}\rule{0.5\linewidth}{0.5pt}\end{center}

\section{6. API
邊界：GraphQL、REST、WebSocket、FlatBuffers}\label{6-api-ux908aux754cgraphqlrestwebsocketflatbuffers}

\subsection{6.1 GraphQL 適用}\label{61-graphql-ux9069ux7528}

\begin{itemize}
\tightlist
\item
  Match、team、player、roster、set。
\item
  Capture session metadata 與完整 timeline availability。
\item
  Rally list/detail、submission history、job status。
\item
  Analysis、filter、statistics、saved view。
\item
  管理設定與 track identity assignment。
\item
  Annotation snapshot/refetch 與低頻 fallback mutation。
\end{itemize}

GraphQL 不傳：

\begin{itemize}
\tightlist
\item
  MP4、HLS segment、FlatBuffers、callback multipart。
\item
  每 frame tracking array。
\item
  MinIO credential。
\end{itemize}

\subsection{6.2 REST 適用}\label{62-rest-ux9069ux7528}

\begin{Shaded}
\begin{Highlighting}[]
\NormalTok{POST /v1/media/playback{-}windows}
\NormalTok{GET  /v1/media/playback{-}windows/\{id\}}
\NormalTok{GET  /v1/media/playback{-}windows/\{id\}/manifest.m3u8}
\NormalTok{GET  /v1/media/playback{-}windows/\{id\}/media.mp4}
\NormalTok{POST /v1/media/resolve{-}cursor}
\NormalTok{POST /v1/media/frame{-}step}

\NormalTok{GET  /v1/ai/jobs/\{ai\_\allowbreak{}job\_\allowbreak{}id\}/clip}
\NormalTok{POST /v1/ai/callback/\{ai\_\allowbreak{}job\_\allowbreak{}id\}}
\NormalTok{GET  /v1/analysis{-}runs/\{id\}/overlay{-}manifest}
\NormalTok{GET  /v1/analysis{-}runs/\{id\}/overlay{-}chunks/\{index\}}

\NormalTok{GET  /health/live}
\NormalTok{GET  /health/ready}
\end{Highlighting}
\end{Shaded}

\subsection{6.3 WebSocket 適用}\label{63-websocket-ux9069ux7528}

單一連線可加入 match room，訊息分為：

\begin{itemize}
\tightlist
\item
  \texttt{annotation.command}
\item
  \texttt{annotation.ack}
\item
  \texttt{annotation.conflict}
\item
  \texttt{annotation.snapshot\_\allowbreak{}changed}
\item
  \texttt{presence.update}
\item
  \texttt{soft\_\allowbreak{}lock.update}
\item
  \texttt{timeline.live\_\allowbreak{}edge}
\item
  \texttt{timeline.segment\_\allowbreak{}added}
\item
  \texttt{processing.status\_\allowbreak{}changed}
\item
  \texttt{analysis.ready}
\end{itemize}

WebSocket 不廣播完整逐幀 overlay，也不應每次小操作廣播整場大型
snapshot。

\subsection{6.4 FlatBuffers 適用}\label{64-flatbuffers-ux9069ux7528}

只負責：

\begin{itemize}
\tightlist
\item
  每 frame 的 track observation。
\item
  frame bbox、foot frame position、court position。
\item
  ball frame position。
\item
  optional action dictionary index。
\item
  flags／optional quantized confidence。
\end{itemize}

Event、rally、team、player、filter 仍使用 JSON/GraphQL。

\begin{center}\rule{0.5\linewidth}{0.5pt}\end{center}

\section{7.
統一時間與座標定義}\label{7-ux7d71ux4e00ux6642ux9593ux8207ux5ea7ux6a19ux5b9aux7fa9}

\subsection{7.1
統一時間欄位}\label{71-ux7d71ux4e00ux6642ux9593ux6b04ux4f4d}

時間分成四層，欄位名稱不得互換：

\begin{enumerate}
\def\labelenumi{\arabic{enumi}.}
\tightlist
\item
  \textbf{來源 epoch
  時間}：\texttt{capture\_\allowbreak{}epoch\_\allowbreak{}id\ +\ source\_\allowbreak{}pts\ +\ source\_\allowbreak{}time\_\allowbreak{}base}。來源重連或
  PTS 重置時建立新 epoch；raw PTS 只在 epoch 內有意義。
\item
  \textbf{整場 canonical 時間}：\texttt{capture\_\allowbreak{}time\_\allowbreak{}us}，從 capture
  session 起點單調增加，資料庫使用 PostgreSQL \texttt{BIGINT}，wire
  使用十進位字串。
\item
  \textbf{整場 canonical frame}：\texttt{capture\_\allowbreak{}frame\_\allowbreak{}index}，由
  sample index 建立；瀏覽器不得以 \texttt{currentTime\ ×\ FPS} 推算。
\item
  \textbf{裁切片段局部時間}：\texttt{clip\_\allowbreak{}time\_\allowbreak{}us}、\texttt{clip\_\allowbreak{}pts}、\texttt{clip\_\allowbreak{}frame\_\allowbreak{}index}，均由前處理服務產生，AI
  只使用這組局部時間。
\end{enumerate}

Browser 只送 \texttt{PlaybackCursor} 觀測值。後端依
\texttt{playback\_\allowbreak{}window\_\allowbreak{}id\ +\ mapping\_\allowbreak{}version\ +\ player\_\allowbreak{}media\_\allowbreak{}time\_\allowbreak{}us}
查 window mapping 與 sample index，解析為 authoritative
anchor。\texttt{requestVideoFrameCallback().mediaTime}
是首選，\texttt{currentTime} 只作 fallback；兩者都仍是 client
observation，不直接寫進 key point canonical 欄位。

GraphQL \texttt{Int} 只有 32-bit，不能承載 microseconds、PTS、frame
count 或 byte size。所有 64-bit wire 欄位使用自訂 \texttt{BigInt}
scalar並序列化為 decimal string；前端 codegen映射為
\texttt{string}，不可映射為 JavaScript \texttt{number}。

\subsection{7.2 Frame Position}\label{72-frame-position}

\begin{Shaded}
\begin{Highlighting}[]
\FunctionTok{\{}\DataTypeTok{"x"}\FunctionTok{:} \DecValTok{0}\ErrorTok{.}\DecValTok{42}\FunctionTok{,} \DataTypeTok{"y"}\FunctionTok{:} \DecValTok{0}\ErrorTok{.}\DecValTok{63}\FunctionTok{\}}
\end{Highlighting}
\end{Shaded}

\begin{itemize}
\tightlist
\item
  \texttt{frame\_\allowbreak{}pos} 原點是影像左上。
\item
  x/y 通常位於 \texttt{{[}0,1{]}}。
\item
  \texttt{frame\_\allowbreak{}bbox} 使用 \texttt{x1/y1/x2/y2} 正規化。
\item
  前端按實際 video content rect 換算，不可把 letterbox/pillarbox
  也當影像區。
\end{itemize}

\subsection{7.3 Court Position}\label{73-court-position}

\begin{Shaded}
\begin{Highlighting}[]
\FunctionTok{\{}\DataTypeTok{"x"}\FunctionTok{:} \DecValTok{0}\ErrorTok{.}\DecValTok{31}\FunctionTok{,} \DataTypeTok{"y"}\FunctionTok{:} \DecValTok{0}\ErrorTok{.}\DecValTok{72}\FunctionTok{\}}
\end{Highlighting}
\end{Shaded}

固定唯一規格：

\begin{itemize}
\tightlist
\item
  \texttt{x=0}：左側端線；\texttt{x=1}：右側端線，對應 18 m 長度。
\item
  \texttt{y=0}：canonical
  俯視球場圖的上側邊線；\texttt{y=1}：下側邊線，對應 9 m
  寬度。此方向不依賴攝影機畫面。
\item
  球場內 \texttt{{[}0,1{]}\ ×\ {[}0,1{]}}。
\item
  發球區、救球區、出界點允許 \texttt{\textless{}0} 或
  \texttt{\textgreater{}1}。
\item
  不因換邊翻轉 canonical data；前端可依教練視角 transform。
\item
  公尺值需要時：\texttt{x\_\allowbreak{}m\ =\ x\ *\ 18}、\texttt{y\_\allowbreak{}m\ =\ y\ *\ 9}。
\item
  A/B 由 AI 已處理的 footprint proxy／terminal representative point
  產生；中央與前端不再投影。
\end{itemize}

\begin{center}\rule{0.5\linewidth}{0.5pt}\end{center}

\section{8. Full-session Live DVR 與 Lazy
Playback}\label{8-full-session-live-dvr-ux8207-lazy-playback}

\subsection{8.1 使用者體驗}\label{81-ux4f7fux7528ux8005ux9ad4ux9a57}

\begin{itemize}
\tightlist
\item
  標註端進入場次後預設位於 live edge，持續播放聲音與影像。
\item
  底部 timeline 範圍是整個 capture session，從第一段可用影像到 live
  edge。
\item
  使用者可像 YouTube Live 一樣拖到任何已保存時間。
\item
  遠端歷史影片不一次下載整場；播放器只保留目前 window 與前後相鄰
  window。
\item
  使用者停留在歷史回放時，伺服器仍持續錄 live；WebSocket 持續更新 live
  edge。
\item
  按「回到現場」時切回 live rolling window。
\item
  若某區段缺檔，timeline 顯示 gap，不允許在 gap 建 marker。
\end{itemize}

\subsection{8.2 Server Timeline}\label{82-server-timeline}

GraphQL \texttt{captureTimeline} 回傳 metadata，而不是 media bytes：

\begin{Shaded}
\begin{Highlighting}[]
\FunctionTok{\{}
  \DataTypeTok{"capture\_\allowbreak{}session\_\allowbreak{}id"}\FunctionTok{:} \StringTok{"..."}\FunctionTok{,}
  \DataTypeTok{"timeline\_\allowbreak{}version"}\FunctionTok{:} \DecValTok{87}\FunctionTok{,}
  \DataTypeTok{"capture\_\allowbreak{}start\_\allowbreak{}time\_\allowbreak{}us"}\FunctionTok{:} \StringTok{"0"}\FunctionTok{,}
  \DataTypeTok{"live\_\allowbreak{}edge\_\allowbreak{}capture\_\allowbreak{}time\_\allowbreak{}us"}\FunctionTok{:} \StringTok{"7132456000"}\FunctionTok{,}
  \DataTypeTok{"available\_\allowbreak{}ranges"}\FunctionTok{:} \OtherTok{[}
    \FunctionTok{\{}\DataTypeTok{"start\_\allowbreak{}us"}\FunctionTok{:} \StringTok{"0"}\FunctionTok{,} \DataTypeTok{"end\_\allowbreak{}us"}\FunctionTok{:} \StringTok{"2410000000"}\FunctionTok{,} \DataTypeTok{"discontinuity"}\FunctionTok{:} \DecValTok{0}\FunctionTok{\}}\OtherTok{,}
    \FunctionTok{\{}\DataTypeTok{"start\_\allowbreak{}us"}\FunctionTok{:} \StringTok{"2416500000"}\FunctionTok{,} \DataTypeTok{"end\_\allowbreak{}us"}\FunctionTok{:} \StringTok{"7132456000"}\FunctionTok{,} \DataTypeTok{"discontinuity"}\FunctionTok{:} \DecValTok{1}\FunctionTok{\}}
  \OtherTok{]}
\FunctionTok{\}}
\end{Highlighting}
\end{Shaded}

UI 可用這些 range 繪完整時間軸；不需要先下載所有 segment list。

\subsection{8.3 Playback Window 與完整
DVR}\label{83-playback-window-ux8207ux5b8cux6574-dvr}

Server 保存整場，client 只取得 bounded playback window。Live manifest 與
archive window 都映射到同一條 \texttt{capture\_\allowbreak{}time\_\allowbreak{}us}；播放器的
\texttt{currentTime} 原點可以不同，因此每個 descriptor 必須攜帶 mapping
origin。

Request：

\begin{Shaded}
\begin{Highlighting}[]
\FunctionTok{\{}
  \DataTypeTok{"capture\_\allowbreak{}session\_\allowbreak{}id"}\FunctionTok{:} \StringTok{"018f..."}\FunctionTok{,}
  \DataTypeTok{"mode"}\FunctionTok{:} \StringTok{"archive"}\FunctionTok{,}
  \DataTypeTok{"target\_\allowbreak{}capture\_\allowbreak{}time\_\allowbreak{}us"}\FunctionTok{:} \StringTok{"2823123000"}\FunctionTok{,}
  \DataTypeTok{"requested\_\allowbreak{}back\_\allowbreak{}us"}\FunctionTok{:} \StringTok{"90000000"}\FunctionTok{,}
  \DataTypeTok{"requested\_\allowbreak{}forward\_\allowbreak{}us"}\FunctionTok{:} \StringTok{"120000000"}
\FunctionTok{\}}
\end{Highlighting}
\end{Shaded}

Response：

\begin{Shaded}
\begin{Highlighting}[]
\FunctionTok{\{}
  \DataTypeTok{"playback\_\allowbreak{}window\_\allowbreak{}id"}\FunctionTok{:} \StringTok{"0190..."}\FunctionTok{,}
  \DataTypeTok{"capture\_\allowbreak{}session\_\allowbreak{}id"}\FunctionTok{:} \StringTok{"018f..."}\FunctionTok{,}
  \DataTypeTok{"mode"}\FunctionTok{:} \StringTok{"archive"}\FunctionTok{,}
  \DataTypeTok{"mapping\_\allowbreak{}version"}\FunctionTok{:} \DecValTok{3}\FunctionTok{,}
  \DataTypeTok{"window\_\allowbreak{}capture\_\allowbreak{}start\_\allowbreak{}us"}\FunctionTok{:} \StringTok{"2733123000"}\FunctionTok{,}
  \DataTypeTok{"window\_\allowbreak{}capture\_\allowbreak{}end\_\allowbreak{}us"}\FunctionTok{:} \StringTok{"2943123000"}\FunctionTok{,}
  \DataTypeTok{"presentation\_\allowbreak{}origin\_\allowbreak{}capture\_\allowbreak{}us"}\FunctionTok{:} \StringTok{"2733123000"}\FunctionTok{,}
  \DataTypeTok{"target\_\allowbreak{}player\_\allowbreak{}media\_\allowbreak{}time\_\allowbreak{}us"}\FunctionTok{:} \StringTok{"90000000"}\FunctionTok{,}
  \DataTypeTok{"manifest\_\allowbreak{}url"}\FunctionTok{:} \StringTok{"/api/v1/media/playback{-}windows/0190.../manifest.m3u8"}\FunctionTok{,}
  \DataTypeTok{"expires\_\allowbreak{}at"}\FunctionTok{:} \StringTok{"2026{-}08{-}07T05:00:00Z"}
\FunctionTok{\}}
\end{Highlighting}
\end{Shaded}

\texttt{capture\_\allowbreak{}time\_\allowbreak{}us\ =\ presentation\_\allowbreak{}origin\_\allowbreak{}capture\_\allowbreak{}us\ +\ player\_\allowbreak{}media\_\allowbreak{}time\_\allowbreak{}us}
只是一階 mapping；後端仍須以 sample index吸附到實際 frame。HLS
playlist可使用 program date time協助跨 rendition／窗口對齊，但
canonical資料仍以自己的 capture/sample index為準。

預設 window可設 3--5 分鐘，這只是 client loading策略，不是回放限制。接近
window邊界時 lazy prefetch下一個 window；使用者在歷史回放時，live錄製與
live-edge更新持續進行。

\subsection{8.4 Client Buffer 策略}\label{84-client-buffer-ux7b56ux7565}

\begin{itemize}
\tightlist
\item
  \texttt{live}：使用 LL-HLS rolling manifest，只保存最近幾分鐘的
  segment references。
\item
  \texttt{archive}：下載固定 window；接近 20--30 秒邊界時 prefetch 相鄰
  window。
\item
  記憶體只保留 current/previous/next window metadata；已遠離的
  MediaSource buffer 與 overlay chunk 必須釋放。
\item
  不在背景維持第二支 live video；以 WebSocket live edge
  更新替代，避免雙倍頻寬。
\item
  回到 live 時再 load live manifest。
\end{itemize}

\subsection{8.5 Playback Cursor 與 key point
對齊}\label{85-playback-cursor-ux8207-key-point-ux5c0dux9f4a}

前端在鍵盤或觸控按鈕被觸發時，送「最後一個真正呈現的影片 frame」觀測：

\begin{Shaded}
\begin{Highlighting}[]
\FunctionTok{\{}
  \DataTypeTok{"playback\_\allowbreak{}window\_\allowbreak{}id"}\FunctionTok{:} \StringTok{"0190..."}\FunctionTok{,}
  \DataTypeTok{"mapping\_\allowbreak{}version"}\FunctionTok{:} \DecValTok{3}\FunctionTok{,}
  \DataTypeTok{"player\_\allowbreak{}media\_\allowbreak{}time\_\allowbreak{}us"}\FunctionTok{:} \StringTok{"90167234"}\FunctionTok{,}
  \DataTypeTok{"observation\_\allowbreak{}source"}\FunctionTok{:} \StringTok{"request\_\allowbreak{}video\_\allowbreak{}frame\_\allowbreak{}callback"}\FunctionTok{,}
  \DataTypeTok{"presented\_\allowbreak{}frames"}\FunctionTok{:} \StringTok{"1842"}\FunctionTok{,}
  \DataTypeTok{"seek\_\allowbreak{}generation"}\FunctionTok{:} \DecValTok{27}\FunctionTok{,}
  \DataTypeTok{"cursor\_\allowbreak{}status"}\FunctionTok{:} \StringTok{"ready"}
\FunctionTok{\}}
\end{Highlighting}
\end{Shaded}

\begin{itemize}
\tightlist
\item
  \texttt{observation\_\allowbreak{}source} 可為
  \texttt{request\_\allowbreak{}video\_\allowbreak{}frame\_\allowbreak{}callback} 或
  \texttt{current\_\allowbreak{}time\_\allowbreak{}fallback}。
\item
  Safari 原生 HLS不保證可暴露 segment ID，因此 \texttt{segment\_\allowbreak{}id}
  不可設為 browser必填欄位。
\item
  \texttt{cursor\_\allowbreak{}status=seeking/stale/gap} 時，Z/Space必須被 UI阻擋。
\item
  \texttt{seek\_\allowbreak{}generation} 防止 seek完成前的舊 callback被誤用。
\end{itemize}

後端解析後回：

\begin{Shaded}
\begin{Highlighting}[]
\FunctionTok{\{}
  \DataTypeTok{"capture\_\allowbreak{}epoch\_\allowbreak{}id"}\FunctionTok{:} \StringTok{"0190..."}\FunctionTok{,}
  \DataTypeTok{"source\_\allowbreak{}pts"}\FunctionTok{:} \StringTok{"169387400"}\FunctionTok{,}
  \DataTypeTok{"source\_\allowbreak{}time\_\allowbreak{}base"}\FunctionTok{:} \FunctionTok{\{}\DataTypeTok{"num"}\FunctionTok{:} \DecValTok{1}\FunctionTok{,} \DataTypeTok{"den"}\FunctionTok{:} \DecValTok{60000}\FunctionTok{\},}
  \DataTypeTok{"capture\_\allowbreak{}time\_\allowbreak{}us"}\FunctionTok{:} \StringTok{"2823290234"}\FunctionTok{,}
  \DataTypeTok{"capture\_\allowbreak{}frame\_\allowbreak{}index"}\FunctionTok{:} \StringTok{"169228"}\FunctionTok{,}
  \DataTypeTok{"resolved\_\allowbreak{}player\_\allowbreak{}media\_\allowbreak{}time\_\allowbreak{}us"}\FunctionTok{:} \StringTok{"90166900"}\FunctionTok{,}
  \DataTypeTok{"mapping\_\allowbreak{}version"}\FunctionTok{:} \DecValTok{3}\FunctionTok{,}
  \DataTypeTok{"snap\_\allowbreak{}distance\_\allowbreak{}us"}\FunctionTok{:} \StringTok{"334"}\FunctionTok{,}
  \DataTypeTok{"timing\_\allowbreak{}precision"}\FunctionTok{:} \StringTok{"frame\_\allowbreak{}exact"}
\FunctionTok{\}}
\end{Highlighting}
\end{Shaded}

這個 server resolve 結果才可寫入 KeyPoint。無法解析時必須拒絕
command，不得退化成 client clock。

\subsection{8.6 Frame Step}\label{86-frame-step}

編輯 key point 時，前端送 key point ID 與方向：

\begin{Shaded}
\begin{Highlighting}[]
\FunctionTok{\{}\DataTypeTok{"key\_\allowbreak{}point\_\allowbreak{}id"}\FunctionTok{:}\StringTok{"..."}\FunctionTok{,}\DataTypeTok{"direction"}\FunctionTok{:}\StringTok{"next"}\FunctionTok{,}\DataTypeTok{"step\_\allowbreak{}count"}\FunctionTok{:}\DecValTok{1}\FunctionTok{\}}
\end{Highlighting}
\end{Shaded}

後端依 sample index 回傳前／後一個 frame 的 authoritative cursor。播放器
seek 到對應 media time；marker 更新仍走 revision command。

\begin{center}\rule{0.5\linewidth}{0.5pt}\end{center}

\section{9. Annotation 狀態機與完整 User
Flow}\label{9-annotation-ux72c0ux614bux6a5fux8207ux5b8cux6574-user-flow}

\subsection{9.1 狀態}\label{91-ux72c0ux614b}

\begin{Shaded}
\begin{Highlighting}[]
\NormalTok{OPEN}
\NormalTok{  ├─ \textless{} CLOSE\_\allowbreak{}RALLY(target, resolved/left) → READY}
\NormalTok{  ├─ \textgreater{} CLOSE\_\allowbreak{}RALLY(target, resolved/right) → READY}
\NormalTok{  └─ ? CLOSE\_\allowbreak{}RALLY(target, unknown/null) → READY}
\NormalTok{READY}
\NormalTok{  ├─ reopen → OPEN（清除terminal與outcome，之後重新CLOSE\_\allowbreak{}RALLY）}
\NormalTok{  └─ Enter → SUBMITTED}
\NormalTok{SUBMITTED}
\NormalTok{  └─ correction → 新 draft / 新 submission}
\NormalTok{VOIDED}
\end{Highlighting}
\end{Shaded}

處理狀態另存：\texttt{IDLE\ /\ CLIP\_\allowbreak{}QUEUED\ /\ CLIPPING\ /\ AI\_\allowbreak{}QUEUED\ /\ AI\_\allowbreak{}PROCESSING\ /\ INGESTING\ /\ COMPLETED\ /\ FAILED\ /\ SUPERSEDED}。

\subsection{9.2 完整操作}\label{92-ux5b8cux6574ux64cdux4f5c}

\subsubsection{Z - 發球}\label{z---ux767cux7403}

\begin{enumerate}
\def\labelenumi{\arabic{enumi}.}
\tightlist
\item
  播放器 cursor 必須 \texttt{ready} 且不在 gap。
\item
  Client送單一且可重試的 \texttt{CREATE\_\allowbreak{}SERVICE\_\allowbreak{}KEY\_\allowbreak{}POINT}
  command；\texttt{marker\_\allowbreak{}kind=service}由command語意固定。
\item
  Server解析 authoritative frame，建立 RallyDraft與 sequence 0 key
  point。
\item
  Timeline開始顯示灰色 open mask。
\item
  若已有 active draft，拒絕，不自動關閉前一 rally。
\end{enumerate}

\subsubsection{Space - 擊球}\label{space---ux64caux7403}

\begin{enumerate}
\def\labelenumi{\arabic{enumi}.}
\tightlist
\item
  僅 \texttt{OPEN} 可新增。
\item
  每次送一個 \texttt{contact} marker。
\item
  Server依 canonical time排序並產生 revision。
\item
  多人幾乎同時打同一球時保留兩點並標
  \texttt{possible\_\allowbreak{}duplicate}，不靜默刪除。
\end{enumerate}

\subsubsection{\texorpdfstring{\texttt{\textless{}}、\texttt{\textgreater{}}、\texttt{?}
- 關閉 Rally 並記錄 Outcome}{\textless、\textgreater、? - 關閉 Rally 並記錄 Outcome}}

\begin{enumerate}
\def\labelenumi{\arabic{enumi}.}
\tightlist
\item
  UI顯示三個按鈕：\texttt{\textless{}\ 左側得分}、\texttt{\textgreater{}\ 右側得分}、\texttt{?\ 未知}。硬體鍵盤使用實際
  \texttt{\textless{}}、\texttt{\textgreater{}}、\texttt{?}；方向鍵保留給逐幀。
\item
  Client從目前server-confirmed snapshot取得最後一個有效key point ID，作為
  \texttt{target\_\allowbreak{}key\_\allowbreak{}point\_\allowbreak{}id}，並送單一
  \texttt{CLOSE\_\allowbreak{}RALLY} command。
\item
  \texttt{\textless{}}傳\texttt{resolved/left}、\texttt{\textgreater{}}傳
  \texttt{resolved/right}、\texttt{?}傳\texttt{unknown/null}。Outcome只屬於rally，不是key point。
\item
  Command不含playback cursor、capture time、frame或score event；terminal時間／frame就是被選定key point既有的authoritative anchor。
\item
  Server在同一transaction驗證target於\texttt{base\_\allowbreak{}revision}仍是最後一個未刪除key point，將它標為terminal，保存outcome並把狀態改為\texttt{READY}。
\item
  若另一協作者已新增點，回\texttt{REVISION\_\allowbreak{}CONFLICT}／
  \texttt{CLOSE\_\allowbreak{}RALLY\_\allowbreak{}TARGET\_\allowbreak{}NOT\_\allowbreak{}LAST}；前端refetch後讓使用者再決定。
\item
  \texttt{left/right}指當下球場左右側，不是永久team identity；Server依immutable submission的side-assignment snapshot解析team ID。
\item
  \texttt{pending}與\texttt{unknown}不同：pending只允許open draft；unknown代表使用者已明確選\texttt{?}且可以提交。
\end{enumerate}

\subsubsection{Enter - 提交}\label{enter---ux63d0ux4ea4}

Server以單一 transaction：

\begin{enumerate}
\def\labelenumi{\arabic{enumi}.}
\tightlist
\item
  驗證至少一個 service、恰好一個 terminal、terminal為最後有效 key
  point。
\item
  驗證 score state是 \texttt{resolved} 或 \texttt{unknown}。
\item
  鎖定 draft revision並建立 immutable \texttt{RallySubmission}。
\item
  複製 immutable submission key points、side assignment、score
  resolution、clip policy與 content hash。
\item
  \texttt{resolved} 才建立 \texttt{PointAward}；\texttt{unknown}
  不改比分 ledger。
\item
  建立 ClipJob與Outbox event。
\item
  廣播 submission／processing狀態；UI mask轉綠。
\end{enumerate}

提交後不可原地修改。修正必須建立 correction draft並新建 submission，保留
supersedes關係。

\subsection{9.3 編輯與逐幀}\label{93-ux7de8ux8f2fux8207ux9010ux5e40}

\begin{itemize}
\tightlist
\item
  灰色 mask可新增、刪除、移動、重開rally；要改terminal或outcome時先
  \texttt{REOPEN\_\allowbreak{}RALLY}，再以新的\texttt{CLOSE\_\allowbreak{}RALLY} atomically關閉。
\item
  方向鍵在 key point edit mode 呼叫 server frame-step API；後端依 sample
  index取得前後 sample，再更新 marker。
\item
  綠色 mask唯讀；開啟修正時從 submission建立新 draft。
\item
  拖曳 marker時可用短暫 soft lock提示其他協作者，但 canonical
  concurrency仍以 revision/CAS為準。
\end{itemize}

\subsection{9.4 UI 必備區域}\label{94-ui-ux5fc5ux5099ux5340ux57df}

\begin{itemize}
\tightlist
\item
  頂部：場次、來源健康、LIVE狀態、與 live
  edge距離、WebSocket、協作者、來源 timecode。
\item
  中央：有聲 video、回到 LIVE、播放/暫停、逐幀、倍速。
\item
  多軌 timeline：完整 capture range、gap、playhead、key points、rally
  masks、score strip、processing badges。
\item
  底部固定 control deck：發球、擊球、左側得分、右側得分、未知、提交六個觸控動作，顯示目前快捷鍵（預設
  \texttt{Z}、\texttt{Space}、\texttt{\textless{}}、\texttt{\textgreater{}}、\texttt{?}、\texttt{Enter}）並依state啟停；不得顯示獨立結束控制。
\item
  Inspector：選取 key point時顯示 canonical
  time、frame、建立者、revision、duplicate/precision資訊。
\end{itemize}

\section{10. Annotation Realtime Schema
v2.0}\label{10-annotation-realtime-schema-v20}

\subsubsection{10.0 協議版本}\label{100-ux5354ux8b70ux7248ux672c}

Annotation Realtime Protocol 的正式版本為\texttt{2.0.0}。這是breaking
change：移除v1.1的獨立terminal與後續score兩階段command，改由一個
\texttt{CLOSE\_\allowbreak{}RALLY} atomically terminalize目前server-confirmed最後key
point並保存rally-level outcome。任何v1.x annotation message都不得被2.0
consumer靜默當成相同語意。

WebSocket用於高頻 annotation
command、ack、revision、presence與處理通知。GraphQL提供
snapshot/refetch，不承擔逐次按鍵。

\subsection{10.1 Client Command
Envelope}\label{101-client-command-envelope}

\begin{Shaded}
\begin{Highlighting}[]
\FunctionTok{\{}
  \DataTypeTok{"schema\_\allowbreak{}version"}\FunctionTok{:} \StringTok{"2.0.0"}\FunctionTok{,}
  \DataTypeTok{"command\_\allowbreak{}id"}\FunctionTok{:} \StringTok{"0190..."}\FunctionTok{,}
  \DataTypeTok{"room\_\allowbreak{}id"}\FunctionTok{:} \StringTok{"match:...:capture:..."}\FunctionTok{,}
  \DataTypeTok{"rally\_\allowbreak{}id"}\FunctionTok{:} \StringTok{"0190..."}\FunctionTok{,}
  \DataTypeTok{"base\_\allowbreak{}revision"}\FunctionTok{:} \StringTok{"27"}\FunctionTok{,}
  \DataTypeTok{"kind"}\FunctionTok{:} \StringTok{"CLOSE\_\allowbreak{}RALLY"}\FunctionTok{,}
  \DataTypeTok{"payload"}\FunctionTok{:} \FunctionTok{\{}
    \DataTypeTok{"target\_\allowbreak{}key\_\allowbreak{}point\_\allowbreak{}id"}\FunctionTok{:} \StringTok{"0190..."}\FunctionTok{,}
    \DataTypeTok{"score\_\allowbreak{}resolution"}\FunctionTok{:} \StringTok{"resolved"}\FunctionTok{,}
    \DataTypeTok{"scoring\_\allowbreak{}court\_\allowbreak{}side"}\FunctionTok{:} \StringTok{"left"}
  \FunctionTok{\}}
\FunctionTok{\}}
\end{Highlighting}
\end{Shaded}

\texttt{user\_\allowbreak{}id}與\texttt{device\_\allowbreak{}session\_\allowbreak{}id}不由每一個command自報。HTTP/WebSocket
upgrade先驗證登入身份並綁定device session；server寫入operation
log時使用connection
context。\texttt{CREATE\_\allowbreak{}SERVICE\_\allowbreak{}KEY\_\allowbreak{}POINT}的\texttt{rally\_\allowbreak{}id}由client在按Z當下產生UUID，讓「建立rally＋解析該frame＋建立service
marker」保持單一idempotent
command，不需要先呼叫另一個會造成時間漂移的API。

所有 revision、time、PTS、global frame在 JSON wire使用 decimal string。

\subsection{10.2 Command Kinds}\label{102-command-kinds}

\begin{longtable}[]{@{}>{\RaggedRight\arraybackslash}p{0.18\linewidth}>{\RaggedRight\arraybackslash}p{0.34\linewidth}>{\RaggedRight\arraybackslash}p{0.42\linewidth}@{}}
\toprule\noalign{}
kind & 主要 payload & 說明 \\
\midrule\noalign{}
\endhead
\bottomrule\noalign{}
\endlastfoot
\texttt{CREATE\_\allowbreak{}SERVICE\_\allowbreak{}KEY\_\allowbreak{}POINT} & \texttt{playback\_\allowbreak{}cursor} &
Z；建立 rally與service marker \\
\texttt{CREATE\_\allowbreak{}CONTACT\_\allowbreak{}KEY\_\allowbreak{}POINT} & \texttt{playback\_\allowbreak{}cursor} &
Space \\
\texttt{CLOSE\_\allowbreak{}RALLY} &
\texttt{target\_\allowbreak{}key\_\allowbreak{}point\_\allowbreak{}id} + strict outcome union &
\texttt{\textless{}}／\texttt{\textgreater{}}／\texttt{?}；atomically terminalize最後key point並保存rally-level outcome；不帶新時間／frame \\
\texttt{MOVE\_\allowbreak{}KEY\_\allowbreak{}POINT} & \texttt{key\_\allowbreak{}point\_\allowbreak{}id},
\texttt{playback\_\allowbreak{}cursor} & drag/frame
step後由server重新解析authoritative anchor \\
\texttt{DELETE\_\allowbreak{}KEY\_\allowbreak{}POINT} & \texttt{key\_\allowbreak{}point\_\allowbreak{}id} & draft-only \\
\texttt{REOPEN\_\allowbreak{}RALLY} & none & 清除terminal與outcome、回OPEN；之後以新的\texttt{CLOSE\_\allowbreak{}RALLY}重新關閉 \\
\texttt{VOID\_\allowbreak{}RALLY} & \texttt{reason} &
明確作廢draft，不等於刪除歷史 \\
\texttt{SUBMIT\_\allowbreak{}RALLY} & none & Enter，建立immutable submission \\
\end{longtable}

\subsection{10.3 Ack}\label{103-ack}

\begin{Shaded}
\begin{Highlighting}[]
\FunctionTok{\{}
  \DataTypeTok{"schema\_\allowbreak{}version"}\FunctionTok{:} \StringTok{"2.0.0"}\FunctionTok{,}
  \DataTypeTok{"type"}\FunctionTok{:} \StringTok{"command\_\allowbreak{}ack"}\FunctionTok{,}
  \DataTypeTok{"command\_\allowbreak{}id"}\FunctionTok{:} \StringTok{"0190..."}\FunctionTok{,}
  \DataTypeTok{"room\_\allowbreak{}id"}\FunctionTok{:} \StringTok{"match:...:capture:..."}\FunctionTok{,}
  \DataTypeTok{"rally\_\allowbreak{}id"}\FunctionTok{:} \StringTok{"0190..."}\FunctionTok{,}
  \DataTypeTok{"operation\_\allowbreak{}kind"}\FunctionTok{:} \StringTok{"CLOSE\_\allowbreak{}RALLY"}\FunctionTok{,}
  \DataTypeTok{"result\_\allowbreak{}revision"}\FunctionTok{:} \StringTok{"28"}\FunctionTok{,}
  \DataTypeTok{"server\_\allowbreak{}sequence"}\FunctionTok{:} \StringTok{"1042"}\FunctionTok{,}
  \DataTypeTok{"effects"}\FunctionTok{:} \FunctionTok{\{}
    \DataTypeTok{"terminal\_\allowbreak{}key\_\allowbreak{}point\_\allowbreak{}id"}\FunctionTok{:} \StringTok{"0190..."}\FunctionTok{,}
    \DataTypeTok{"annotation\_\allowbreak{}status"}\FunctionTok{:} \StringTok{"ready"}\FunctionTok{,}
    \DataTypeTok{"score\_\allowbreak{}resolution"}\FunctionTok{:} \StringTok{"resolved"}\FunctionTok{,}
    \DataTypeTok{"scoring\_\allowbreak{}court\_\allowbreak{}side"}\FunctionTok{:} \StringTok{"left"}
  \FunctionTok{\},}
  \DataTypeTok{"resolved\_\allowbreak{}anchor"}\FunctionTok{:} \KeywordTok{null}
\FunctionTok{\}}
\end{Highlighting}
\end{Shaded}

同一\texttt{command\_\allowbreak{}id}重送必須回相同結果。對
\texttt{CLOSE\_\allowbreak{}RALLY}，effects必須同時回傳terminal key point、
\texttt{ready}狀態與strict rally outcome，且
\texttt{resolved\_\allowbreak{}anchor}必須為null；schema不允許score time、score
frame或其他新event anchor。Server先寫DB transaction/outbox，再廣播。

\subsection{10.4 Conflict}\label{104-conflict}

\begin{Shaded}
\begin{Highlighting}[]
\FunctionTok{\{}
  \DataTypeTok{"type"}\FunctionTok{:} \StringTok{"command\_\allowbreak{}rejected"}\FunctionTok{,}
  \DataTypeTok{"command\_\allowbreak{}id"}\FunctionTok{:} \StringTok{"0190..."}\FunctionTok{,}
  \DataTypeTok{"code"}\FunctionTok{:} \StringTok{"CLOSE\_\allowbreak{}RALLY\_\allowbreak{}TARGET\_\allowbreak{}NOT\_\allowbreak{}LAST"}\FunctionTok{,}
  \DataTypeTok{"expected\_\allowbreak{}revision"}\FunctionTok{:} \StringTok{"27"}\FunctionTok{,}
  \DataTypeTok{"actual\_\allowbreak{}revision"}\FunctionTok{:} \StringTok{"28"}\FunctionTok{,}
  \DataTypeTok{"snapshot\_\allowbreak{}refetch\_\allowbreak{}required"}\FunctionTok{:} \KeywordTok{true}
\FunctionTok{\}}
\end{Highlighting}
\end{Shaded}

\texttt{CLOSE\_\allowbreak{}RALLY\_\allowbreak{}TARGET\_\allowbreak{}NOT\_\allowbreak{}LAST}、\texttt{REVISION\_\allowbreak{}CONFLICT}、\texttt{CURSOR\_\allowbreak{}STALE}、\texttt{CURSOR\_\allowbreak{}IN\_\allowbreak{}GAP}、\texttt{RALLY\_\allowbreak{}NOT\_\allowbreak{}OPEN}、\texttt{SCORE\_\allowbreak{}PENDING}等均用明確domain
code。Close target不再是最後key point時必須CAS失敗並要求snapshot
refetch，不可自動terminalize新點。

\section{11. GraphQL Schema
使用規格}\label{11-graphql-schema-ux4f7fux7528ux898fux683c}

\subsection{11.0 Code-first
產生規則}\label{110-code-first-ux7522ux751fux898fux5247}

\begin{itemize}
\tightlist
\item
  Pothos modules 是 source of truth；建議依 domain co-locate
  type/query/mutation。
\item
  \texttt{server/src/graphql/builder.ts} 只設定
  plugins/scalars/context，不定義 domain types。
\item
  \texttt{server/src/graphql/schema.ts} 只 import type modules並
  \texttt{builder.toSchema()}。
\item
  Prisma generator同時產生 Prisma Client與
  \texttt{prisma-pothos-types}。
\item
  \texttt{server/src/graphql/export-schema.ts} 使用
  \texttt{printSchema(lexicographicSortSchema(schema))}輸出 SDL。
\item
  CI先 \texttt{prisma\ generate}，再輸出SDL、跑 GraphQL
  Inspector，再執行 Web Codegen。
\item
  API內部錯誤使用穩定 error code；不要把 DB exception或 MinIO
  secret洩漏到 GraphQL message。
\end{itemize}

GraphQL 唯一原始碼位於 \texttt{server/src/graphql/**}，由 Pothos
code-first 建構。\texttt{packages/contracts/graphql/schema.graphql} 是
CI 產生並提交的 contract artifact，不可手動編輯。以下是前後端共同的核心
shape。

\subsection{11.1 BigInt}\label{111-bigint}

\begin{Shaded}
\begin{Highlighting}[]
\KeywordTok{scalar}\NormalTok{ BigInt }\CommentTok{\# JSON wire representation is a decimal string}
\KeywordTok{scalar}\NormalTok{ DateTime}
\KeywordTok{scalar}\NormalTok{ JSON}
\end{Highlighting}
\end{Shaded}

所有 PTS、microseconds、frame index、byte size使用 BigInt，不使用
GraphQL Int。

\subsection{11.2 Query}\label{112-query}

\begin{Shaded}
\begin{Highlighting}[]
\KeywordTok{type}\NormalTok{ Query \{}
\NormalTok{  me: User!}
\NormalTok{  matches(filter: MatchFilter): [Match!]!}
\NormalTok{  match(id: }\DataTypeTok{ID}\NormalTok{!): Match}
\NormalTok{  captureTimeline(captureSessionId: }\DataTypeTok{ID}\NormalTok{!): CaptureTimeline!}
\NormalTok{  rallies(matchId: }\DataTypeTok{ID}\NormalTok{!, filter: RallyFilter, page: PageInput): RallyConnection!}
\NormalTok{  rally(id: }\DataTypeTok{ID}\NormalTok{!): Rally}
\NormalTok{  analysisView(}\KeywordTok{input}\NormalTok{: AnalysisViewInput!): AnalysisView!}
\NormalTok{  featureAvailability(matchId: }\DataTypeTok{ID}\NormalTok{!): FeatureAvailability!}
\NormalTok{  savedAnalysisViews(matchId: }\DataTypeTok{ID}\NormalTok{!): [SavedAnalysisView!]!}
\NormalTok{\}}
\end{Highlighting}
\end{Shaded}

\subsection{11.3 Mutation}\label{113-mutation}

\begin{Shaded}
\begin{Highlighting}[]
\KeywordTok{type}\NormalTok{ Mutation \{}
\NormalTok{  createMatch(}\KeywordTok{input}\NormalTok{: CreateMatchInput!): Match!}
\NormalTok{  startCapture(}\KeywordTok{input}\NormalTok{: StartCaptureInput!): CaptureSession!}
\NormalTok{  stopCapture(captureSessionId: }\DataTypeTok{ID}\NormalTok{!): CaptureSession!}

\NormalTok{  applyAnnotationCommand(}\KeywordTok{input}\NormalTok{: AnnotationCommandInput!): AnnotationCommandResult!}
\NormalTok{  createCorrectionDraft(submissionId: }\DataTypeTok{ID}\NormalTok{!): Rally!}

\NormalTok{  assignTrackIdentity(}\KeywordTok{input}\NormalTok{: AssignTrackIdentityInput!): TrackIdentityAssignment!}
\NormalTok{  retryProcessing(}\KeywordTok{input}\NormalTok{: RetryProcessingInput!): ProcessingState!}
\NormalTok{  saveAnalysisView(}\KeywordTok{input}\NormalTok{: SaveAnalysisViewInput!): SavedAnalysisView!}
\NormalTok{\}}
\end{Highlighting}
\end{Shaded}

\texttt{applyAnnotationCommand} 是 WS 不可用時的 fallback，必須走相同
domain handler，不能有另一套語意。

\subsection{11.4 Rally Snapshot
範例}\label{114-rally-snapshot-ux7bc4ux4f8b}

\begin{Shaded}
\begin{Highlighting}[]
\FunctionTok{\{}
  \DataTypeTok{"id"}\FunctionTok{:} \StringTok{"018f..."}\FunctionTok{,}
  \DataTypeTok{"match\_\allowbreak{}id"}\FunctionTok{:} \StringTok{"..."}\FunctionTok{,}
  \DataTypeTok{"set\_\allowbreak{}id"}\FunctionTok{:} \StringTok{"..."}\FunctionTok{,}
  \DataTypeTok{"court\_\allowbreak{}side\_\allowbreak{}assignment"}\FunctionTok{:} \FunctionTok{\{}
    \DataTypeTok{"left\_\allowbreak{}team\_\allowbreak{}id"}\FunctionTok{:} \StringTok{"team{-}a"}\FunctionTok{,}
    \DataTypeTok{"right\_\allowbreak{}team\_\allowbreak{}id"}\FunctionTok{:} \StringTok{"team{-}b"}
  \FunctionTok{\},}
  \DataTypeTok{"annotation\_\allowbreak{}status"}\FunctionTok{:} \StringTok{"READY"}\FunctionTok{,}
  \DataTypeTok{"processing\_\allowbreak{}status"}\FunctionTok{:} \StringTok{"NOT\_\allowbreak{}REQUESTED"}\FunctionTok{,}
  \DataTypeTok{"revision"}\FunctionTok{:} \DecValTok{14}\FunctionTok{,}
  \DataTypeTok{"score\_\allowbreak{}resolution\_\allowbreak{}state"}\FunctionTok{:} \StringTok{"unknown"}\FunctionTok{,}
  \DataTypeTok{"scoring\_\allowbreak{}court\_\allowbreak{}side"}\FunctionTok{:} \KeywordTok{null}\FunctionTok{,}
  \DataTypeTok{"scoring\_\allowbreak{}team\_\allowbreak{}id"}\FunctionTok{:} \KeywordTok{null}\FunctionTok{,}
  \DataTypeTok{"key\_\allowbreak{}points"}\FunctionTok{:} \OtherTok{[}
    \FunctionTok{\{}
      \DataTypeTok{"id"}\FunctionTok{:} \StringTok{"..."}\FunctionTok{,}
      \DataTypeTok{"sequence\_\allowbreak{}index"}\FunctionTok{:} \DecValTok{0}\FunctionTok{,}
      \DataTypeTok{"marker\_\allowbreak{}kind"}\FunctionTok{:} \StringTok{"service"}\FunctionTok{,}
      \DataTypeTok{"is\_\allowbreak{}terminal"}\FunctionTok{:} \KeywordTok{false}\FunctionTok{,}
      \DataTypeTok{"capture\_\allowbreak{}time\_\allowbreak{}us"}\FunctionTok{:} \StringTok{"10000000"}\FunctionTok{,}
      \DataTypeTok{"capture\_\allowbreak{}frame\_\allowbreak{}index"}\FunctionTok{:} \StringTok{"599"}
    \FunctionTok{\}}
  \OtherTok{]}\FunctionTok{,}
  \DataTypeTok{"latest\_\allowbreak{}submission"}\FunctionTok{:} \KeywordTok{null}
\FunctionTok{\}}
\end{Highlighting}
\end{Shaded}

\subsection{11.5 Feature Availability}\label{115-feature-availability}

\begin{Shaded}
\begin{Highlighting}[]
\FunctionTok{\{}
  \DataTypeTok{"dimensions"}\FunctionTok{:} \FunctionTok{\{}
    \DataTypeTok{"player"}\FunctionTok{:} \KeywordTok{true}\FunctionTok{,}
    \DataTypeTok{"team"}\FunctionTok{:} \KeywordTok{true}\FunctionTok{,}
    \DataTypeTok{"action"}\FunctionTok{:} \KeywordTok{false}\FunctionTok{,}
    \DataTypeTok{"court\_\allowbreak{}zone"}\FunctionTok{:} \KeywordTok{true}\FunctionTok{,}
    \DataTypeTok{"direct\_\allowbreak{}outcome"}\FunctionTok{:} \KeywordTok{false}
  \FunctionTok{\},}
  \DataTypeTok{"metrics"}\FunctionTok{:} \OtherTok{[}
    \StringTok{"rally\_\allowbreak{}count"}\OtherTok{,}
    \StringTok{"event\_\allowbreak{}count"}\OtherTok{,}
    \StringTok{"rally\_\allowbreak{}win\_\allowbreak{}rate"}\OtherTok{,}
    \StringTok{"source\_\allowbreak{}heatmap"}\OtherTok{,}
    \StringTok{"target\_\allowbreak{}heatmap"}\OtherTok{,}
    \StringTok{"unresolved\_\allowbreak{}rate"}
  \OtherTok{]}\FunctionTok{,}
  \DataTypeTok{"action\_\allowbreak{}taxonomy"}\FunctionTok{:} \KeywordTok{null}
\FunctionTok{\}}
\end{Highlighting}
\end{Shaded}

前端不可假設 action／confidence 一定存在。

\begin{center}\rule{0.5\linewidth}{0.5pt}\end{center}

\section{12. PostgreSQL／Prisma
資料模型}\label{12-postgresqlprisma-ux8cc7ux6599ux6a21ux578b}

\subsection{12.1 核心分層}\label{121-ux6838ux5fc3ux5206ux5c64}

\begin{itemize}
\tightlist
\item
  Match／MatchSet／Team／Player／Roster／CourtSideAssignment。
\item
  CaptureSession／CaptureEpoch／DvrProgram／DvrSegment／MediaAsset。
\item
  RallyDraft（\texttt{Rally}）／mutable KeyPoint／AnnotationOperation。
\item
  immutable RallySubmission／RallySubmissionKeyPoint／PointAward。
\item
  ClipJob／ClipKeyPointMapping。
\item
  AiIntegration／AiJob／AiCallbackReceipt。
\item
  AnalysisRun／AnalysisTrack／ContactEvent／Actor／Candidate／BallPathSegment／AnalysisArtifact。
\item
  TrackIdentityAssignment／SavedAnalysisView／OutboxEvent。
\end{itemize}

\subsection{12.2
必要欄位與約束}\label{122-ux5fc5ux8981ux6b04ux4f4dux8207ux7d04ux675f}

\begin{itemize}
\tightlist
\item
  \texttt{CaptureEpoch} 保存 \texttt{source\_\allowbreak{}time\_\allowbreak{}base}、epoch起點 raw
  PTS、epoch對應 capture time；raw PTS允許負值；\texttt{DvrSegment}
  必須指向 epoch。
\item
  Draft KeyPoint保存
  \texttt{capture\_\allowbreak{}epoch\_\allowbreak{}id}、\texttt{source\_\allowbreak{}pts}、\texttt{capture\_\allowbreak{}time\_\allowbreak{}us}、\texttt{capture\_\allowbreak{}frame\_\allowbreak{}index}、\texttt{original\_\allowbreak{}playback\_\allowbreak{}cursor}與
  timing precision。
\item
  \texttt{RallySubmission} 快照必須保存
  \texttt{score\_\allowbreak{}resolution\_\allowbreak{}state}；unknown時 scoring side/team與score
  before/after允許 null。
\item
  \texttt{PointAward} 只存在於 resolved submission，且不可由 unknown
  submission建立；保存 score revision before/after，並對
  \texttt{(set\_\allowbreak{}id,\ score\_\allowbreak{}revision\_\allowbreak{}after)} 建唯一約束以序列化比分
  ledger。
\item
  ClipJob、AiJob、AnalysisRun都直接指向 immutable submission，不得只指
  mutable Rally。
\item
  Submitted資料的 immutable語意由 service layer、transaction與可選DB
  trigger共同保障；Prisma schema本身不足以表達所有check constraints。
\item
  \texttt{track\_\allowbreak{}id} 的primary key scope為
  \texttt{(analysis\_\allowbreak{}run\_\allowbreak{}id,\ track\_\allowbreak{}id)}；\texttt{AnalysisTrack.mean\_\allowbreak{}confidence}為optional，不得缺值時補0或1。
\item
  \texttt{AnalysisRun}將AI
  \texttt{producer.name}、\texttt{producer.build\_\allowbreak{}id}與optional
  \texttt{producer.sdk\_\allowbreak{}version}保存成可查詢欄位；完整原始結果仍保存為artifact。
\item
  ContactEvent primary relation使用 immutable submission key point
  ID；每個分析run每個key point恰好一筆；optional
  \texttt{resolved\_\allowbreak{}frame\_\allowbreak{}index}保存AI實際解析事件的clip frame。
\item
  每筆已解析Actor保存必填
  \texttt{observation\_\allowbreak{}frame\_\allowbreak{}index}，讓bbox、footprint與\texttt{court\_\allowbreak{}pos}的取樣frame不需由consumer猜測。
\item
  \texttt{RallySubmission.service\_\allowbreak{}key\_\allowbreak{}point\_\allowbreak{}id}與\texttt{terminal\_\allowbreak{}key\_\allowbreak{}point\_\allowbreak{}id}指向同一份immutable
  submission snapshot中的key point；DB以unique FK保證單一引用，service
  layer另驗證它們屬於該submission且terminal為序列最後一點。
\item
  64-bit時間、PTS、frame與byte size使用PostgreSQL BIGINT。
\end{itemize}

\subsection{12.3 JSONB
使用限制}\label{123-jsonb-ux4f7fux7528ux9650ux5236}

JSONB只適合：provider extensions、尚未固定的 action payload、原始 client
cursor audit、saved view
filters與非核心metadata。以下不得只放JSONB：ID關係、revision、score、time
anchor、artifact狀態、job狀態與可索引的分析主欄位。

\section{13. Clip
與時間碼前處理合約}\label{13-clip-ux8207ux6642ux9593ux78bcux524dux8655ux7406ux5408ux7d04}

\subsection{13.1 Submission → Clip Job}\label{131-submission--clip-job}

普通使用者不能傳 pre/post-roll。Server依 match clip profile 建：

\begin{Shaded}
\begin{Highlighting}[]
\FunctionTok{\{}
  \DataTypeTok{"clip\_\allowbreak{}job\_\allowbreak{}id"}\FunctionTok{:} \StringTok{"0190..."}\FunctionTok{,}
  \DataTypeTok{"rally\_\allowbreak{}submission\_\allowbreak{}id"}\FunctionTok{:} \StringTok{"018f..."}\FunctionTok{,}
  \DataTypeTok{"annotation\_\allowbreak{}revision"}\FunctionTok{:} \DecValTok{14}\FunctionTok{,}
  \DataTypeTok{"source\_\allowbreak{}range"}\FunctionTok{:} \FunctionTok{\{}
    \DataTypeTok{"logical\_\allowbreak{}start\_\allowbreak{}time\_\allowbreak{}us"}\FunctionTok{:} \StringTok{"2823290234"}\FunctionTok{,}
    \DataTypeTok{"logical\_\allowbreak{}end\_\allowbreak{}time\_\allowbreak{}us"}\FunctionTok{:} \StringTok{"2838123410"}\FunctionTok{,}
    \DataTypeTok{"requested\_\allowbreak{}start\_\allowbreak{}time\_\allowbreak{}us"}\FunctionTok{:} \StringTok{"2818290234"}\FunctionTok{,}
    \DataTypeTok{"requested\_\allowbreak{}end\_\allowbreak{}time\_\allowbreak{}us"}\FunctionTok{:} \StringTok{"2848123410"}
  \FunctionTok{\},}
  \DataTypeTok{"profile\_\allowbreak{}id"}\FunctionTok{:} \StringTok{"canonical{-}rally{-}v1"}
\FunctionTok{\}}
\end{Highlighting}
\end{Shaded}

\subsection{13.2 Canonical Clip}\label{132-canonical-clip}

\begin{itemize}
\tightlist
\item
  保持音訊。
\item
  固定旋轉與 pixel aspect ratio。
\item
  Canonical FPS/profile由部署設定；若輸入 VFR，mapping manifest
  必須從實際 output sample table建立。
\item
  Clip media timeline 從 0 開始。
\item
  每個 submission key point都要得到：

  \begin{itemize}
  \tightlist
  \item
    \texttt{clip\_\allowbreak{}pts}
  \item
    \texttt{clip\_\allowbreak{}time\_\allowbreak{}us}
  \item
    \texttt{clip\_\allowbreak{}frame\_\allowbreak{}index}
  \end{itemize}
\item
  不使用 \texttt{source\_\allowbreak{}time\ -\ start} 再乘 FPS 的近似方式產 frame
  index。
\end{itemize}

\subsection{13.3 Timing Manifest}\label{133-timing-manifest}

\begin{Shaded}
\begin{Highlighting}[]
\FunctionTok{\{}
  \DataTypeTok{"schema\_\allowbreak{}version"}\FunctionTok{:} \StringTok{"1.0.0"}\FunctionTok{,}
  \DataTypeTok{"clip\_\allowbreak{}id"}\FunctionTok{:} \StringTok{"..."}\FunctionTok{,}
  \DataTypeTok{"source"}\FunctionTok{:} \FunctionTok{\{}
    \DataTypeTok{"capture\_\allowbreak{}session\_\allowbreak{}id"}\FunctionTok{:} \StringTok{"..."}\FunctionTok{,}
    \DataTypeTok{"requested\_\allowbreak{}start\_\allowbreak{}time\_\allowbreak{}us"}\FunctionTok{:} \StringTok{"..."}\FunctionTok{,}
    \DataTypeTok{"actual\_\allowbreak{}start\_\allowbreak{}time\_\allowbreak{}us"}\FunctionTok{:} \StringTok{"..."}\FunctionTok{,}
    \DataTypeTok{"actual\_\allowbreak{}end\_\allowbreak{}time\_\allowbreak{}us"}\FunctionTok{:} \StringTok{"..."}
  \FunctionTok{\},}
  \DataTypeTok{"video"}\FunctionTok{:} \FunctionTok{\{}
    \DataTypeTok{"width"}\FunctionTok{:} \DecValTok{1920}\FunctionTok{,}
    \DataTypeTok{"height"}\FunctionTok{:} \DecValTok{1080}\FunctionTok{,}
    \DataTypeTok{"fps"}\FunctionTok{:} \FunctionTok{\{}\DataTypeTok{"num"}\FunctionTok{:} \DecValTok{60000}\FunctionTok{,} \DataTypeTok{"den"}\FunctionTok{:} \DecValTok{1001}\FunctionTok{\},}
    \DataTypeTok{"time\_\allowbreak{}base"}\FunctionTok{:} \FunctionTok{\{}\DataTypeTok{"num"}\FunctionTok{:} \DecValTok{1}\FunctionTok{,} \DataTypeTok{"den"}\FunctionTok{:} \DecValTok{60000}\FunctionTok{\},}
    \DataTypeTok{"total\_\allowbreak{}frames"}\FunctionTok{:} \StringTok{"1079"}\FunctionTok{,}
    \DataTypeTok{"duration\_\allowbreak{}us"}\FunctionTok{:} \StringTok{"18001333"}\FunctionTok{,}
    \DataTypeTok{"has\_\allowbreak{}audio"}\FunctionTok{:} \KeywordTok{true}
  \FunctionTok{\},}
  \DataTypeTok{"key\_\allowbreak{}points"}\FunctionTok{:} \OtherTok{[}
    \FunctionTok{\{}
      \DataTypeTok{"key\_\allowbreak{}point\_\allowbreak{}id"}\FunctionTok{:} \StringTok{"..."}\FunctionTok{,}
      \DataTypeTok{"source\_\allowbreak{}pts"}\FunctionTok{:} \StringTok{"..."}\FunctionTok{,}
      \DataTypeTok{"capture\_\allowbreak{}time\_\allowbreak{}us"}\FunctionTok{:} \StringTok{"..."}\FunctionTok{,}
      \DataTypeTok{"capture\_\allowbreak{}frame\_\allowbreak{}index"}\FunctionTok{:} \StringTok{"..."}\FunctionTok{,}
      \DataTypeTok{"clip\_\allowbreak{}pts"}\FunctionTok{:} \StringTok{"300300"}\FunctionTok{,}
      \DataTypeTok{"clip\_\allowbreak{}time\_\allowbreak{}us"}\FunctionTok{:} \StringTok{"5005000"}\FunctionTok{,}
      \DataTypeTok{"clip\_\allowbreak{}frame\_\allowbreak{}index"}\FunctionTok{:} \StringTok{"300"}
    \FunctionTok{\}}
  \OtherTok{]}
\FunctionTok{\}}
\end{Highlighting}
\end{Shaded}

\subsection{13.4 Clip Result}\label{134-clip-result}

\begin{itemize}
\tightlist
\item
  \texttt{canonical\_\allowbreak{}clip}：AI 分析來源。
\item
  \texttt{preview\_\allowbreak{}clip}：可低 bitrate，但必須有明確 mapping；若 frame
  count/FPS相同可直接對應。
\item
  \texttt{timing\_\allowbreak{}manifest}。
\item
  SHA-256、byte size、codec metadata。
\end{itemize}

只有 clip與所有 key point mapping驗證通過後，才可建立 AI Job。

\begin{center}\rule{0.5\linewidth}{0.5pt}\end{center}

\section{14. 外部 AI
串接：保留方法定義，不寫死模型細節}\label{14-ux5916ux90e8-ai-ux4e32ux63a5ux4fddux7559ux65b9ux6cd5ux5b9aux7fa9ux4e0dux5bebux6b7bux6a21ux578bux7d30ux7bc0}

AI團隊已經有球員追蹤、球追蹤、球場追蹤／投影、個體動作，以及可能存在的群體動作
work。本 repository不實作AI，只制定串接目標與結果語意。

\subsection{14.1
必要方法骨架}\label{141-ux5fc5ux8981ux65b9ux6cd5ux9aa8ux67b6}

\begin{enumerate}
\def\labelenumi{\arabic{enumi}.}
\tightlist
\item
  讀取 canonical rally clip與人工 key point anchors。
\item
  利用既有 work 取得可用的球員 tracks、球
  observation、frame座標、球場投影與 optional action evidence。
\item
  對每個輸入 key point產生恰好一筆 contact event，不跳過、不重編ID。
\item
  在該 key
  point附近建立球與球員的事件歸屬：單人、多位共同參與、多位候選、無法解析，或事件本身不適用球員（例如terminal落地／出界）。
\item
  以actor footprint的 \texttt{court\_\allowbreak{}pos}作代表點；若
  terminal沒有球員，AI可提供其既有方法產生的terminal representative
  court position，或明確回傳缺失。
\item
  由相鄰 contact events建立 A→B path segment。
\item
  回傳結構化Analysis JSON與逐幀FlatBuffers overlay。
\end{enumerate}

\subsection{14.2
不在本規格寫死的事項}\label{142-ux4e0dux5728ux672cux898fux683cux5bebux6b7bux7684ux4e8bux9805}

\begin{itemize}
\tightlist
\item
  模型架構、模組執行順序、時間窗大小。
\item
  球距離、bbox distance、IoU、pose、action的融合公式或門檻。
\item
  補點、插值、平滑與confidence threshold。
\item
  action label清單、模型原始輸出欄位、群體phase label。
\end{itemize}

\subsection{14.3
球歸屬與多球員重疊}\label{143-ux7403ux6b78ux5c6cux8207ux591aux7403ux54e1ux91cdux758a}

\begin{itemize}
\tightlist
\item
  \texttt{resolved\_\allowbreak{}single}：明確一位 actor。
\item
  \texttt{resolved\_\allowbreak{}multiple}：AI明確判定多人共同參與，例如雙人攔網；不是因為bbox重疊就自動成立。
\item
  \texttt{ambiguous}：有候選但無法決定；\texttt{actors={[}{]}}、\texttt{actor\_\allowbreak{}candidates{[}{]}}有值。
\item
  \texttt{unresolved}：無法提供有效actor或候選。
\item
  \texttt{no\_\allowbreak{}player}：事件語意本身不需要球員，例如terminal球落地／出界；與模型失敗不同。
\end{itemize}

高IoU或球位於重疊框只是一種evidence。actors使用陣列以保留多人事件，但AI團隊仍需自行判定何時是共同參與、何時只是ambiguous。

\subsection{14.4
Action與群體phase待確認}\label{144-actionux8207ux7fa4ux9ad4phaseux5f85ux78baux8a8d}

\begin{itemize}
\tightlist
\item
  \texttt{action}為optional extension；baseline result不要求。
\item
  label由provider定義，並帶optional taxonomy
  ID/version；中央與前端不能硬編碼serve/receive/set/spike等清單。
\item
  confidence也是optional；沒有就省略，不得假造。
\item
  群體動作／比賽phase目前沒有已確認產品用途，不列第一版required output。
\item
  AI團隊後續需回答：真實label清單、粒度、confidence語意，以及group
  phase可支援哪個教練決策。
\end{itemize}

\section{15. AI Provider API 與 Python
SDK}\label{15-ai-provider-api-ux8207-python-sdk}

\subsection{15.1 Provider Capabilities}\label{151-provider-capabilities}

\texttt{GET\ /v1/capabilities} 只聲明支援的contract版本與optional
extension，不遠端配置模型參數。

\begin{Shaded}
\begin{Highlighting}[]
\FunctionTok{\{}
  \DataTypeTok{"schema\_\allowbreak{}version"}\FunctionTok{:} \StringTok{"1.0.0"}\FunctionTok{,}
  \DataTypeTok{"provider\_\allowbreak{}name"}\FunctionTok{:} \StringTok{"team{-}ai"}\FunctionTok{,}
  \DataTypeTok{"provider\_\allowbreak{}build\_\allowbreak{}id"}\FunctionTok{:} \StringTok{"2026{-}08{-}07.1"}\FunctionTok{,}
  \DataTypeTok{"supported\_\allowbreak{}job\_\allowbreak{}schema\_\allowbreak{}versions"}\FunctionTok{:} \OtherTok{[}\StringTok{"1.1.0"}\OtherTok{]}\FunctionTok{,}
  \DataTypeTok{"supported\_\allowbreak{}result\_\allowbreak{}schema\_\allowbreak{}versions"}\FunctionTok{:} \OtherTok{[}\StringTok{"1.0.0"}\OtherTok{]}\FunctionTok{,}
  \DataTypeTok{"supported\_\allowbreak{}overlay\_\allowbreak{}formats"}\FunctionTok{:} \OtherTok{[}\StringTok{"flatbuffers\_\allowbreak{}v1"}\OtherTok{]}\FunctionTok{,}
  \DataTypeTok{"optional\_\allowbreak{}extensions"}\FunctionTok{:} \FunctionTok{\{}
    \DataTypeTok{"action"}\FunctionTok{:} \KeywordTok{true}\FunctionTok{,}
    \DataTypeTok{"group\_\allowbreak{}phase"}\FunctionTok{:} \KeywordTok{false}\FunctionTok{,}
    \DataTypeTok{"confidence"}\FunctionTok{:} \KeywordTok{false}
  \FunctionTok{\},}
  \DataTypeTok{"action\_\allowbreak{}taxonomies"}\FunctionTok{:} \OtherTok{[]}
\FunctionTok{\}}
\end{Highlighting}
\end{Shaded}

\subsection{15.2 AI Job Request}\label{152-ai-job-request}

AI Job Request 的正式版本為
\texttt{1.1.0}；\texttt{AnalysisResult}、callback envelope、provider
capabilities envelope 與 \texttt{JobAccepted} 仍各自維持
\texttt{1.0.0}。Job 1.1 新增必填的
\texttt{clip.download\_\allowbreak{}url\_\allowbreak{}expires\_\allowbreak{}at}，且明確分離 clip signed URL 與
callback bearer token。Capabilities 必須宣告
\texttt{supported\_\allowbreak{}job\_\allowbreak{}schema\_\allowbreak{}versions} 包含 \texttt{1.1.0} 才能接收。

\begin{Shaded}
\begin{Highlighting}[]
\NormalTok{POST /v1/jobs}
\NormalTok{Authorization: Bearer \textless{}provider{-}token\textgreater{}}
\NormalTok{Idempotency{-}Key: \textless{}ai\_\allowbreak{}job\_\allowbreak{}id\textgreater{}}
\NormalTok{Content{-}Type: application/json}
\end{Highlighting}
\end{Shaded}

\begin{Shaded}
\begin{Highlighting}[]
\FunctionTok{\{}
  \DataTypeTok{"schema\_\allowbreak{}version"}\FunctionTok{:} \StringTok{"1.1.0"}\FunctionTok{,}
  \DataTypeTok{"ai\_\allowbreak{}job\_\allowbreak{}id"}\FunctionTok{:} \StringTok{"018f6d86{-}2d8c{-}7f10{-}9c1b{-}0f0c32aa6001"}\FunctionTok{,}
  \DataTypeTok{"rally\_\allowbreak{}submission\_\allowbreak{}id"}\FunctionTok{:} \StringTok{"018f6d86{-}2d8c{-}7f10{-}9c1b{-}0f0c32aa3101"}\FunctionTok{,}
  \DataTypeTok{"rally\_\allowbreak{}id"}\FunctionTok{:} \StringTok{"018f6d86{-}2d8c{-}7f10{-}9c1b{-}0f0c32aa3001"}\FunctionTok{,}
  \DataTypeTok{"match\_\allowbreak{}id"}\FunctionTok{:} \StringTok{"018f6d86{-}2d8c{-}7f10{-}9c1b{-}0f0c32aa1001"}\FunctionTok{,}
  \DataTypeTok{"annotation\_\allowbreak{}revision"}\FunctionTok{:} \StringTok{"14"}\FunctionTok{,}
  \DataTypeTok{"clip"}\FunctionTok{:} \FunctionTok{\{}
    \DataTypeTok{"clip\_\allowbreak{}asset\_\allowbreak{}id"}\FunctionTok{:} \StringTok{"018f6d86{-}2d8c{-}7f10{-}9c1b{-}0f0c32aa5001"}\FunctionTok{,}
    \DataTypeTok{"download\_\allowbreak{}url"}\FunctionTok{:} \StringTok{"https://media.example/signed/canonical{-}clip.mp4?..."}\FunctionTok{,}
    \DataTypeTok{"download\_\allowbreak{}url\_\allowbreak{}expires\_\allowbreak{}at"}\FunctionTok{:} \StringTok{"2026{-}08{-}07T06:00:00Z"}\FunctionTok{,}
    \DataTypeTok{"sha256"}\FunctionTok{:} \StringTok{"9a31d43d2c2aeb1a1a98beec4ff8bbd649f68f16eb7345a51fd61a3b302543aa"}\FunctionTok{,}
    \DataTypeTok{"byte\_\allowbreak{}length"}\FunctionTok{:} \StringTok{"83124491"}\FunctionTok{,}
    \DataTypeTok{"content\_\allowbreak{}type"}\FunctionTok{:} \StringTok{"video/mp4"}\FunctionTok{,}
    \DataTypeTok{"video"}\FunctionTok{:} \FunctionTok{\{}
      \DataTypeTok{"width"}\FunctionTok{:} \DecValTok{1920}\FunctionTok{,}
      \DataTypeTok{"height"}\FunctionTok{:} \DecValTok{1080}\FunctionTok{,}
      \DataTypeTok{"fps"}\FunctionTok{:} \FunctionTok{\{}\DataTypeTok{"num"}\FunctionTok{:} \DecValTok{60000}\FunctionTok{,} \DataTypeTok{"den"}\FunctionTok{:} \DecValTok{1001}\FunctionTok{\},}
      \DataTypeTok{"time\_\allowbreak{}base"}\FunctionTok{:} \FunctionTok{\{}\DataTypeTok{"num"}\FunctionTok{:} \DecValTok{1}\FunctionTok{,} \DataTypeTok{"den"}\FunctionTok{:} \DecValTok{60000}\FunctionTok{\},}
      \DataTypeTok{"total\_\allowbreak{}frames"}\FunctionTok{:} \StringTok{"1079"}\FunctionTok{,}
      \DataTypeTok{"duration\_\allowbreak{}us"}\FunctionTok{:} \StringTok{"18001333"}\FunctionTok{,}
      \DataTypeTok{"has\_\allowbreak{}audio"}\FunctionTok{:} \KeywordTok{true}
    \FunctionTok{\}}
  \FunctionTok{\},}
  \DataTypeTok{"key\_\allowbreak{}points"}\FunctionTok{:} \OtherTok{[}
    \FunctionTok{\{}
      \DataTypeTok{"key\_\allowbreak{}point\_\allowbreak{}id"}\FunctionTok{:} \StringTok{"...4000"}\FunctionTok{,}
      \DataTypeTok{"sequence\_\allowbreak{}index"}\FunctionTok{:} \DecValTok{0}\FunctionTok{,}
      \DataTypeTok{"marker\_\allowbreak{}kind"}\FunctionTok{:} \StringTok{"service"}\FunctionTok{,}
      \DataTypeTok{"is\_\allowbreak{}terminal"}\FunctionTok{:} \KeywordTok{false}\FunctionTok{,}
      \DataTypeTok{"clip\_\allowbreak{}pts"}\FunctionTok{:} \StringTok{"300300"}\FunctionTok{,}
      \DataTypeTok{"clip\_\allowbreak{}time\_\allowbreak{}us"}\FunctionTok{:} \StringTok{"5005000"}\FunctionTok{,}
      \DataTypeTok{"clip\_\allowbreak{}frame\_\allowbreak{}index"}\FunctionTok{:} \StringTok{"300"}
    \FunctionTok{\}}\OtherTok{,}
    \FunctionTok{\{}
      \DataTypeTok{"key\_\allowbreak{}point\_\allowbreak{}id"}\FunctionTok{:} \StringTok{"...4010"}\FunctionTok{,}
      \DataTypeTok{"sequence\_\allowbreak{}index"}\FunctionTok{:} \DecValTok{1}\FunctionTok{,}
      \DataTypeTok{"marker\_\allowbreak{}kind"}\FunctionTok{:} \StringTok{"contact"}\FunctionTok{,}
      \DataTypeTok{"is\_\allowbreak{}terminal"}\FunctionTok{:} \KeywordTok{true}\FunctionTok{,}
      \DataTypeTok{"clip\_\allowbreak{}pts"}\FunctionTok{:} \StringTok{"480480"}\FunctionTok{,}
      \DataTypeTok{"clip\_\allowbreak{}time\_\allowbreak{}us"}\FunctionTok{:} \StringTok{"8008000"}\FunctionTok{,}
      \DataTypeTok{"clip\_\allowbreak{}frame\_\allowbreak{}index"}\FunctionTok{:} \StringTok{"480"}
    \FunctionTok{\}}
  \OtherTok{]}\FunctionTok{,}
  \DataTypeTok{"outcome"}\FunctionTok{:} \FunctionTok{\{}
    \DataTypeTok{"score\_\allowbreak{}resolution"}\FunctionTok{:} \StringTok{"resolved"}\FunctionTok{,}
    \DataTypeTok{"scoring\_\allowbreak{}court\_\allowbreak{}side"}\FunctionTok{:} \StringTok{"left"}
  \FunctionTok{\},}
  \DataTypeTok{"callback"}\FunctionTok{:} \FunctionTok{\{}
    \DataTypeTok{"url"}\FunctionTok{:} \StringTok{"https://central.example/api/v1/ai/callback/018f..."}\FunctionTok{,}
    \DataTypeTok{"token"}\FunctionTok{:} \StringTok{"job{-}scoped{-}opaque{-}token"}\FunctionTok{,}
    \DataTypeTok{"expires\_\allowbreak{}at"}\FunctionTok{:} \StringTok{"2026{-}08{-}07T06:00:00Z"}
  \FunctionTok{\}}
\FunctionTok{\}}
\end{Highlighting}
\end{Shaded}

Unknown outcome：

\begin{Shaded}
\begin{Highlighting}[]
\FunctionTok{\{}\DataTypeTok{"score\_\allowbreak{}resolution"}\FunctionTok{:}\StringTok{"unknown"}\FunctionTok{,}\DataTypeTok{"scoring\_\allowbreak{}court\_\allowbreak{}side"}\FunctionTok{:}\KeywordTok{null}\FunctionTok{\}}
\end{Highlighting}
\end{Shaded}

Job不包含court definition、team IDs、model
requirements、threshold、MinIO upload
URI或中央DB設定。球場座標規格是固定contract，不需要每個job重複傳。

\texttt{clip.download\_\allowbreak{}url}必須是獨立的短期signed
URL，SDK下載時不得夾帶\texttt{callback.token}。\texttt{callback.token}只用於對中央callback
URL送出Bearer token；不得送往MinIO/S3、不得出現在browser或log。

\subsubsection{AI欄位 ownership}\label{aiux6b04ux4f4d-ownership}

\begin{longtable}[]{@{}>{\RaggedRight\arraybackslash}p{0.18\linewidth}>{\RaggedRight\arraybackslash}p{0.34\linewidth}>{\RaggedRight\arraybackslash}p{0.42\linewidth}@{}}
\toprule\noalign{}
欄位 & 來源 & AI端行為 \\
\midrule\noalign{}
\endhead
\bottomrule\noalign{}
\endlastfoot
\texttt{ai\_\allowbreak{}job\_\allowbreak{}id} & CENTRAL + PASSTHROUGH &
result/callback原樣回傳 \\
\texttt{rally\_\allowbreak{}submission\_\allowbreak{}id} & CENTRAL + PASSTHROUGH &
原樣回傳；AI結果的版本錨點 \\
\texttt{rally\_\allowbreak{}id}, \texttt{match\_\allowbreak{}id} & CENTRAL + PASSTHROUGH &
原樣回傳 \\
\texttt{annotation\_\allowbreak{}revision} & SUBMISSION SNAPSHOT + PASSTHROUGH &
原樣回傳 \\
\texttt{clip\_\allowbreak{}asset\_\allowbreak{}id} & PREPROCESS + PASSTHROUGH & 原樣回傳 \\
clip URL/checksum/video metadata & PREPROCESS &
下載並驗證；不回寫修改 \\
\texttt{key\_\allowbreak{}point\_\allowbreak{}id} & SUBMISSION SNAPSHOT + PASSTHROUGH &
每個恰好出現一次 \\
marker/sequence/terminal/clip timing & PREPROCESS + PASSTHROUGH &
作anchor，不改寫 \\
\texttt{outcome} & SUBMISSION SNAPSHOT &
可使用，不需要在result重複回傳 \\
callback URL/token & CENTRAL & token只授權中央callback
POST；clip使用獨立signed URL；兩者皆不放result、不下發browser \\
track/event/action/position & AI\_\allowbreak{}GENERATED & 由AI輸出 \\
\end{longtable}

\subsection{15.3 Provider Accepted}\label{153-provider-accepted}

\begin{Shaded}
\begin{Highlighting}[]
\FunctionTok{\{}
  \DataTypeTok{"schema\_\allowbreak{}version"}\FunctionTok{:} \StringTok{"1.0.0"}\FunctionTok{,}
  \DataTypeTok{"ai\_\allowbreak{}job\_\allowbreak{}id"}\FunctionTok{:} \StringTok{"..."}\FunctionTok{,}
  \DataTypeTok{"provider\_\allowbreak{}job\_\allowbreak{}id"}\FunctionTok{:} \StringTok{"provider{-}123"}\FunctionTok{,}
  \DataTypeTok{"state"}\FunctionTok{:} \StringTok{"accepted"}\FunctionTok{,}
  \DataTypeTok{"accepted\_\allowbreak{}at"}\FunctionTok{:} \StringTok{"2026{-}08{-}07T05:30:00Z"}
\FunctionTok{\}}
\end{Highlighting}
\end{Shaded}

相同Idempotency-Key重送應回同一provider job。

\subsection{15.4 Python SDK}\label{154-python-sdk}

GitHub安裝：

\begin{Shaded}
\begin{Highlighting}[]
\ExtensionTok{pip}\NormalTok{ install }\StringTok{"volleyball{-}monitoring{-}ai{-}sdk @ git+https://github.com/\textless{}owner\textgreater{}/volleyball{-}monitoring{-}ai.git@v0.1.0\#subdirectory=sdk"}
\end{Highlighting}
\end{Shaded}

SDK必須提供Pydantic job/result/capabilities/accepted models、JSON Schema
validation、signed clip下載與checksum、callback client、FlatBuffers
helper、passthrough/invariant validator、FastAPI provider adapter
optional extra與fake
fixtures。SDK不依賴torch/CUDA/OpenCV，也不實作AI模型。

\section{16. AI Analysis Result
Schema}\label{16-ai-analysis-result-schema}

\subsection{16.1 結果範例}\label{161-ux7d50ux679cux7bc4ux4f8b}

\begin{Shaded}
\begin{Highlighting}[]
\FunctionTok{\{}
  \DataTypeTok{"schema\_\allowbreak{}version"}\FunctionTok{:} \StringTok{"1.0.0"}\FunctionTok{,}
  \DataTypeTok{"analysis\_\allowbreak{}id"}\FunctionTok{:} \StringTok{"analysis{-}0190"}\FunctionTok{,}
  \DataTypeTok{"analysis\_\allowbreak{}version"}\FunctionTok{:} \StringTok{"provider{-}build{-}2026{-}08{-}07.1"}\FunctionTok{,}
  \DataTypeTok{"ai\_\allowbreak{}job\_\allowbreak{}id"}\FunctionTok{:} \StringTok{"018f...6001"}\FunctionTok{,}
  \DataTypeTok{"rally\_\allowbreak{}submission\_\allowbreak{}id"}\FunctionTok{:} \StringTok{"018f...3101"}\FunctionTok{,}
  \DataTypeTok{"rally\_\allowbreak{}id"}\FunctionTok{:} \StringTok{"018f...3001"}\FunctionTok{,}
  \DataTypeTok{"match\_\allowbreak{}id"}\FunctionTok{:} \StringTok{"018f...1001"}\FunctionTok{,}
  \DataTypeTok{"annotation\_\allowbreak{}revision"}\FunctionTok{:} \StringTok{"14"}\FunctionTok{,}
  \DataTypeTok{"clip\_\allowbreak{}asset\_\allowbreak{}id"}\FunctionTok{:} \StringTok{"018f...5001"}\FunctionTok{,}
  \DataTypeTok{"input\_\allowbreak{}clip\_\allowbreak{}sha256"}\FunctionTok{:} \StringTok{"9a31d43d2c2aeb1a1a98beec4ff8bbd649f68f16eb7345a51fd61a3b302543aa"}\FunctionTok{,}
  \DataTypeTok{"producer"}\FunctionTok{:} \FunctionTok{\{}\DataTypeTok{"name"}\FunctionTok{:}\StringTok{"team{-}ai"}\FunctionTok{,}\DataTypeTok{"build\_\allowbreak{}id"}\FunctionTok{:}\StringTok{"2026{-}08{-}07.1"}\FunctionTok{,}\DataTypeTok{"sdk\_\allowbreak{}version"}\FunctionTok{:}\StringTok{"0.1.0"}\FunctionTok{\},}
  \DataTypeTok{"tracks"}\FunctionTok{:} \OtherTok{[}
    \FunctionTok{\{}\DataTypeTok{"track\_\allowbreak{}id"}\FunctionTok{:}\DecValTok{8}\FunctionTok{,}\DataTypeTok{"court\_\allowbreak{}side"}\FunctionTok{:}\StringTok{"left"}\FunctionTok{,}\DataTypeTok{"first\_\allowbreak{}frame\_\allowbreak{}index"}\FunctionTok{:}\StringTok{"0"}\FunctionTok{,}\DataTypeTok{"last\_\allowbreak{}frame\_\allowbreak{}index"}\FunctionTok{:}\StringTok{"1078"}\FunctionTok{\}}
  \OtherTok{]}\FunctionTok{,}
  \DataTypeTok{"contact\_\allowbreak{}events"}\FunctionTok{:} \OtherTok{[}
    \FunctionTok{\{}
      \DataTypeTok{"key\_\allowbreak{}point\_\allowbreak{}id"}\FunctionTok{:}\StringTok{"...4000"}\FunctionTok{,}
      \DataTypeTok{"sequence\_\allowbreak{}index"}\FunctionTok{:}\DecValTok{0}\FunctionTok{,}
      \DataTypeTok{"marker\_\allowbreak{}kind"}\FunctionTok{:}\StringTok{"service"}\FunctionTok{,}
      \DataTypeTok{"is\_\allowbreak{}terminal"}\FunctionTok{:}\KeywordTok{false}\FunctionTok{,}
      \DataTypeTok{"anchor\_\allowbreak{}frame\_\allowbreak{}index"}\FunctionTok{:}\StringTok{"300"}\FunctionTok{,}
      \DataTypeTok{"resolved\_\allowbreak{}frame\_\allowbreak{}index"}\FunctionTok{:}\StringTok{"301"}\FunctionTok{,}
      \DataTypeTok{"association\_\allowbreak{}state"}\FunctionTok{:}\StringTok{"resolved\_\allowbreak{}single"}\FunctionTok{,}
      \DataTypeTok{"actors"}\FunctionTok{:}\OtherTok{[}\FunctionTok{\{}
        \DataTypeTok{"track\_\allowbreak{}id"}\FunctionTok{:}\DecValTok{8}\FunctionTok{,}
        \DataTypeTok{"observation\_\allowbreak{}frame\_\allowbreak{}index"}\FunctionTok{:}\StringTok{"301"}\FunctionTok{,}
        \DataTypeTok{"frame\_\allowbreak{}bbox"}\FunctionTok{:\{}\DataTypeTok{"x1"}\FunctionTok{:}\DecValTok{0}\ErrorTok{.}\DecValTok{12}\FunctionTok{,}\DataTypeTok{"y1"}\FunctionTok{:}\DecValTok{0}\ErrorTok{.}\DecValTok{30}\FunctionTok{,}\DataTypeTok{"x2"}\FunctionTok{:}\DecValTok{0}\ErrorTok{.}\DecValTok{21}\FunctionTok{,}\DataTypeTok{"y2"}\FunctionTok{:}\DecValTok{0}\ErrorTok{.}\DecValTok{81}\FunctionTok{\},}
        \DataTypeTok{"frame\_\allowbreak{}foot\_\allowbreak{}pos"}\FunctionTok{:\{}\DataTypeTok{"x"}\FunctionTok{:}\DecValTok{0}\ErrorTok{.}\DecValTok{165}\FunctionTok{,}\DataTypeTok{"y"}\FunctionTok{:}\DecValTok{0}\ErrorTok{.}\DecValTok{81}\FunctionTok{\},}
        \DataTypeTok{"court\_\allowbreak{}pos"}\FunctionTok{:\{}\DataTypeTok{"x"}\FunctionTok{:}\DecValTok{{-}0}\ErrorTok{.06}\FunctionTok{,}\DataTypeTok{"y"}\FunctionTok{:}\DecValTok{0}\ErrorTok{.}\DecValTok{24}\FunctionTok{\}}
      \FunctionTok{\}}\OtherTok{]}\FunctionTok{,}
      \DataTypeTok{"actor\_\allowbreak{}candidates"}\FunctionTok{:}\OtherTok{[]}\FunctionTok{,}
      \DataTypeTok{"ball"}\FunctionTok{:\{}\DataTypeTok{"state"}\FunctionTok{:}\StringTok{"observed"}\FunctionTok{,}\DataTypeTok{"sample\_\allowbreak{}frame\_\allowbreak{}index"}\FunctionTok{:}\StringTok{"301"}\FunctionTok{,}\DataTypeTok{"frame\_\allowbreak{}pos"}\FunctionTok{:\{}\DataTypeTok{"x"}\FunctionTok{:}\DecValTok{0}\ErrorTok{.}\DecValTok{182}\FunctionTok{,}\DataTypeTok{"y"}\FunctionTok{:}\DecValTok{0}\ErrorTok{.}\DecValTok{614}\FunctionTok{\}\},}
      \DataTypeTok{"representative\_\allowbreak{}court\_\allowbreak{}positions"}\FunctionTok{:}\OtherTok{[}
        \FunctionTok{\{}\DataTypeTok{"track\_\allowbreak{}id"}\FunctionTok{:}\DecValTok{8}\FunctionTok{,}\DataTypeTok{"basis"}\FunctionTok{:}\StringTok{"player\_\allowbreak{}footprint\_\allowbreak{}proxy"}\FunctionTok{,}\DataTypeTok{"court\_\allowbreak{}pos"}\FunctionTok{:\{}\DataTypeTok{"x"}\FunctionTok{:}\DecValTok{{-}0}\ErrorTok{.06}\FunctionTok{,}\DataTypeTok{"y"}\FunctionTok{:}\DecValTok{0}\ErrorTok{.}\DecValTok{24}\FunctionTok{\}\}}
      \OtherTok{]}\FunctionTok{,}
      \DataTypeTok{"quality\_\allowbreak{}flags"}\FunctionTok{:}\OtherTok{[]}
    \FunctionTok{\}}
  \OtherTok{]}\FunctionTok{,}
  \DataTypeTok{"path\_\allowbreak{}segments"}\FunctionTok{:} \OtherTok{[]}\FunctionTok{,}
  \DataTypeTok{"summary"}\FunctionTok{:\{}\DataTypeTok{"track\_\allowbreak{}count"}\FunctionTok{:}\DecValTok{1}\FunctionTok{,}\DataTypeTok{"contact\_\allowbreak{}event\_\allowbreak{}count"}\FunctionTok{:}\DecValTok{1}\FunctionTok{,}\DataTypeTok{"path\_\allowbreak{}segment\_\allowbreak{}count"}\FunctionTok{:}\DecValTok{0}\FunctionTok{,}\DataTypeTok{"unresolved\_\allowbreak{}event\_\allowbreak{}count"}\FunctionTok{:}\DecValTok{0}\FunctionTok{\},}
  \DataTypeTok{"extensions"}\FunctionTok{:\{\}}
\FunctionTok{\}}
\end{Highlighting}
\end{Shaded}

\subsection{16.2 Passthrough與1:1
invariant}\label{162-passthroughux820711-invariant}

\begin{itemize}
\tightlist
\item
  \texttt{ai\_\allowbreak{}job\_\allowbreak{}id/rally\_\allowbreak{}submission\_\allowbreak{}id/rally\_\allowbreak{}id/match\_\allowbreak{}id/annotation\_\allowbreak{}revision/clip\_\allowbreak{}asset\_\allowbreak{}id}原樣回傳。
\item
  每個input key point恰好一個contact event。
\item
  \texttt{key\_\allowbreak{}point\_\allowbreak{}id/sequence\_\allowbreak{}index/marker\_\allowbreak{}kind/is\_\allowbreak{}terminal}與job完全一致。
\item
  N個events通常有N-1個segments；只有單一service且terminal的回合可為0段。
\item
  無法判斷actor仍回event，不得刪掉。
\end{itemize}

\subsection{16.3 Actor observation
frame}\label{163-actor-observation-frame}

每個已解析 actor 必須帶 \texttt{observation\_\allowbreak{}frame\_\allowbreak{}index}，表示
\texttt{frame\_\allowbreak{}bbox}、\texttt{frame\_\allowbreak{}foot\_\allowbreak{}pos}與
\texttt{court\_\allowbreak{}pos}取樣的clip
frame。AI可在人工anchor附近選擇更可靠的frame，但不得讓consumer猜測這些座標來自哪一幀。

AI wire result不需要生成跨系統 \texttt{segment\_\allowbreak{}id}；每段路徑以
\texttt{sequence\_\allowbreak{}index\ +\ start\_\allowbreak{}key\_\allowbreak{}point\_\allowbreak{}id\ +\ end\_\allowbreak{}key\_\allowbreak{}point\_\allowbreak{}id}識別，中央ingest後才建立自己的DB
ID。\texttt{producer.sdk\_\allowbreak{}version}為optional，因AI
provider可不用SDK實作同一contract。

\subsection{16.4 Association state}\label{164-association-state}

\begin{longtable}[]{@{}>{\RaggedRight\arraybackslash}p{0.11\linewidth}>{\RaggedRight\arraybackslash}p{0.19\linewidth}>{\RaggedRight\arraybackslash}p{0.31\linewidth}>{\RaggedRight\arraybackslash}p{0.33\linewidth}@{}}
\toprule\noalign{}
state & actors & candidates & 語意 \\
\midrule\noalign{}
\endhead
\bottomrule\noalign{}
\endlastfoot
\texttt{resolved\_\allowbreak{}single} & 1 & 0 & 明確單人 \\
\texttt{resolved\_\allowbreak{}multiple} & \textgreater=2 & 0 & 明確共同參與 \\
\texttt{ambiguous} & 0 & \textgreater=1 & 有候選但無法判定 \\
\texttt{unresolved} & 0 & 0 & 模型無法解析 \\
\texttt{no\_\allowbreak{}player} & 0 & 0 & terminal落地／出界等本來就無球員 \\
\end{longtable}

\texttt{actors{[}{]}}、\texttt{actor\_\allowbreak{}candidates{[}{]}}的association
confidence均optional。Action可放在actor中作optional provider-defined
extension；沒有action時省略。

\subsection{16.5 座標}\label{165-ux5ea7ux6a19}

\begin{itemize}
\tightlist
\item
  \texttt{frame\_\allowbreak{}pos}：影像平面正規化點，x/y在0..1。
\item
  \texttt{frame\_\allowbreak{}bbox}：影像平面正規化框，x1\textless=x2、y1\textless=y2。
\item
  \texttt{frame\_\allowbreak{}foot\_\allowbreak{}pos}：影像中的球員腳底代表點。
\item
  \texttt{court\_\allowbreak{}pos}：AI已投影的球場座標；球場內x/y
  0..1，但場外允許負值或\textgreater1，不可clamp。
\item
  固定球場定義：x=0左端線、x=1右端線（18m）；y=0畫面標準化球場上側邊線、y=1下側邊線（9m）。顯示層可翻轉，資料層不可因教練視角改寫。
\item
  缺值用null／省略與quality flag，不用(0,0)代表unknown。
\end{itemize}

\subsection{16.6 Path Segment}\label{166-path-segment}

Segment只引用相鄰key
point並保存起終代表點陣列；多人事件可有多個position。\texttt{render\_\allowbreak{}state=complete/partial/unavailable}說明能否畫線。Rally
outcome由中央submission join，不在每段重複；避免AI
result與人工得分互相矛盾。

\section{17. AI Callback}\label{17-ai-callback}

\subsection{17.1 Token}\label{171-token}

\begin{itemize}
\tightlist
\item
  Callback URL job-specific。
\item
  Bearer token只對該 job有效，DB只保存 hash。
\item
  可在 TTL內重試，不是用一次即銷毀。
\item
  \texttt{callback\_\allowbreak{}id} 每次 callback唯一；相同 callback
  ID重送回相同結果。
\end{itemize}

\subsection{17.2 Progress／Failed}\label{172-progressfailed}

\begin{Shaded}
\begin{Highlighting}[]
\FunctionTok{\{}
  \DataTypeTok{"callback\_\allowbreak{}id"}\FunctionTok{:} \StringTok{"0190..."}\FunctionTok{,}
  \DataTypeTok{"kind"}\FunctionTok{:} \StringTok{"progress"}\FunctionTok{,}
  \DataTypeTok{"schema\_\allowbreak{}version"}\FunctionTok{:} \StringTok{"1.0.0"}\FunctionTok{,}
  \DataTypeTok{"ai\_\allowbreak{}job\_\allowbreak{}id"}\FunctionTok{:} \StringTok{"..."}\FunctionTok{,}
  \DataTypeTok{"progress"}\FunctionTok{:} \DecValTok{0}\ErrorTok{.}\DecValTok{55}\FunctionTok{,}
  \DataTypeTok{"phase"}\FunctionTok{:} \StringTok{"provider{-}defined optional string"}\FunctionTok{,}
  \DataTypeTok{"message"}\FunctionTok{:} \KeywordTok{null}
\FunctionTok{\}}
\end{Highlighting}
\end{Shaded}

\begin{Shaded}
\begin{Highlighting}[]
\FunctionTok{\{}
  \DataTypeTok{"callback\_\allowbreak{}id"}\FunctionTok{:} \StringTok{"0190..."}\FunctionTok{,}
  \DataTypeTok{"kind"}\FunctionTok{:} \StringTok{"failed"}\FunctionTok{,}
  \DataTypeTok{"schema\_\allowbreak{}version"}\FunctionTok{:} \StringTok{"1.0.0"}\FunctionTok{,}
  \DataTypeTok{"ai\_\allowbreak{}job\_\allowbreak{}id"}\FunctionTok{:} \StringTok{"..."}\FunctionTok{,}
  \DataTypeTok{"error"}\FunctionTok{:} \FunctionTok{\{}
    \DataTypeTok{"code"}\FunctionTok{:} \StringTok{"PROVIDER\_\allowbreak{}PROCESSING\_\allowbreak{}FAILED"}\FunctionTok{,}
    \DataTypeTok{"message"}\FunctionTok{:} \StringTok{"..."}\FunctionTok{,}
    \DataTypeTok{"retryable"}\FunctionTok{:} \KeywordTok{true}
  \FunctionTok{\}}
\FunctionTok{\}}
\end{Highlighting}
\end{Shaded}

中央前端不能依 \texttt{phase}作狀態機，只依 kind/progress。

\subsection{17.3 Completed Multipart}\label{173-completed-multipart}

\begin{Shaded}
\begin{Highlighting}[]
\NormalTok{POST /v1/ai/callback/\{ai\_\allowbreak{}job\_\allowbreak{}id\}}
\NormalTok{Authorization: Bearer \textless{}callback{-}token\textgreater{}}
\NormalTok{Content{-}Type: multipart/form{-}data}
\end{Highlighting}
\end{Shaded}

Parts：

\begin{itemize}
\tightlist
\item
  \texttt{metadata}：callback envelope JSON，含 callback ID與
  checksums。
\item
  \texttt{analysis}：\texttt{application/json}。
\item
  \texttt{overlay}：\texttt{application/vnd.volleyball.overlay+flatbuffers;version=1}。
\end{itemize}

中央以 streaming方式寫 temporary object，不將整個 overlay載入 Node
heap。全部驗證通過後才 commit analysis run。

\subsection{17.4 HTTP Semantics}\label{174-http-semantics}

\begin{longtable}[]{@{}>{\RaggedRight\arraybackslash}p{0.30\linewidth}>{\RaggedRight\arraybackslash}p{0.64\linewidth}@{}}
\toprule\noalign{}
Status & 語意 \\
\midrule\noalign{}
\endhead
\bottomrule\noalign{}
\endlastfoot
200 & 已處理或相同 callback id已處理 \\
202 & 接受，artifact ingestion仍在背景進行 \\
401/403 & token錯誤／過期 \\
409 & callback與 job/submission版本衝突 \\
413 & payload過大 \\
415 & content type/FlatBuffers version不支援 \\
422 & Schema/invariant失敗，不可重試原 payload \\
503 & 暫時不可用，可 retry \\
\end{longtable}

\begin{center}\rule{0.5\linewidth}{0.5pt}\end{center}

\section{18. FlatBuffers Overlay v1}\label{18-flatbuffers-overlay-v1}

\subsection{18.1 原則}\label{181-ux539fux5247}

\begin{itemize}
\tightlist
\item
  AI 回傳一個 rally完整 sequence binary。
\item
  Central驗證後依固定 frame chunk切成 Web用 chunks。
\item
  採 Structure of Arrays，避免大量 nested JS object。
\item
  frame座標量化 U16；court座標 float32，因為可場外。
\item
  缺值由 flags/sentinel表示，不偽造 0。
\end{itemize}

\subsection{18.2 Chunk 內容}\label{182-chunk-ux5167ux5bb9}

\begin{Shaded}
\begin{Highlighting}[]
\NormalTok{start\_\allowbreak{}frame\_\allowbreak{}index}
\NormalTok{frame\_\allowbreak{}count}
\NormalTok{frame\_\allowbreak{}offsets[frame\_\allowbreak{}count + 1]}
\NormalTok{track\_\allowbreak{}ids[detection\_\allowbreak{}count]}
\NormalTok{bbox\_\allowbreak{}x1/y1/x2/y2[detection\_\allowbreak{}count]}
\NormalTok{foot\_\allowbreak{}x/foot\_\allowbreak{}y[detection\_\allowbreak{}count]}
\NormalTok{court\_\allowbreak{}x/court\_\allowbreak{}y[detection\_\allowbreak{}count]}
\NormalTok{action\_\allowbreak{}label\_\allowbreak{}ids[detection\_\allowbreak{}count]}
\NormalTok{confidence\_\allowbreak{}q[detection\_\allowbreak{}count]}
\NormalTok{flags[detection\_\allowbreak{}count]}

\NormalTok{ball\_\allowbreak{}x/ball\_\allowbreak{}y[frame\_\allowbreak{}count]}
\NormalTok{ball\_\allowbreak{}confidence\_\allowbreak{}q[frame\_\allowbreak{}count]}
\NormalTok{ball\_\allowbreak{}flags[frame\_\allowbreak{}count]}
\end{Highlighting}
\end{Shaded}

Invariant：

\begin{itemize}
\tightlist
\item
  \texttt{frame\_\allowbreak{}offsets{[}0{]}\ =\ 0}
\item
  last offset = detection count。
\item
  所有 detection column長度相同。
\item
  action label id \texttt{65535}表示缺少。
\item
  confidence \texttt{-1}表示缺少。
\item
  court validity由 flag表示；x/y本身允許負值或 \textgreater1。
\end{itemize}

\subsection{18.3 Web Manifest}\label{183-web-manifest}

\begin{Shaded}
\begin{Highlighting}[]
\FunctionTok{\{}
  \DataTypeTok{"schema\_\allowbreak{}version"}\FunctionTok{:}\StringTok{"1.0.0"}\FunctionTok{,}
  \DataTypeTok{"analysis\_\allowbreak{}run\_\allowbreak{}id"}\FunctionTok{:}\StringTok{"..."}\FunctionTok{,}
  \DataTypeTok{"overlay\_\allowbreak{}version"}\FunctionTok{:}\DecValTok{2}\FunctionTok{,}
  \DataTypeTok{"video"}\FunctionTok{:\{}\DataTypeTok{"width"}\FunctionTok{:}\DecValTok{1920}\FunctionTok{,}\DataTypeTok{"height"}\FunctionTok{:}\DecValTok{1080}\FunctionTok{,}\DataTypeTok{"total\_\allowbreak{}frames"}\FunctionTok{:}\StringTok{"1079"}\FunctionTok{\},}
  \DataTypeTok{"chunk\_\allowbreak{}frame\_\allowbreak{}count"}\FunctionTok{:}\DecValTok{120}\FunctionTok{,}
  \DataTypeTok{"action\_\allowbreak{}dictionary"}\FunctionTok{:}\OtherTok{[]}\FunctionTok{,}
  \DataTypeTok{"chunks"}\FunctionTok{:}\OtherTok{[}
    \FunctionTok{\{}
      \DataTypeTok{"index"}\FunctionTok{:}\DecValTok{0}\FunctionTok{,}
      \DataTypeTok{"start\_\allowbreak{}frame"}\FunctionTok{:}\StringTok{"0"}\FunctionTok{,}
      \DataTypeTok{"end\_\allowbreak{}frame"}\FunctionTok{:}\StringTok{"119"}\FunctionTok{,}
      \DataTypeTok{"url"}\FunctionTok{:}\StringTok{"/v1/analysis{-}runs/.../overlay{-}chunks/0"}\FunctionTok{,}
      \DataTypeTok{"byte\_\allowbreak{}length"}\FunctionTok{:}\StringTok{"48211"}\FunctionTok{,}
      \DataTypeTok{"sha256"}\FunctionTok{:}\StringTok{"..."}
    \FunctionTok{\}}
  \OtherTok{]}
\FunctionTok{\}}
\end{Highlighting}
\end{Shaded}

Web只預載當前與下一 chunk；seek時取消無用 request並載目標 chunk。

\begin{center}\rule{0.5\linewidth}{0.5pt}\end{center}

\section{19. 教練／裁判端
UX}\label{19-ux6559ux7df4ux88c1ux5224ux7aef-ux}

\subsection{19.1 Routes}\label{191-routes}

\begin{Shaded}
\begin{Highlighting}[]
\NormalTok{/                         場次選單／最近紀錄}
\NormalTok{/matches/:matchId/live    現場總覽}
\NormalTok{/matches/:matchId/paths   球路與熱點}
\NormalTok{/matches/:matchId/players 球員列表與個人頁}
\NormalTok{/matches/:matchId/stats   統計與篩選}
\NormalTok{/matches/:matchId/history Rally／保存視圖歷史}
\NormalTok{/matches/:matchId/replay/:rallyId  單 Rally 回放}
\NormalTok{/annotate/:matchId        標註工作台}
\NormalTok{/settings                 系統設定（依角色）}
\end{Highlighting}
\end{Shaded}

\subsection{19.2 現場頁}\label{192-ux73feux5834ux9801}

\begin{itemize}
\tightlist
\item
  場次、比分、局數、live edge與最新完成分析。
\item
  最新 rally建議卡；必須附資料來源與樣本數。
\item
  快速進入最新回放。
\item
  底部導覽固定但不遮住播放器/球場。
\end{itemize}

\subsection{19.3 球路頁}\label{193-ux7403ux8defux9801}

篩選：

\begin{itemize}
\tightlist
\item
  Team。
\item
  Player（identity mapping存在時）。
\item
  Set／時間範圍。
\item
  Start/target court zone。
\item
  Rally outcome。
\item
  Association quality。
\item
  Action（feature available才顯示）。
\end{itemize}

2D court顯示 A/B、方向、熱點；點擊路徑開對應 rally/time。

\subsection{19.4 回放頁}\label{194-ux56deux653eux9801}

\begin{itemize}
\tightlist
\item
  Video + transparent Canvas overlay。
\item
  2D court同步。
\item
  全部符合篩選路徑淡灰，當前 segment高亮。
\item
  Overlay modes：

  \begin{itemize}
  \tightlist
  \item
    Off
  \item
    Tracking
  \item
    Coach
  \item
    Tactical
  \item
    Debug（只對授權角色）
  \end{itemize}
\item
  Layer toggle：bbox、track
  ID、player、action、ball、trail、footprint、confidence。
\item
  Track ID在未綁 player時顯示 \texttt{Track\ 8}，不能冒充背號。
\end{itemize}

\subsection{19.5
歷史與保存紀錄}\label{195-ux6b77ux53f2ux8207ux4fddux5b58ux7d00ux9304}

\begin{itemize}
\tightlist
\item
  頂部返回場次選單。
\item
  History可按 set、score、日期、processing state、quality篩選。
\item
  Saved Analysis View保存
  filter/query設定，不保存一份過時統計快照；開啟時用當前資料重算，顯示
  saved at/schema version。
\end{itemize}

\subsection{19.6 Offline／Reconnect}\label{196-offlinereconnect}

\begin{itemize}
\tightlist
\item
  WS中斷仍可看已載入影片/資料，但新增標記要顯示 pending。
\item
  Command先放 local outbox，帶 command ID。
\item
  重連後比較 server revision。
\item
  Mapping window過期時，尚未送出的 marker cursor必須重新
  resolve或要求使用者確認。
\end{itemize}

\begin{center}\rule{0.5\linewidth}{0.5pt}\end{center}

\section{20. 分析與統計}\label{20-ux5206ux6790ux8207ux7d71ux8a08}

\subsection{20.1 Baseline
一定可做}\label{201-baseline-ux4e00ux5b9aux53efux505a}

只依人工 outcome、AI contact association與 court positions：

\begin{itemize}
\tightlist
\item
  Rally count。
\item
  Team rally win/loss（排除 unknown）。
\item
  Contact event count。
\item
  Participant event count。
\item
  Source/target heatmap。
\item
  A→B方向分布。
\item
  區域轉移 matrix。
\item
  Terminal位置分布。
\item
  Ambiguous/unresolved比例。
\item
  Samples數、資料品質與缺值率。
\end{itemize}

\subsection{20.2 Identity 後可做}\label{202-identity-ux5f8cux53efux505a}

\begin{itemize}
\tightlist
\item
  每球員 contact數。
\item
  每球員起點／目標熱點。
\item
  每球員 rally participation與 rally outcome split。
\item
  跨 rally球員路徑。
\end{itemize}

\subsection{20.3 Action taxonomy
後才可做}\label{203-action-taxonomy-ux5f8cux624dux53efux505a}

\begin{itemize}
\tightlist
\item
  Attack／serve／receive等 action filter。
\item
  Kill/error/ace/direct outcome。
\item
  Hitting efficiency。
\item
  Setter distribution。
\item
  Side-out／breakpoint細分。
\end{itemize}

在 direct outcome contract未確認前，不能把「包含此事件的
rally最後贏球」當成該事件直接得分。

\subsection{20.4 API 回應}\label{204-api-ux56deux61c9}

每個 metric必須附：

\begin{itemize}
\tightlist
\item
  \texttt{value}
\item
  \texttt{sample\_\allowbreak{}count}
\item
  \texttt{excluded\_\allowbreak{}count}
\item
  \texttt{unknown\_\allowbreak{}count}
\item
  \texttt{quality\_\allowbreak{}breakdown}
\item
  \texttt{feature\_\allowbreak{}dependencies}
\end{itemize}

教練卡片不能只顯示一個無樣本數百分比。

\begin{center}\rule{0.5\linewidth}{0.5pt}\end{center}

\section{21.
後端／前端實作優先順序}\label{21-ux5f8cux7aefux524dux7aefux5be6ux4f5cux512aux5148ux9806ux5e8f}

\subsection{Phase 0 --- Contract Freeze 與
Scaffold}\label{phase-0--contract-freeze-ux8207-scaffold}

\subsubsection{Contracts／SDK Agent}\label{contractssdk-agent}

\begin{itemize}
\tightlist
\item
  建 GraphQL、REST、WS、AI JSON Schema、FlatBuffers。
\item
  建正常、缺球、多 actor、ambiguous、terminal、場外 service fixtures。
\item
  SDK models與validator。
\end{itemize}

\subsubsection{Backend Agent}\label{backend-agent}

\begin{itemize}
\tightlist
\item
  Prisma schema、Compose、health、MinIO init、MediaMTX baseline。
\item
  不先實作完整功能。
\end{itemize}

\subsubsection{Luna Frontend Agent}\label{luna-frontend-agent}

\begin{itemize}
\tightlist
\item
  Nuxt shell、routes、bottom/top nav、design tokens。
\item
  Fixture-driven player/timeline/court prototype。
\end{itemize}

\subsubsection{Exit}\label{exit}

\begin{itemize}
\tightlist
\item
  相同 fixtures通過 TS/Python schema validation。
\item
  Directory ownership不衝突。
\item
  ADR紀錄 API boundary。
\end{itemize}

\subsection{Phase 1 --- Core
Domain／Auth／DB}\label{phase-1--core-domainauthdb}

Backend：

\begin{itemize}
\tightlist
\item
  Match/team/player/set/side assignment。
\item
  Prisma migration。
\item
  GraphQL Yoga endpoint。
\item
  開發 auth與 role guard。
\end{itemize}

Frontend：

\begin{itemize}
\tightlist
\item
  場次選單、history shell、route guards。
\end{itemize}

Exit：可建立場次、set、左右隊與 roster；換邊資料可查。

\subsection{Phase 2 --- Media Ingest／完整
Timeline}\label{phase-2--media-ingestux5b8cux6574-timeline}

Backend／Media：

\begin{itemize}
\tightlist
\item
  MediaMTX ingest/record。
\item
  Segment index → MinIO/DB。
\item
  Timeline GraphQL。
\item
  Live rolling playback window。
\item
  Archive window／lazy seek。
\item
  PlaybackCursor resolve/frame step。
\end{itemize}

Frontend：

\begin{itemize}
\tightlist
\item
  Live video/audio。
\item
  Full timeline、live edge、gap。
\item
  Seek遠端 window與回到 live。
\item
  requestVideoFrameCallback cursor。
\end{itemize}

Exit：一小時以上
capture可持續錄；client記憶體不隨整場線性成長；任意已保存 frame可
resolve。

\subsection{Phase 3 --- Annotation Vertical
Slice}\label{phase-3--annotation-vertical-slice}

Backend：

\begin{itemize}
\tightlist
\item
  WS command/revision/idempotency。
\item
  Draft mutable model + operation log。
\item
  Z/Space/CLOSE\_RALLY(left/right/unknown)/reopen/move/delete/Enter。
\item
  immutable submission。
\end{itemize}

Frontend：

\begin{itemize}
\tightlist
\item
  剪輯式 timeline與 tracks。
\item
  灰／綠 mask、score strip、快捷鍵模式。
\item
  多人 presence/conflict/reconnect。
\end{itemize}

Exit：兩個獨立 PC browser session 同時標註；鍵盤、滑鼠與六個同 command-path touch controls 可用；refresh/reconnect後一致；close不建timestamp／score
frame；unknown可提交。Coach iPad PWA 在 Phase 5 另行驗收，不作為 annotation editor 的主要裝置。

\subsection{Phase 4 --- Clip／Fake
AI／SDK}\label{phase-4--clipfake-aisdk}

\begin{itemize}
\tightlist
\item
  Clip worker與timing manifest。
\item
  AI Integration/capabilities/job/callback。
\item
  SDK v0.1.0可 GitHub安裝。
\item
  Fake provider回 fixtures。
\end{itemize}

Exit：Enter後自動 clip → fake AI → callback → analysis run。

\subsection{Phase 5 --- Coach
Replay／Overlay}\label{phase-5--coach-replayoverlay}

\begin{itemize}
\tightlist
\item
  Overlay sequence ingest/chunk。
\item
  Manifest REST。
\item
  Canvas overlay。
\item
  2D court、A/B path與 history。
\end{itemize}

Exit：點路徑可跳影片；seek只載附近 chunks；action缺失不破版。

\subsection{Phase 6 ---
Identity／Analytics}\label{phase-6--identityanalytics}

\begin{itemize}
\tightlist
\item
  Track-player mapping。
\item
  Baseline metrics/filters/saved views。
\item
  Quality與sample count。
\end{itemize}

Exit：跨 rally player view不依賴 track ID；side switch後 team統計正確。

\subsection{Phase 7 --- Hardening}\label{phase-7--hardening}

\begin{itemize}
\tightlist
\item
  2--3小時 ingest soak。
\item
  Source reconnect/discontinuity。
\item
  Postgres/Redis/MinIO/server/worker restart。
\item
  Callback duplicate/retry/bad checksum。
\item
  iPad memory/performance。
\item
  Backup、retention、metrics、audit。
\end{itemize}

\begin{center}\rule{0.5\linewidth}{0.5pt}\end{center}

\section{22. 三個 Subagent
的協作規則}\label{22-ux4e09ux500b-subagent-ux7684ux5354ux4f5cux898fux5247}

\subsection{22.1 固定 Ownership}\label{221-ux56faux5b9a-ownership}

\begin{longtable}[]{@{}>{\RaggedRight\arraybackslash}p{0.30\linewidth}>{\RaggedRight\arraybackslash}p{0.64\linewidth}@{}}
\toprule\noalign{}
Agent & 可主動修改 \\
\midrule\noalign{}
\endhead
\bottomrule\noalign{}
\endlastfoot
A Contracts/SDK & \texttt{packages/contracts/**}, \texttt{sdk/**} \\
B Backend/Media/Infra & \texttt{server/**}, \texttt{worker/**},
\texttt{packages/db/**}, \texttt{infra/**} \\
C Luna Frontend & \texttt{web/**} \\
Main Agent & root、docs、CI、跨目錄整合與 contract批准 \\
\end{longtable}

任何 Agent若需改別人責任範圍，先回報主 Agent，不直接動手。

\subsection{22.2 第一輪任務}\label{222-ux7b2cux4e00ux8f2aux4efbux52d9}

\subsubsection{Agent A}\label{agent-a}

\begin{quote}
建立 v1 contract baseline與 Python SDK skeleton。只實作
Schema、validation、fake provider adapter與 fixtures，不實作 AI
模型。完成後回報 contract版本、生成物、測試與未定 action/phase事項。
\end{quote}

\subsubsection{Agent B}\label{agent-b}

\begin{quote}
建立 TypeScript
server/worker、Prisma、PostgreSQL、MinIO、Redis、MediaMTX與Compose
skeleton。GraphQL使用 Yoga，binary走REST，高頻
annotation走自訂WS。不得自行發明與 contracts不同的欄位。
\end{quote}

\subsubsection{Agent C - Luna Frontend}\label{agent-c---luna-frontend}

\begin{quote}
建立 PC-first Nuxt 4 annotation workstation，以及 coach/history 多頁 shell、教練端底部導覽、頂部狀態列與 fixture-driven 元件。只有 coach/viewer 顯示面板以 installable iPad landscape PWA 驗收；annotation editor 保留 touch parity 但以 desktop keyboard/mouse 高資訊密度工作流為主。只參考舊 MVP UI libraries/framework，不沿用 domain flow/schema/statistics。
\end{quote}

\subsection{22.3
每次交付格式}\label{223-ux6bcfux6b21ux4ea4ux4ed8ux683cux5f0f}

\begin{Shaded}
\begin{Highlighting}[]
\FunctionTok{\#\# Completed}
\SpecialStringTok{{-} }\NormalTok{...}

\FunctionTok{\#\# Changed files}
\SpecialStringTok{{-} }\NormalTok{...}

\FunctionTok{\#\# Contract/schema version consumed}
\SpecialStringTok{{-} }\NormalTok{...}

\FunctionTok{\#\# Tests run}
\SpecialStringTok{{-} }\NormalTok{command + result}

\FunctionTok{\#\# Known issues / blockers}
\SpecialStringTok{{-} }\NormalTok{...}

\FunctionTok{\#\# Decisions needed from main agent}
\SpecialStringTok{{-} }\NormalTok{...}
\end{Highlighting}
\end{Shaded}

\subsection{22.4 Main Agent Merge Gate}\label{224-main-agent-merge-gate}

\begin{itemize}
\tightlist
\item
  Contract變更先於實作。
\item
  Fixture要同時由 TS/Python讀取。
\item
  Prisma migration與GraphQL field一致。
\item
  UI field必須由 generated type取得。
\item
  每一 Phase都有可執行的 end-to-end evidence。
\item
  不接受「已建立 component/table」作為 flow完成證據。
\end{itemize}

\begin{center}\rule{0.5\linewidth}{0.5pt}\end{center}

\section{23. 主 Agent 可直接使用的
Prompt}\label{23-ux4e3b-agent-ux53efux76f4ux63a5ux4f7fux7528ux7684-prompt}

\begin{Shaded}
\begin{Highlighting}[]
\NormalTok{你現在位於新 Repository \textasciigrave{}volleyball{-}monitoring{-}ai\textasciigrave{} 的 root。}

\NormalTok{你的身份是主要 PM Agent、系統架構師與最終整合者。先完整閱讀：}
\NormalTok{1. \textasciigrave{}docs/MASTER\_\allowbreak{}IMPLEMENTATION\_\allowbreak{}SPEC.md\textasciigrave{}}
\NormalTok{2. \textasciigrave{}AGENTS.md\textasciigrave{}}
\NormalTok{3. \textasciigrave{}packages/contracts/README.md\textasciigrave{}}
\NormalTok{4. 現有 git status、目錄與 placeholder}

\NormalTok{本專案不實作 AI 模型，只實作：}
\NormalTok{{-} Nuxt 標註端與教練端}
\NormalTok{{-} Fastify + GraphQL Yoga + Pothos/Prisma 的 GraphQL/REST/WebSocket 中央後端}
\NormalTok{{-} MediaMTX + FFmpeg 完整 DVR、windowed lazy playback與 clip服務}
\NormalTok{{-} Prisma/PostgreSQL、MinIO、Redis與durable workers}
\NormalTok{{-} 外部 AI Job/callback contracts}
\NormalTok{{-} 可從 GitHub subdirectory安裝的 Python SDK}
\NormalTok{{-} Fake AI Provider與端到端fixtures}

\NormalTok{最多同時派出三個 subagent，固定分工：}
\NormalTok{A. Contracts + Python SDK}
\NormalTok{B. Backend + Media + Infra}
\NormalTok{C. Nuxt Frontend}

\NormalTok{你必須保留架構與Schema決策權。不要讓subagent修改彼此責任目錄；跨目錄變更由你完成。}

\NormalTok{執行規則：}
\NormalTok{{-} 先稽核 scaffold是否與 MASTER spec一致，再依 Phase 0開始。}
\NormalTok{{-} Contract{-}first；跨系統欄位只以 \textasciigrave{}packages/contracts\textasciigrave{} 為真實來源。}
\NormalTok{{-} 不要寫死 action labels、confidence、AI模型參數或群體phase。}
\NormalTok{{-} \textasciigrave{}CLOSE\_RALLY\textasciigrave{}只將指定的server{-}confirmed最後key point標terminal並保存rally{-}level outcome，不建立時間點或得分event。}
\NormalTok{{-} \textasciigrave{}?\textasciigrave{} 是明確unknown，可送出；pending不可送出。}
\NormalTok{{-} Left/right只表示court side，不是team identity。}
\NormalTok{{-} Browser只送PlaybackCursor；authoritative PTS/frame由後端resolve。}
\NormalTok{{-} 整場可回放，但client只lazy{-}load目前數分鐘，不下載全場。}
\NormalTok{{-} GraphQL管structured domain data；REST管media/binary/callback；WebSocket管annotation command/revision；FlatBuffers管逐幀overlay。}
\NormalTok{{-} Submitted revision immutable。}
\NormalTok{{-} Tracker ID只在單rally有效。}
\NormalTok{{-} AI系統只透過Job/Result/Callback/SDK串接。}

\NormalTok{先建立/確認 baseline commit。接著提出 Phase 0 三個subagent的精確任務、路徑ownership與exit criteria，派出最多三個subagent。你自己同時負責root/docs/CI與整合測試。}

\NormalTok{每個Phase結束：}
\NormalTok{{-} 實際執行測試與build}
\NormalTok{{-} 更新 \textasciigrave{}docs/progress.md\textasciigrave{}}
\NormalTok{{-} 列出完成flow、未完成flow、風險與下一Phase}
\NormalTok{{-} 建立小而清楚的commit，不製造一個巨大commit}

\NormalTok{不得把placeholder、mock畫面或單一layer宣稱成已完成。完整flow至少需具備：DB/migration → domain/service → GraphQL/REST/WS → frontend → success/error/reconnect → test。}
\end{Highlighting}
\end{Shaded}

\begin{center}\rule{0.5\linewidth}{0.5pt}\end{center}

\section{24. Luna Worker 定義與建立
Prompt}\label{24-luna-worker-ux5b9aux7fa9ux8207ux5efaux7acb-prompt}

\subsection{24.1 建議 TOML}\label{241-ux5efaux8b70-toml}

\begin{Shaded}
\begin{Highlighting}[]
\DataTypeTok{name} \OperatorTok{=} \StringTok{"luna\_\allowbreak{}worker"}
\DataTypeTok{description} \OperatorTok{=} \StringTok{"Execution{-}focused worker for one bounded repository task with explicit path ownership, acceptance criteria, and no authority to alter architecture or public contracts."}
\DataTypeTok{model} \OperatorTok{=} \StringTok{"gpt{-}5.6{-}luna"}
\DataTypeTok{model\_\allowbreak{}reasoning\_\allowbreak{}effort} \OperatorTok{=} \StringTok{"medium"}
\DataTypeTok{sandbox\_\allowbreak{}mode} \OperatorTok{=} \StringTok{"workspace{-}write"}

\DataTypeTok{developer\_\allowbreak{}instructions} \OperatorTok{=} \StringTok{"""}
\StringTok{Work only on the exact task delegated by the parent agent and only in the explicitly assigned files or directories.}

\StringTok{Before editing:}
\StringTok{1. Read the nearest AGENTS.md files.}
\StringTok{2. Inspect the actual current implementation; do not assume the parent prompt is still accurate.}
\StringTok{3. Restate the task boundary, owned paths, consumed contract/schema version, and acceptance criteria internally before making changes.}

\StringTok{You may perform bounded implementation, file discovery, data organization, fixture work, focused tests, documentation updates, and small fixes that are fully specified by the parent.}

\StringTok{You must not:}
\StringTok{{-} Change product goals, architecture, public contracts, database semantics, or cross{-}team APIs.}
\StringTok{{-} Add, remove, or upgrade dependencies unless explicitly authorized.}
\StringTok{{-} Modify files owned by another active worker.}
\StringTok{{-} Refactor unrelated code, broaden scope, or create speculative abstractions.}
\StringTok{{-} Spawn additional subagents.}
\StringTok{{-} Implement AI models or invent action labels, confidence semantics, or model parameters.}
\StringTok{{-} Hide blockers by guessing.}

\StringTok{If the task requires a cross{-}cutting decision, contradicts the current contract, lacks required input, or would touch unowned paths, stop and report the exact blocker to the parent agent.}

\StringTok{Use the smallest complete change. Run focused validation for changed behavior. Do not claim a test was run unless you actually ran it.}

\StringTok{Return exactly these sections:}
\StringTok{{-} Completed}
\StringTok{{-} Changed files}
\StringTok{{-} Contract/schema version consumed}
\StringTok{{-} Tests run (commands and results)}
\StringTok{{-} Known issues/blockers}
\StringTok{{-} Decision needed from parent (or None)}
\StringTok{"""}
\end{Highlighting}
\end{Shaded}

\texttt{sandbox\_\allowbreak{}mode=workspace-write} 是合理的，因
worker需要實際完成小型修改；若某個 worker只做搜尋/審查，主 Agent可在
spawn時改用 read-only agent。

\subsection{24.2 建立 Prompt}\label{242-ux5efaux7acb-prompt}

\begin{Shaded}
\begin{Highlighting}[]
\NormalTok{請建立個人層級 Codex 自訂 Agent：}

\NormalTok{\textasciitilde{}/.codex/agents/luna{-}worker.toml}

\NormalTok{先使用目前已安裝的 Codex CLI 執行 \textasciigrave{}codex {-}{-}version\textasciigrave{} 與 \textasciigrave{}codex {-}{-}help\textasciigrave{}，並查閱該版本可用的本機/官方 custom agent設定格式；不要假設舊格式仍相同。}

\NormalTok{請建立以下語意的 Agent：}
\NormalTok{{-} name = "luna\_\allowbreak{}worker"}
\NormalTok{{-} model = "gpt{-}5.6{-}luna"}
\NormalTok{{-} model\_\allowbreak{}reasoning\_\allowbreak{}effort = "medium"}
\NormalTok{{-} sandbox\_\allowbreak{}mode = "workspace{-}write"}
\NormalTok{{-} description與developer\_\allowbreak{}instructions採本 repository \textasciigrave{}docs/MASTER\_\allowbreak{}IMPLEMENTATION\_\allowbreak{}SPEC.md\textasciigrave{} 中的 Luna Worker定義；若目前版本欄位名稱不同，只做格式相容調整，不改變行為邊界。}

\NormalTok{luna\_\allowbreak{}worker只負責由主Agent明確委派、邊界清楚、具備path ownership與驗收條件的執行型任務。它不得修改產品目標、架構、public schema、database語意或未授權dependency；不得擴張範圍或再派subagent。遇到跨團隊決策或缺少contract時必須停止並回報。}

\NormalTok{請保留現有其他Codex設定，不覆蓋或刪除無關內容。若目錄或檔案不存在才建立。}

\NormalTok{建立完成後：}
\NormalTok{1. 顯示新增檔案完整內容。}
\NormalTok{2. 顯示只針對該檔案的diff。}
\NormalTok{3. 使用目前系統可用的TOML parser驗證語法。}
\NormalTok{4. 確認 \textasciigrave{}name\textasciigrave{} 是Codex辨識自訂Agent的來源，並確認檔案位於個人Agent目錄。}
\NormalTok{5. 在新的Codex session中以最小的read{-}only測試任務要求主Agent spawn \textasciigrave{}luna\_\allowbreak{}worker\textasciigrave{}；不要用不存在於目前版本help中的虛構CLI子命令。}
\NormalTok{6. 回報版本、驗證命令、驗證結果與任何相容性調整。}
\NormalTok{7. 不修改任何與此Agent無關的設定。}
\end{Highlighting}
\end{Shaded}

\subsection{24.3 Project Config}\label{243-project-config}

\begin{Shaded}
\begin{Highlighting}[]
\CommentTok{\#:schema https://developers.openai.com/codex/config{-}schema.json}

\KeywordTok{[agents]}
\DataTypeTok{enabled} \OperatorTok{=} \ConstantTok{true}
\DataTypeTok{max\_\allowbreak{}concurrent\_\allowbreak{}threads\_\allowbreak{}per\_\allowbreak{}session} \OperatorTok{=} \DecValTok{3}
\end{Highlighting}
\end{Shaded}

\begin{center}\rule{0.5\linewidth}{0.5pt}\end{center}

\section{25. 測試與 Definition of
Done}\label{25-ux6e2cux8a66ux8207-definition-of-done}

\subsection{25.1 Contract}\label{251-contract}

\begin{itemize}
\tightlist
\item
  JSON Schema fixtures通過 Python與TypeScript validator。
\item
  FlatBuffers file identifier/version不符會 fail。
\item
  Passthrough ID改動會 fail。
\item
  Action完全缺失仍通過。
\item
  Confidence缺失仍通過。
\end{itemize}

\subsection{25.2 Media／Time}\label{252-mediatime}

\begin{itemize}
\tightlist
\item
  30fps、59.94fps、VFR input。
\item
  一小時以上 live ingest。
\item
  Source reconnect/discontinuity。
\item
  Timeline gap顯示。
\item
  遠端seek建立新 playback window。
\item
  Client memory不隨整場時間線性成長。
\item
  PlaybackCursor resolve在目標呈現frame可重現。
\item
  Frame step前後一幀正確。
\end{itemize}

\subsection{25.3 Annotation}\label{253-annotation}

\begin{itemize}
\tightlist
\item
  Z/Space/left-close/right-close/unknown-close/Enter六個固定語意。
\item
  \texttt{CLOSE\_\allowbreak{}RALLY}不新增timestamp、score frame或score event，terminal anchor沿用target key point。
\item
  只有 service的service error。
\item
  Unknown可提交，pending不可。
\item
  Reopen與change terminal。
\item
  Duplicate command idempotency。
\item
  多人revision conflict/refetch。
\item
  Submitted immutable/correction draft。
\end{itemize}

\subsection{25.4 Clip／AI}\label{254-clipai}

\begin{itemize}
\tightlist
\item
  pre/post-roll被source開始/結束截斷。
\item
  所有 submission key point有clip mapping。
\item
  AI Job從SDK model解析。
\item
  Clip checksum錯誤。
\item
  Capabilities不相容。
\item
  Job retry/idempotency。
\item
  Callback duplicate、過期token、bad schema、bad FlatBuffers、oversize。
\item
  1:1 key point invariant。
\item
  resolved multiple與ambiguous分開。
\end{itemize}

\subsection{25.5 Coach}\label{255-coach}

\begin{itemize}
\tightlist
\item
  Bottom nav與歷史。
\item
  Video overlay可切層。
\item
  2D path與video seek同步。
\item
  Track未綁player時不顯示假背號。
\item
  Action unavailable時filter/metric隱藏。
\item
  Unknown outcome排除win rate。
\item
  換邊後team統計正確。
\item
  樣本數與品質顯示。
\end{itemize}

\subsection{25.6
第一版完成條件}\label{256-ux7b2cux4e00ux7248ux5b8cux6210ux689dux4ef6}

第一版只有在以下全部成立時才完成：

\begin{itemize}
\tightlist
\item
  伺服器可持續錄整場且前端可任意回看已保存區間。
\item
  Client只windowed lazy-load，live持續錄製。
\item
  Key point authoritative time對齊影片呈現frame。
\item
  多人標註與revision恢復正確。
\item
  灰／綠mask與score state符合使用流程。
\item
  Enter建立immutable submission。
\item
  Clip mapping正確。
\item
  SDK可由GitHub安裝，Fake與真AI使用同一contract。
\item
  Callback/FlatBuffers可驗證、保存與重試。
\item
  教練端多頁、歷史、overlay、A/B與baseline分析可用。
\item
  所有durable state位於PostgreSQL/MinIO；API/worker重啟不遺失。
\end{itemize}

\begin{center}\rule{0.5\linewidth}{0.5pt}\end{center}

\section{26. 舊 MVP
可參考與禁止範圍}\label{26-ux820a-mvp-ux53efux53c3ux8003ux8207ux7981ux6b62ux7bc4ux570d}

只允許參考：

\begin{itemize}
\tightlist
\item
  Nuxt 4／Vue 3／TypeScript。
\item
  Vant/Tailwind/Motion/Lucide等 UI library。
\item
  GraphQL Yoga + Pothos + Prisma + PostgreSQL + MinIO + FFmpeg
  的工程組合。
\item
  Binary upload／stream不走GraphQL的分工。
\item
  WebSocket command ID／server revision概念。
\item
  Docker Compose、Kubernetes、health endpoint、虛擬化列表等工程作法。
\end{itemize}

禁止沿用：

\begin{itemize}
\tightlist
\item
  舊手動球員/事件/球路輸入流程。
\item
  舊 action enum、統計公式或比分/輪轉語意。
\item
  舊 DB schema與 memory store。
\item
  舊頁面資訊架構或 demo data。
\item
  任何與本文件fixed semantics衝突的實作。
\end{itemize}

\begin{center}\rule{0.5\linewidth}{0.5pt}\end{center}

\section{27. 尚待人類／AI 團隊確認，但不阻塞
Baseline}\label{27-ux5c1aux5f85ux4ebaux985eai-ux5718ux968aux78baux8a8dux4f46ux4e0dux963bux585e-baseline}

\begin{enumerate}
\def\labelenumi{\arabic{enumi}.}
\tightlist
\item
  GitHub owner與正式 repository URL。
\item
  正式 auth/SSO。
\item
  第一優先直播輸入協議與 capture硬體。
\item
  Raw segment、playback window、clip、analysis retention。
\item
  Canonical clip FPS/解析度。
\item
  AI action taxonomy、粒度、confidence。
\item
  群體 phase是否有產品用途。
\item
  AI result人工 review UI是否第一版需要。
\item
  正式教練 recommendation規則。
\item
  賽事比分是否需完整 set勝負／暫停／換人功能；baseline至少保存 rally
  outcome與score snapshot。
\end{enumerate}

主 Agent不得為了看起來完整而擅自填補這些未知項目。

\begin{center}\rule{0.5\linewidth}{0.5pt}\end{center}

\section{28.
官方工程參考與選型依據}\label{28-ux5b98ux65b9ux5de5ux7a0bux53c3ux8003ux8207ux9078ux578bux4f9dux64da}

\begin{itemize}
\tightlist
\item
  Codex custom agents / subagents:
  \url{https://developers.openai.com/codex/agent-configuration/subagents}
\item
  Codex config reference:
  \url{https://developers.openai.com/codex/config-reference}
\item
  Nuxt 4: \url{https://nuxt.com/docs/4.x/}
\item
  GraphQL Yoga: \url{https://the-guild.dev/graphql/yoga-server}
\item
  Prisma PostgreSQL:
  \url{https://www.prisma.io/docs/orm/overview/databases/postgresql}
\item
  MediaMTX: \url{https://mediamtx.org/docs/}
\item
  MediaMTX playback: \url{https://mediamtx.org/docs/features/playback}
\item
  HTMLVideoElement requestVideoFrameCallback:
  \url{https://developer.mozilla.org/docs/Web/API/HTMLVideoElement/requestVideoFrameCallback}
\item
  FlatBuffers: \url{https://flatbuffers.dev/}
\item
  pip VCS/subdirectory:
  \url{https://pip.pypa.io/en/stable/topics/vcs-support/}
\end{itemize}

\begin{center}\rule{0.5\linewidth}{0.5pt}\end{center}

\subsection{v3.2 交付註記}\label{v32-ux4ea4ux4ed8ux8a3bux8a18}

本ZIP是可供主Agent啟動Phase 0的「contract-first
scaffold」，不是已完成產品。它包含可驗證的Schema、fixtures、Python
SDK模型、Prisma資料模型、Nuxt PWA shell、Fastify/Yoga/Pothos
shell、worker角色與Docker/Traefik骨架；media ingest、annotation
persistence、clip pipeline與coach
overlay仍須依第21章逐階段完成。不得把placeholder、echo WebSocket或route
shell宣稱為vertical slice完成。

本規格已依實作與使用流程重新稽核；starter
repository只建立contract-first骨架與可驗證fixtures，不代表產品功能已完成。主Agent必須依Phase逐條完成vertical
slice，不得把placeholder當成果。

\section{29. 最終欄位稽核與 Team Contract
Catalog}\label{29-ux6700ux7d42ux6b04ux4f4dux7a3dux6838ux8207-team-contract-catalog}

本章是開工前的最後欄位檢查。若前文範例與本章衝突，以本章、repository中的versioned
schema與golden fixtures為準。

\subsection{29.1 三個團隊的 Schema
交付}\label{291-ux4e09ux500bux5718ux968aux7684-schema-ux4ea4ux4ed8}

\begin{longtable}[]{@{}>{\RaggedRight\arraybackslash}p{0.10\linewidth}>{\RaggedRight\arraybackslash}p{0.16\linewidth}>{\RaggedRight\arraybackslash}p{0.18\linewidth}>{\RaggedRight\arraybackslash}p{0.24\linewidth}>{\RaggedRight\arraybackslash}p{0.24\linewidth}@{}}
\toprule\noalign{}
Producer & Contract & Consumer & Transport & Source of truth \\
\midrule\noalign{}
\endhead
\bottomrule\noalign{}
\endlastfoot
前端／後端 & PlaybackWindow request/descriptor & 後端／前端 & REST JSON
& \texttt{packages/contracts/media/playback-window-request.schema.json}
+ \texttt{playback-window-descriptor.schema.json} \\
前端 & PlaybackCursor & 後端 & WS/REST JSON &
\texttt{packages/contracts/media/playback-cursor.schema.json} \\
前端／後端 & Canonical FrameStep & 後端／前端 & REST JSON &
\texttt{packages/contracts/media/frame-step-request.schema.json} +
\texttt{canonical-frame-anchor.schema.json} \\
後端 & ResolvedMediaAnchor & 前端/DB & WS/REST JSON &
\texttt{packages/contracts/media/resolved-media-anchor.schema.json} \\
前端 & AnnotationCommand & 後端 & WebSocket JSON &
\texttt{packages/contracts/annotation/realtime.schema.json} \\
後端 & Ack/Conflict/Snapshot & 前端 & WebSocket + GraphQL & schema +
Pothos SDL \\
AI & ProviderCapabilities & 後端 & REST JSON &
\texttt{packages/contracts/ai/capabilities.schema.json} \\
後端前處理 & AIJobRequest & AI SDK/provider & REST JSON &
\texttt{packages/contracts/ai/job.schema.json} \\
AI & JobAccepted & 後端 & REST JSON &
\texttt{packages/contracts/ai/job-accepted.schema.json} \\
AI & AnalysisResult & 後端 & callback JSON part &
\texttt{packages/contracts/ai/result.schema.json} \\
AI & OverlaySequence & 後端 & callback FlatBuffers part &
\texttt{flatbuffers/overlay.fbs} \\
後端 & OverlayManifest/Chunk & 前端 & REST JSON + FlatBuffers & manifest
schema + \texttt{overlay-chunk.fbs} \\
後端 & Domain query/mutation & 前端 & GraphQL & Pothos code-first
generated SDL \\
\end{longtable}

\subsection{29.2 前端送後端的
PlaybackCursor}\label{292-ux524dux7aefux9001ux5f8cux7aefux7684-playbackcursor}

\begin{longtable}[]{@{}>{\RaggedRight\arraybackslash}p{0.10\linewidth}>{\RaggedRight\arraybackslash}p{0.16\linewidth}>{\RaggedRight\arraybackslash}p{0.18\linewidth}>{\RaggedRight\arraybackslash}p{0.24\linewidth}>{\RaggedRight\arraybackslash}p{0.24\linewidth}@{}}
\toprule\noalign{}
欄位 & 型別 & 必填 & Ownership & 說明 \\
\midrule\noalign{}
\endhead
\bottomrule\noalign{}
\endlastfoot
\texttt{playback\_\allowbreak{}window\_\allowbreak{}id} & UUID/string & 是 & CLIENT\_\allowbreak{}OBSERVED +
PASSTHROUGH & 當前播放器window \\
\texttt{mapping\_\allowbreak{}version} & integer & 是 & BACKEND產生、client
passthrough & 防止過期mapping \\
\texttt{player\_\allowbreak{}media\_\allowbreak{}time\_\allowbreak{}us} & decimal string & 是 &
CLIENT\_\allowbreak{}OBSERVED & 該window的player timeline \\
\texttt{observation\_\allowbreak{}source} & enum & 是 & CLIENT\_\allowbreak{}OBSERVED &
rvfc或fallback \\
\texttt{presented\_\allowbreak{}frames} & decimal string/null & 否 & CLIENT\_\allowbreak{}OBSERVED
& 只供audit/掉幀診斷 \\
\texttt{seek\_\allowbreak{}generation} & integer & 是 & CLIENT\_\allowbreak{}OBSERVED &
避免舊frame callback \\
\texttt{cursor\_\allowbreak{}status} & enum & 是 & CLIENT\_\allowbreak{}OBSERVED &
ready才允許標記 \\
\end{longtable}

不應從browser傳
\texttt{source\_\allowbreak{}pts}、\texttt{capture\_\allowbreak{}time\_\allowbreak{}us}或\texttt{capture\_\allowbreak{}frame\_\allowbreak{}index}作canonical值。

\subsection{29.3 後端解析後的 KeyPoint
Anchor}\label{293-ux5f8cux7aefux89e3ux6790ux5f8cux7684-keypoint-anchor}

\begin{longtable}[]{@{}>{\RaggedRight\arraybackslash}p{0.11\linewidth}>{\RaggedRight\arraybackslash}p{0.19\linewidth}>{\RaggedRight\arraybackslash}p{0.31\linewidth}>{\RaggedRight\arraybackslash}p{0.33\linewidth}@{}}
\toprule\noalign{}
欄位 & 型別 & Ownership & 說明 \\
\midrule\noalign{}
\endhead
\bottomrule\noalign{}
\endlastfoot
\texttt{capture\_\allowbreak{}epoch\_\allowbreak{}id} & UUID & BACKEND\_\allowbreak{}DERIVED & raw
PTS有效範圍 \\
\texttt{source\_\allowbreak{}pts} & decimal string & BACKEND\_\allowbreak{}DERIVED &
epoch內來源PTS \\
\texttt{source\_\allowbreak{}time\_\allowbreak{}base} & rational & BACKEND\_\allowbreak{}DERIVED & PTS單位 \\
\texttt{capture\_\allowbreak{}time\_\allowbreak{}us} & decimal string & BACKEND\_\allowbreak{}DERIVED &
整場canonical時間 \\
\texttt{capture\_\allowbreak{}frame\_\allowbreak{}index} & decimal string & BACKEND\_\allowbreak{}DERIVED &
整場canonical frame \\
\texttt{timing\_\allowbreak{}precision} & enum & BACKEND\_\allowbreak{}DERIVED &
frame\_\allowbreak{}exact/pts\_\allowbreak{}exact/estimated \\
\texttt{snap\_\allowbreak{}distance\_\allowbreak{}us} & decimal string/null & BACKEND\_\allowbreak{}DERIVED &
client observation到sample距離 \\
\texttt{original\_\allowbreak{}playback\_\allowbreak{}cursor} & JSON audit & CLIENT\_\allowbreak{}OBSERVED &
不作canonical查詢欄位 \\
\end{longtable}

\subsection{29.4 Immutable Submission
Snapshot}\label{294-immutable-submission-snapshot}

必要：submission ID、rally ID、annotation revision、content hash、score
resolution、nullable scoring court side/team、left/right team
snapshot、nullable scores、service/terminal key point IDs、clip policy
version與實際pre/post-roll快照、submitted by/at、supersedes
ID與完整submission key points。

\texttt{unknown} submission：score resolution=unknown；scoring
side/team、score
before/after與PointAward皆為null/不存在。\texttt{resolved}
submission：side/team與score ledger必須完整且transaction一致。

\subsection{29.5 AI
Job／Result最小原則}\label{295-ai-jobresultux6700ux5c0fux539fux5247}

\begin{itemize}
\tightlist
\item
  Job只傳已裁切影片、局部key points、人工得分解析與callback。
\item
  Job不傳court definition、requirements、模型參數、MinIO prefix或team
  metadata。
\item
  Result不重複人工得分、side assignment或team truth；中央以submission
  join。
\item
  所有跨AI邊界ID在欄位表標成PASSTHROUGH。
\item
  AI生成欄位只包含tracks、contact events、path
  segments、overlay與optional extensions。
\end{itemize}

\subsection{29.6
已知待確認但不阻塞介面}\label{296-ux5df2ux77e5ux5f85ux78baux8a8dux4f46ux4e0dux963bux585eux4ecbux9762}

\begin{enumerate}
\def\labelenumi{\arabic{enumi}.}
\tightlist
\item
  AI action真實taxonomy、粒度與confidence。
\item
  群體動作／比賽phase是否有可驗收產品用途。
\item
  第一個正式來源protocol與canonical recording profile。
\item
  pre/post-roll預設秒數。
\item
  production auth provider與retention policy。
\end{enumerate}

這些不得被subagent猜成既定事實。先以optional extension、設定值或ADR open
decision保留。

\section{30. v3.2
最終交付與啟動}\label{30-v32-ux6700ux7d42ux4ea4ux4ed8ux8207ux555fux52d5}

Repository ZIP內必須同時包含：

\begin{itemize}
\tightlist
\item
  \texttt{docs/SYSTEM\_\allowbreak{}SPEC\_\allowbreak{}V3\_\allowbreak{}2.md}
\item
  \texttt{docs/SYSTEM\_\allowbreak{}SPEC\_\allowbreak{}V3\_\allowbreak{}2.tex}
\item
  \texttt{docs/SYSTEM\_\allowbreak{}SPEC\_\allowbreak{}V3\_\allowbreak{}2.pdf}
\item
  \texttt{docs/MAIN\_\allowbreak{}AGENT\_\allowbreak{}PROMPT.md}
\item
  \texttt{CODEX\_\allowbreak{}SOL\_\allowbreak{}PROMPT.txt}
\item
  \texttt{.codex/agents/contracts-sdk-worker.toml}
\item
  \texttt{.codex/agents/backend-worker.toml}
\item
  \texttt{.codex/agents/luna-worker.toml}
\item
  完整starter目錄、contract fixtures、Python SDK、Compose與Traefik設定。
\end{itemize}

本地啟動：

\begin{Shaded}
\begin{Highlighting}[]
\FunctionTok{cp}\NormalTok{ .env.example .env}
\ExtensionTok{docker}\NormalTok{ compose }\AttributeTok{{-}f}\NormalTok{ infra/compose.yaml }\AttributeTok{{-}{-}profile}\NormalTok{ app }\AttributeTok{{-}{-}profile}\NormalTok{ dev{-}ai up }\AttributeTok{{-}{-}build}
\end{Highlighting}
\end{Shaded}

iPad以Safari開啟Docker host的LAN
IP或DNS，確認PWA、WebSocket、GraphQL、HLS均為同源。Traefik
\texttt{:8080} dashboard與MinIO
\texttt{:9001}只供本地開發，不是正式公開入口。

\end{document}
